\RequirePackage{xstring}
\IfSubStr{\jobname}{\detokenize{with_notes}}{%
  % The case with comments
  \xdef\myDevelopVariable{1}
  \xdef\myHideVariable{0}
  \documentclass[fleqn
    ,fontsize=12pt
    ,DIV=13
    ,BCOR=8mm
    ,openany
  ]{./myClass}
}{%
  \xdef\myDevelopVariable{0}
  \xdef\myHideVariable{0}
  \documentclass[fleqn
    ,fontsize=12pt
    ,DIV=13
    ,BCOR=8mm
    ,openany
  ]{./myClass}
  % \ifnum\myHideVariable=1
  % \usepackage{showframe}
  % \fi
}
\usepackage{layout}

\bibliography{literatur.bib}

%TODO: replace def --> newcommand
\def\bme{\boldsymbol{e}}

% {{{ bb
\def\A{\ensuremath \mathbb{A}}
%\def\B{\ensuremath \mathbb{B}}
\def\C{\ensuremath \mathbb{C}}
%\def\D{\ensuremath \mathbb{D}}
%\def\E{\ensuremath \mathbb{E}}
%\def\F{\ensuremath \mathbb{F}}
%\def\G{\ensuremath \mathbb{G}}
%\def\H{\ensuremath \mathbb{H}}
\def\I{\ensuremath \mathbb{I}}
%\def\J{\ensuremath \mathbb{J}}
%\def\K{\ensuremath \mathbb{K}}
%\def\L{\ensuremath \mathbb{L}}
%\def\M{\ensuremath \mathbb{M}}
\def\N{\ensuremath \mathbb{N}}
%\def\O{\ensuremath \mathbb{O}}
\def\P{\ensuremath \mathbb{P}}
\def\Q{\ensuremath \mathbb{Q}}
\def\R{\ensuremath \mathbb{R}}
%\def\S{\ensuremath \mathbb{S}}
%\def\T{\ensuremath \mathbb{T}}
%\def\U{\ensuremath \mathbb{U}}
%\def\V{\ensuremath \mathbb{V}}
%\def\W{\ensuremath \mathbb{W}}
%\def\X{\ensuremath \mathbb{X}}
%\def\Y{\ensuremath \mathbb{Y}}
\def\Z{\ensuremath \mathbb{Z}}
% }}}

% {{{ cal
%\def\cA{\ensuremath \mathcal{A}}
%\def\cB{\ensuremath \mathcal{B}}
\def\cC{\ensuremath \mathcal{C}}
\def\cD{\ensuremath \mathcal{D}}
\def\cE{\ensuremath \mathcal{E}}
\def\cF{\ensuremath \mathcal{F}}
\def\cG{\ensuremath \mathcal{G}}
%\def\cH{\ensuremath \mathcal{H}}
%\def\cI{\ensuremath \mathcal{I}}
%\def\cJ{\ensuremath \mathcal{J}}
%\def\cK{\ensuremath \mathcal{K}}
\def\cL{\ensuremath \mathcal{L}}
\def\cM{\ensuremath \mathcal{M}}
\def\cN{\ensuremath \mathcal{N}}
\def\cO{\ensuremath \mathcal{O}}
\def\cP{\ensuremath \mathcal{P}}
%\def\cQ{\ensuremath \mathcal{Q}}
\def\cR{\ensuremath \mathcal{R}}
%\def\cS{\ensuremath \mathcal{S}}
%\def\cT{\ensuremath \mathcal{T}}
%\def\cU{\ensuremath \mathcal{U}}
\def\cV{\ensuremath \mathcal{V}}
%\def\cW{\ensuremath \mathcal{W}}
%\def\cX{\ensuremath \mathcal{X}}
%\def\cY{\ensuremath \mathcal{Y}}
%\def\cZ{\ensuremath \mathcal{Z}}
% }}}

% {{{ scr
%\def\sA{\ensuremath \mathscr{A}}
%\def\sB{\ensuremath \mathscr{B}}
%\def\sC{\ensuremath \mathscr{C}}
%\def\sD{\ensuremath \mathscr{D}}
\def\sE{\ensuremath \mathscr{E}}
%\def\sF{\ensuremath \mathscr{F}}
%\def\sG{\ensuremath \mathscr{G}}
%\def\sH{\ensuremath \mathscr{H}}
%\def\sI{\ensuremath \mathscr{I}}
%\def\sJ{\ensuremath \mathscr{J}}
%\def\sK{\ensuremath \mathscr{K}}
%\def\sL{\ensuremath \mathscr{L}}
%\def\sM{\ensuremath \mathscr{M}}
%\def\sN{\ensuremath \mathscr{N}}
\def\sO{\ensuremath \mathscr{O}}
%\def\sP{\ensuremath \mathscr{P}}
%\def\sQ{\ensuremath \mathscr{Q}}
%\def\sR{\ensuremath \mathscr{R}}
%\def\sS{\ensuremath \mathscr{S}}
%\def\sT{\ensuremath \mathscr{T}}
%\def\sU{\ensuremath \mathscr{U}}
%\def\sV{\ensuremath \mathscr{V}}
%\def\sW{\ensuremath \mathscr{W}}
%\def\sX{\ensuremath \mathscr{X}}
%\def\sY{\ensuremath \mathscr{Y}}
%\def\sZ{\ensuremath \mathscr{Z}}
% }}}

%\renewcommand{\headrulewidth}{0.2pt}

\newcommand{\myhr}{\rule{0.3\textwidth}{1pt}}

\let\epsilon\varepsilon
\let\phi\varphi
\let\hat\widehat

% {{{
%\DeclareMathOperator{\defn}{def}
\DeclareMathOperator{\modulo}{mod}
\DeclareMathOperator{\conv}{conv}
\DeclareMathOperator{\Cof}{Cof}
\DeclareMathOperator{\grad}{grad}
\DeclareMathOperator{\Id}{Id}
\DeclareMathOperator{\id}{id}
\DeclareMathOperator{\dist}{dist}
\DeclareMathOperator{\diam}{diam}
\DeclareMathOperator{\co}{co}
\DeclareMathOperator{\supp}{supp}
\DeclareMathOperator{\graph}{graph}
\DeclareMathOperator{\slopes}{slopes}
\DeclareMathOperator{\Ob}{Ob}
\DeclareMathOperator{\Mor}{Mor}
\DeclareMathOperator{\im}{im}
\DeclareMathOperator{\Car}{Car}
\DeclareMathOperator{\coker}{coker}
\DeclareMathOperator{\ord}{ord}
% }}}

%\def\cP{\text{slopes}}

% {{{
\newcommand{\Cfx}{{\mathbb C\llbracket x\rrbracket}}
\newcommand{\Cft}{{\mathbb C\llbracket t\rrbracket}}
\newcommand{\Cfu}{{\mathbb C\llbracket u\rrbracket}}
\newcommand{\Cfxl}{{\mathbb C (\!(x)\!)}}
\newcommand{\Cftl}{{\mathbb C (\!(t)\!)}}
\newcommand{\Cful}{{\mathbb C (\!(u)\!)}}
\newcommand{\Cfzl}{{\mathbb C (\!(z)\!)}}
\newcommand{\Ckx}{{\mathbb C \{x\}}}
\newcommand{\Ckt}{{\mathbb C \{t\}}}
\newcommand{\Cku}{{\mathbb C \{u\}}}
\newcommand{\Ckxl}{{\mathbb C (\!\{x\}\!)}}
\newcommand{\Cktl}{{\mathbb C (\!\{t\}\!)}}
\newcommand{\Ckul}{{\mathbb C (\!\{u\}\!)}}

\newcommand{\bydef}{\overset{\mbox{\tiny def}}{=}}

\newcommand{\anfzeichen}[1]{\glqq{#1}\grqq}     % Anf\"uhrungszeichen

\providecommand*{\Floor}[1]{\left\lfloor #1 \right\rfloor}         % gro\ss{}es ganzes Abrunden
\providecommand*{\Ceil}[1]{\left\lceil #1 \right\rceil} % gro\ss{}es ganzes

\makeatother

\newcommand{\overbox}[2]{\ensuremath\begin{array}[b]{c}%
\makebox[0pt]{\fbox{\scriptsize#2}}\\[-2pt]\text{\small$\downarrow$}\\[-3pt]%
{\displaystyle#1}\end{array}}

% vim:set ft=tex foldmethod=marker foldmarker={{{,}}}:


% %%%%%%%%%%%%%%%%%%%%%%%%%%%%%%%%%%%%%%%%%%%%%%%%%%%%%%%%%%%%%%%%%%%%%
% \DeclareMathOperator{\Gl}{GL}
% \DeclareMathOperator{\GL}{GL}
% \DeclareMathOperator{\Ad}{Ad}
% \DeclareMathOperator{\ad}{ad}
% \DeclareMathOperator{\End}{End}
% \DeclareMathOperator{\Aut}{Aut}
% \DeclareMathOperator{\Proj}{Proj}
% \DeclareMathOperator{\Spec}{Spec}
% \DeclareMathOperator{\codim}{codim}
% \DeclareMathOperator{\Lie}{Lie}
% \DeclareMathOperator{\Res}{Res}
% \DeclareMathOperator{\Tr}{Tr}
% \DeclareMathOperator{\diag}{diag}

\let\tilde\widetilde
\newcommand{\bdelta}{\boldsymbol{\delta}}

%%%%%%%%%%%%%%%%%%%%%%%%%%%%%%%%%%%%%%%%%%%%%%%%%%%%%%%%%%%%%%%%%%%%%
%%%%  testing                                                    %%%%
%%%%%%%%%%%%%%%%%%%%%%%%%%%%%%%%%%%%%%%%%%%%%%%%%%%%%%%%%%%%%%%%%%%%%

\usepackage{paracol}
\usepackage{multicol}
\usepackage{longtable}
\usepackage{relsize}

\usepackage{graphicx}

\usepackage[pages=some]{background}

% \newcommand{\myrel}[1]{\underset{#1,\max}{\prec}}
\newcommand{\myrelPrime}[2]{%
    \scalerel*{%
      \mathlarger{\mathlarger{\prec}}
      \kern-#2
      \raisebox{#1}{$\mathsmaller{\prec}$}
      \,
    }{\theta}
  }
\newcommand{\myrelPrimeB}[2]{%
    \scalerel*{%
      \mathlarger{\mathlarger{\prec}}
      \kern-#2
      \raisebox{#1}{$\prec$}
      \,
    }{\theta}
  }
\newcommand{\myrel}[1]{\underset{#1}{%
    \mathchoice{\myrelPrimeB{1.3125pt}{.70em}}
      {\myrelPrime{1.56pt}{.54em}}
      {\myrelPrime{.975pt}{.45em}}
      {\myrelPrime{1pt}{.45em}}
}}

\begin{document}

\ifnum\myDevelopVariable=0
  \frontmatter
  % \newgeometry{left=3cm,right=3cm}

\KOMAoptions{twoside = false}
\backgroundsetup{scale=1,placement=bottom,contents={%
    \begin{tikzpicture}[remember picture,overlay
      ,scale=6
      ,transform canvas={xshift = -6cm,yshift = 5.5cm}]
    \pgfmathsetmacro{\abs}{0.4};
    \foreach \x in {-2,-1,...,3}{
        \foreach \argA/\absA in {-0.5/0.4
                                ,0.3/0.35
                                ,0.1/0.27
                                ,0.65/0.23}{
        \begin{scope}
            % \clip (-0.65,{\abs+0.1}) rectangle (2.63,{-\abs-0.1});
            \pgfmathsetmacro{\s}{{\argA + \x* 2 - 1/2}};
            \draw[blue!40!white] (\s,0)
                sin ({\s + 1/2},{-\absA})
                cos ({\s + 2/2},0)
                sin ({\s + 3/2},\absA)
                cos ({\s + 4/2},0);
            \draw[dotted] ({\s + 1/2},{-\absA}) -- ({\s + 1/2},-0.5);
            \fill[white] ({\s + 1/2},{-\absA}) circle (1pt);
            \fill[red] ({\s + 1/2},{-\absA}) circle (.4pt);
            \draw[blue!40!black] (\s,0)
                sin ({\s + 1/2},\absA)
                cos ({\s + 2/2},0)
                sin ({\s + 3/2},{-\absA})
                cos ({\s + 4/2},0);
            \draw[dotted] ({\s + 3/2},{-\absA}) -- ({\s + 3/2},-0.5);
            \fill[white] ({\s + 3/2},{-\absA}) circle (1pt);
            \fill[red] ({\s + 3/2},{-\absA}) circle (.4pt);
        \end{scope}
        }
    }
    \end{tikzpicture}%
}}
\begin{titlepage}
  \BgThispage
  \thispagestyle{empty}
  % \begin{flushright}
    \vspace{-1cm}
    \includegraphics[width=0.6\textwidth]{logo/Uni_Aug_Logo_IFM_RGB-crop.png}
    \vspace{1cm}
  % \end{flushright}
  \vspace{2cm}
  \newcommand{\Rule}{%
    \textcolor{myheadLineColor}{\rule{\textwidth}{0.5mm}}%
  }
  \begin{center}\sffamily
    \normalfont\sffamily\large
    Master's thesis
    %\vfill
    \Rule
    \vspace{5mm}
    \textcolor{myheadLineColor}{
      \Huge{Classification of meromorphic connections using Stokes structures
        and Stokes groups}
    }
    \vspace{1mm}
    \Rule
  \end{center}
    %\vfill
  \begin{center}
    \makebox[.2\textwidth][r]{\normalfont\sffamily\large by}
    \makebox[.5\textwidth][l]{\normalfont\sffamily\Large Maximilian Huber}
    \\ \vspace{1cm}
    \makebox[.2\textwidth][r]{\normalfont\sffamily\large submitted to the}
    \makebox[.5\textwidth][l]{\normalfont\sffamily\Large Institut für
    Mathematik}
    \\
    \makebox[.2\textwidth][r]{\normalfont\sffamily\large of}
    \makebox[.5\textwidth][l]{\normalfont\sffamily\Large Universität Augsburg}
    \\ \vspace{1cm}
    \makebox[.2\textwidth][r]{\normalfont\sffamily\large approved by}
    \makebox[.5\textwidth][l]{\normalfont\sffamily\Large Prof.\ Dr.\ Marco Hien}
    \\
    \makebox[.2\textwidth][r]{\normalfont\sffamily\large and}
    \makebox[.5\textwidth][l]{\normalfont\sffamily\Large Prof.\ Dr.\ Marc
    Nieper-Wißkirchen}
    \\ \vspace{1cm}
    \makebox[.2\textwidth][r]{\normalfont\sffamily\large submitted on}
    \makebox[.5\textwidth][l]{\normalfont\sffamily\Large \today}
  \end{center}
  \vfill
  \vfill
\end{titlepage}
\KOMAoptions{twoside}

% \restoregeometry

%%% Local Variables:
%%% TeX-master: "Maximilian_Huber-Masters_Thesis-with_notes.tex"
%%% End:

  \newpage \thispagestyle{empty} \mbox{}
  \setcounter{page}{0}
\else
  \mainmatter
\fi

\tableofcontents{}
\begin{comment}
  \today
\end{comment}

% \ifnum\myHideVariable=0
  \chapter{Introduction}
\TODO[see~\cite{van2003galois} for better text]
Let $t$ be a local coordinate around $0$ in $\C$.
The basic problem, from which this master's thesis arises, is that there are
differential equations $(\frac{d}{dt}-A)\hat v=w$ with coefficients in
convergent powers series, which have no solution with converging entries.
By looking at asymptotic solutions on sectors and their relations on overlaps
one obtains the Stokes phenomenon.
One can use the language of meromorphic connections or the language of systems,
to look at this problem and the interesting object is then the classifying
set\comm{, which is defined as\dots} (cf.\ Section~\ref{sec:classifyingSet}).
The information contained in the Stokes phenomenon will be exactly the
information needed to classify meromorphic connections up to convergent
transformation.

We will only be interested in local classification and thus only in local
information of meromorphic connections. This classification can be splitted in
the coarse formal classification and the fine meromorphic classification.

The formal classification problem was solved by the Levelt-Turittin Theorem
(cf.\ Theorem~\ref{thm:leveltTurittin}). It states that in every formal
equivalence class are some meromorphic connections of special form, which will
be called models.
These models are meromorphic connections, which are defined to be isomorphic to
some direct sum of elementary meromorphic connections and elementary
meromorphic connections are well understood.

Starting with a formal equivalence class corresponding to some model $A^0$, we
can apply the meromorphic classification.
The obtained set of meromorphic classes will be called the classifying set, but
we will look at the slightly larger space $\cH(A^0)$ of meromorphic pairs
(cf.\ Section~\ref{sec:classifyingSet}).
The tool to describe the set of meromorphic pairs will be the Stokes structures.
% which turn out to provide exactly the needed information.
This idea is formulated in the Malgrange-Sibuya Theorem
(cf.\ Theorem~\ref{thm:mainThm1}), where the Stokes structures appear as the
first cohomology $H^1(S^1;\Lambda(A^0))$ of the Stokes sheaf
$\Lambda(A^0)$ on $S^1$ (cf.\ Definition~\ref{defn:StokesSheaf}).

The Malgrange-Sibuya theorem can be improved by showing that in each element
in the $H^1(S^1;\Lambda(A^0))$ contains a unique cocycle called the Stokes
cocycle (cf.\ Definition~\ref{defn:stokesCocycle}) of special form.
These Stokes cocycles are given by the elements in the product over some
special directions $\theta\in\A\subset S^1$ determined by $A^0$
(cf.\ Definition~\ref{defn:antiStokesDir}) of Stokes groups
$\Sto_\theta(A^0)\subset\Lambda_\theta(A^0)$
(cf.\ Definition~\ref{defn:stokesGroup}).
We will see that $\Sto_\theta(A^0)$ has the faithful representation
$\SSto_\theta(A^0)$, which has a rather simple definition
(cf.\ Proposition~\ref{prop:representation}).
The elements of $\SSto_\theta(A^0)$ are the so-called Stokes matrices and it is
easy to see that they are nilpotent.
Since the Stokes matrices also provide the required information to describe a
meromorphic class of a meromorphic connection one obtains an isomorphism
$\C^N\to H^1(S^1;\Lambda(A^0))$, where $N$ is the irregularity of $A^0$
(cf.\ Section~\ref{sec:WhichInformationIsNeeded}).
This characteristic can be used to define the structure of a nilpotent Lie
group on the classifying set $\cH(A^0)$.

%%% Local Variables:
%%% TeX-master: "Maximilian_Huber-Masters_Thesis-with_notes.tex"
%%% End:

% \fi
\ifnum\myHideVariable=0
  \addchap*{List of used symbols}
\textbf{1.~Poincaré asymptotic expansions}
\begin{longtable}[h]{>{\raggedright}p{4cm}@{\hspace{.2cm}}p{10cm}}
  $\mathfrak{s}_{a,b}(r)$ \dotfill~& the sector
    $\left\{t\in\C \mid a<\arg(t)<b, 0<|t|<r\right\}$\\
  $\mathfrak{s}_{I}(r)$ \dotfill~& the sector $\mathfrak{s}_{a,b}(r)$ for $I=(a,b)$ or
    the arc $I$\\
  $\bar{\mathfrak{s}}_{a,b}(r)$ \dotfill~& the closure of $\mathfrak{s}_{a,b}(r)$ in
    $\C^*$.\\
  $(\theta,\theta')$ \dotfill~& the arc from $\theta\in S^1$ to $\theta'$\\
  $U(\theta,\epsilon)$ \dotfill~& the arc $\left(\theta-\frac{\epsilon}{2}
    ,\theta+\frac{\epsilon}{2}\right)\subset S^1$ resp.\ the corresponding
    sector\\
  $\mathfrak{s}\Subset\mathfrak{s}'$ \dotfill~&
    $\mathfrak{s}$ is proper sub-sector of $\mathfrak{s}'$\\
  $\hat f$ \dotfill~& the asymptotic expanson of $f$ (if not used in another way)\\
  $f\sim_{\mathfrak{s}}\hat f$ \dotfill~& $f$ has $\hat f$ as asymptotic expanson on the
    sector $\mathfrak{s}$\\
  $f\sim_{I}\hat f$ \dotfill~& $f$ has $\hat f$ as asymptotic expanson on the arc $I$\\
  $f\sim_{\theta}\hat f$ \dotfill~& $f$ has $\hat f$ as asymptotic expanson in the
    direction $\theta$\\
  $\cA$ \dotfill~& the sheaf of asymptotic expansions\\
  $T_I$ \dotfill~& the Taylor map on the arc $I$\\
  $T_{\mathfrak{s}}$ \dotfill~& the Taylor map on the sector $\mathfrak{s}$\\
  $\cA^{<0}(I)$ \dotfill~& $\ker(T_I)$; the functions asymptotic to $0$ on $I$\\
  $\cA^{<0}(\mathfrak{s})$ \dotfill~& $\ker(T_{\mathfrak{s}})$; the functions asymptotic
    to $0$ on $\mathfrak{s}$\\
\end{longtable}

\textbf{2.~Systems and meromorphic connections}
\begin{longtable}[h]{>{\raggedright}p{4cm}@{\hspace{.2cm}}p{10cm}}
  $M$ \dotfill~& a riemanian surface\\
  $Z$ \dotfill~& a effective divisor on $M$\\
  $\sM$ \dotfill~& a holomorphic bundle over $M$\\
  $\sM$; $(\sM,\nabla)$ \dotfill~& a meromorphic connection on $\sM$ with poles
    on $Z$\\
  $\cM$; $ (\cM,\nabla)$ \dotfill~& germ of a meromorphic connection
    $(\sM,\nabla)$ at $0\in M$\\
  $\triangle$ \dotfill~& a differential operator\\
  $[A]$ \dotfill~& the system corresponding to the connection matrix $A$\\
  $G[t]$ \dotfill~& $\Gl_n(\C[t])$\\
  $G\{t\}$ \dotfill~& $\Gl_n(\C\{t\})$; the analytic transformations\\
  $G(\!\{t\}\!)$ \dotfill~& $\Gl_n\left(\C\{t\}[t^{-1}]\right)$; the meromorphic
    transformations\\
  $G\llbracket t\rrbracket$ \dotfill~& $\Gl_n\left(\C\llbracket t\rrbracket\right)$;
    the maybe not applicable formal transformations\\
  $G(\!(t)\!)$ \dotfill~& $\Gl_n\left(\C\llbracket t\rrbracket[t^{-1}]\right)$;
    the maybe not applicable formal meromorphic transformations\\
  ${}^F\!A$ \dotfill~& $(dF)F^{-1}+FAF^{-1}$; the transforamtion of $A$ by $F$\\
  $\hat G(A)$ \dotfill~& the set of all (applicable) formal transformations\\
  $[A,B]$ \dotfill~& the linear differential system $\frac{dF}{dt}=BF-FA$\\
  $[\End A]$ \dotfill~& $[A^0,A^0]$\\
  $G_0(A)$ \dotfill~& the set of all isotropies of $A$\\
  $\cY$ \dotfill~& a fundamental solution\\
  $\cE^{\phi}$ \dotfill~& the germ $(\C(\!\{t\}\!),d-\phi')$\\
  $\cN_{\alpha,0}$ \dotfill~& the germ $\left(\C(\!\{t\}\!),d+\frac{\alpha}{t}\right)$\\
  $\cN_{\alpha,d}$ \dotfill~& an elementary regular model\\
  $(\cE^\phi,\nabla)\otimes(\cR,\nabla)$ \dotfill~& a elementary meromorphic
    connection\\
  $(\cM^{nf},\nabla^{nf})$ \dotfill~& a model\\
  $\lambda$ \dotfill~& the isomorphism from definition~\ref{defn:model}\\
  $Q(t^{-1})$ \dotfill~& the irregular part of a system $[A]$\\
  $L\in \GL_n(\C)$ \dotfill~& the matrix of formal monodromy of $[A]$\\
  ${}^0C(\cM^{nf},\nabla^{nf})$ \dotfill~& all isomorphism classes of meromorphic
    connections, which are formally isomorphic to $(\cM^{nf},\nabla^{nf})$\\
  ${}^0C(A^0)$ \dotfill~& the system variant of ${}^0C(\cM^{nf},\nabla^{nf})$\\
  $\cH(\cM^{nf},\nabla^{nf})$ \dotfill~& the space of marked (meromorphic)
    pairs\\
  $\Syst_m(A^0)$ \dotfill~& the systems formally meromorphic equivalent to $A^0$\\
  $\cH(A^0)$ \dotfill~& the system version of $\cH(\cM^{nf},\nabla^{nf})$\\
  $\hat\Syst_m(A^0)$ \dotfill~& the marked pairs corresponding to $A^0$\\
\end{longtable}

\textbf{3.~Stokes Structures}
\begin{longtable}[h]{>{\raggedright}p{4cm}@{\hspace{.2cm}}p{10cm}}
  $\Lambda(A^0)$ \dotfill~& the Stokes sheaf (on $S^1$)\\
  $\St(A^0)$ \dotfill~& $H^1(S^1;\Lambda(A^0))$; the first non abelien cohomology of the
    Stokes sheaf\\
  $\exp$ \dotfill~& the Malgrange-Sibuya isomorphism\\
  $\exp_{A}$ \dotfill~& the Malgrange-Sibuya isomorphism corresponding to the system
    $[A]$\\
  $\cG$ \dotfill~& the sheaf of \rewrite{flat functions}\\
  $\Theta$ \dotfill~& $G\llbracket t\rrbracket/G\{t\}\to H^1(S^1;\cG)$; the isomorphism
    from the other Malgrange-Sibuya theorem\\
  $\cQ(A^0)$ \dotfill~& the set of all determining polynomials of $[\End A^0]$\\
  $a_{jl}\in\C\backslash\{0\}$ \dotfill~& the leading factor of
    $q_j-q_l\in\cQ(A^0)$\\
  $q_{jl}(t^{-1})$ \dotfill~& the leading coefficient of
    $q_j-q_l\in\cQ(A^0)$\\
  $k_{jl}$ \dotfill~& the degree of $q_j-q_l\in\cQ(A^0)$\\
  $\cK$ \dotfill~& $\{k_1<\dots<k_r\}$; the set of all levels of a system.\\
  $q_j \underset{\theta}{\prec} q_l$ \dotfill~& the first relation from
    definition~\ref{defn:definingRelations}\\
  $q_j\myrel{\theta}q_l$ \dotfill~& the second relation from
    definition~\ref{defn:definingRelations}\\
  $\Re(z)$ \dotfill~& the real part of a complex number $z$\\
  $\Im(z)$ \dotfill~& the imaginary part of a complex number $z$\\
  $\arg(z)$ \dotfill~& the \rewrite{argument} of a complex number $z$\\
  $\A$ \dotfill~& $\{\alpha_1,\dots,\alpha_\nu\}$; the set of all anti-Stokes
    directions\\
  $\S$ \dotfill~& the set of all Stokes directions\\
  $\Sto_\theta(A^0)$ \dotfill~& the Stokes group of $A^0$ in direction $\theta$ whose
    elements are Stokes germs\\
  $\vartheta_\theta$ & \\
  $\bdelta_{jl}$ \dotfill~& a block matrix version of Kronecker's delta,
    corresponding to the structure of $Q$\\
  $\SSto_\theta(A^0)$ \dotfill~& the group of all Stokes matrices of $A^0$ in direction
    $\theta$\\
  $\hat\SSto_\theta(A^0)$ \dotfill~& \TODO\\
  $\hat\vartheta_\theta$ & \\
  $\rho_\theta$ \dotfill~& $\Sto_\theta(A^0)\to\SSto_\theta(A^0)$; the map from
    proposition~\ref{prop:representation}\\
  $C_{\phi_\theta}$ \dotfill~& $\rho_\theta(\phi_\theta)$; the Stokes matrix
    corresponding to $\phi_\theta$\\
  $\Lambda^{k}(A^0)$ \dotfill~& the subsheaf of $\Lambda(A^0)$ of all germs,
    which are generated by $k$-germs\\
  $\Lambda^{\leq k}(A^0)$ \dotfill~& the subsheaf of $\Lambda(A^0)$ generated by
    $k'$-germs for all $k'\leq k$\\
  $\Lambda^{<k}(A^0)$ \dotfill~& the subsheaf of $\Lambda(A^0)$ generated by
    $k'$-germs for all $k'<k$\\
  $\Lambda^{\geq k}(A^0)$ \dotfill~& the subsheaf of $\Lambda(A^0)$ generated by
    $k'$-germs for all $k'geq k$\\
  $\Sto_\theta^\star(A^0)$ \dotfill~& $\Sto_\theta(A^0)\cap\Lambda_\theta^{\star}(A^0)$;
    \rewrite{the restriction of the Stokes sheaf for
    $\star\in\{k,<k,\leq k,\dots\}$}\\
  $\SSto^\star_\theta(A^0)$ \dotfill~& the groups of representations, which
    correspond to elements of $\Sto^\star_\theta(A^0)$\\
  $\A^k$ \dotfill~& the set of anti-Stokes directions bearing the level $k$\\
  $\A^{\leq k}$ \dotfill~& $\underset{k'\leq k}\bigcup\A^{k'}$\\
  $\A^{<k}$ \dotfill~& $\underset{k'<k}\bigcup\A^{k'}$\\
  $\A^{\geq k}$ \dotfill~& $\underset{k'\geq k}\bigcup\A^{k'}$\\
  $\cK_\alpha$ \dotfill~& the set of levels beared by $\alpha\in\A$\\
  $i_\alpha$ \dotfill~& $\Sto_\alpha(A^0)\to\prod_{k\in\cK_\alpha}\Sto_\alpha^k(A^0)$;
    the isomorphism corresponding to a factorization by a given order in a
    semidirect product\\
  $\rho_{\alpha}^k$ \dotfill~& $\Sto^k_\alpha(A^0)\to\SSto^k_\alpha(A^0)$;
    the \rewrite{restriction} of the map $\rho_{\alpha}$
    to the level $k$.\\
  $i_\alpha$ \dotfill~& $\SSto_\alpha(A^0)\to\prod_{k\in\cK_\alpha}\SSto_\alpha^k(A^0)$;
    the Stokes matrix version of the $i_\alpha$ above\\
  $\dot\cU$ \dotfill~& the nerve of the covering $\cU$\\
  $h$ \dotfill~& $\prod_{\alpha\in\A}\Sto_\alpha(A^0) \longrightarrow \St(A^0)$; the
    isomorphism from theorem~\ref{thm:mainThm2}\\
  $\cU^\star$ \dotfill~& for $\star\in\{k,<k,\leq k\}$ the adequate coverings defined in
    section~\ref{sec:proofOfMatrixThm}\\
  $\Gamma(\dot U^\star;\Lambda^\star(A^0))$ \dotfill~&
    $\prod_{\alpha\in\A^\star}\Gamma(\dot U_\alpha^\star;\Lambda^\star(A^0))$
    for every $\star\in\{k,<k,\leq k,\dots\}$\\
  $s^k$ \dotfill~& $\Gamma(\dot\cU^k;\Lambda^k(A^0))
    \to H^1(S^1;\Lambda^k(A^0))$; the quotient map\\
  $\sigma^k$ \dotfill~& $\Gamma(\dot\cU^k;\Lambda^k(A^0))
    \to\Gamma(\dot\cU^{\leq k};\Lambda^{\leq k}(A^0))$; an injective map
    defined in definition~\ref{defn:firstSetOfInlusions}\\
  $\sigma^{<k}$ \dotfill~& $\Gamma(\dot\cU^{<k};\Lambda^{<k}(A^0))
    \to\Gamma(\dot\cU^{\leq k};\Lambda^{\leq k}(A^0))$; an injective map
    defined in definition~\ref{defn:firstSetOfInlusions}\\
  $\mathfrak{S}^{\leq k}$ \dotfill~& $\Gamma(\dot\cU^{<k};\Lambda^{<k}(A^0)) \times
    \Gamma(\dot\cU^{k};\Lambda^{k}(A^0))
    \to \Gamma(\dot\cU^{\leq k};\Lambda^{\leq k}(A^0))$; the product map of
    cocycles\\
  $\alpha^+$ \dotfill~& the nearest anti-Stokes direction on the right of $\alpha$\\
  $\tau$ \dotfill~& $\prod_{k\in\cK}\Gamma(\dot \cU^k;\Lambda^k(A^0))\to
    \Gamma(\dot\cU;\Lambda(A^0))$; the product map of single-leveled cocycles\\
  $\cT$ \dotfill~& $\prod_{k\in\cK}\Gamma(\dot\cU^k;\Lambda^k(A^0))\to
    H^1(\cU;\Lambda(A^0))$; the isomorphism induced by $\tau$ on the
    cohomology\\
  $\mathfrak{T}$ \dotfill~& a map introduced in the proof of
    theorem~\ref{thm:mainThm2}\\
  % Weiter ab 4.3.2
\end{longtable}

\fi
\KOMAoptions{open=right}

%%%%%%%%%%%%%%%%%%%%%%%%%%%%%%%%%%%%%%%%%%%%%%%%%%%%%%%%%%%%%%%%%%%%%
\ifnum\myDevelopVariable=0
  \mainmatter
\fi

\ifnum\myHideVariable=0
%%%%%%%%%%%%%%%%%%%%%%%%%%%%%%%%%%%%%%%%%%%%%%%%%%%%%%%%%%%%%%%%%%%%%%%%%%%%%%%
\chapter{Poincaré asymptotic expansions}
\begin{comment}
  \begin{multicols}{2}
    ´classical'
    \begin{itemize}
      \item \cite[60]{sabbah_cimpa90} Chapter II.2.2
        \begin{itemize}
          \item \cite{sabbah2000equations}
        \end{itemize}
      \item \textbf{\textcolor{blue}{Van der Put:
        \cite[Chapter 7]{van2003galois}: Exact Asymptotics}}
      \item \cite{majima1984asymptotic}
      \item \cite[Sec.4.4]{Balser2000Formal}
      \item \cite{Loday1994}
      \item \textbf{\textcolor{blue}{\cite{Loday2014} Chapter 2}}
      \item \cite{hsieh2012basic}
    \end{itemize}
    \columnbreak
    ´sheafical'
    \begin{itemize}
      \item \textbf{\textcolor{red}{\cite[II.5]{sabbah2007isomonodromic}}}
    \end{itemize}
  \end{multicols}
  \TODO[\cite{sibuya1990Linear} Appendix A.3]
\end{comment}
In this chapter, we want to look at Poincaré asymptotic expansions (at $t_0$),
or for short, asymptotic expansions.
There are \rewrite{various literature references} for asymptotic expansions.
We will mostly refer to Loday-Richaud's book~\cite[chapter 2]{Loday2014},
\cite[XI-1-13]{hsieh2012basic} from Hsieh and Sibuya and the
book~\cite[chapter 7]{van2003galois} written by van der Put and Singer although
this topic is scratched in many publications.
We will assume that $t_0=0$ and this is without loss of generality, since
asymptotic expansions at $t_0\in\C$ reduce to asymptotic expansions at $0$
after the change of variable $t\mapsto s=t-t_0$ and asymptotic expansions at
$t_0=\infty$ are obtained by $t\mapsto s=\frac{1}{t}$.

We introduce the following terminology:
\begin{defn} We denote by
  \begin{itemize}
    \item $\mathfrak{s}_{a,b}(r)$, $a,b\in\S^1$ and $r\in\R_{>0}$ the open
      sector
      \[
        \mathfrak{s}_{a,b}(r):=\left\{t\in\C \mid a<\arg(t)<b, 0<|t|<r\right\}
      \]
      of all points $t\in\C$ satisfying $a<\arg(t)<b$ and $0<|t|<r$;
    \item $\mathfrak{s}_{I}(r)=\mathfrak{s}_{a,b}(r)$ for an open
      arc\footnote{i.e.\ an open interval of $S^1$} $I=(a,b)$;
    \item $\bar{\mathfrak{s}}_{a,b}(r)$ the closure of
      $\mathfrak{s}_{a,b}(r)$ in $\C^*:=\C\backslash\{0\}$.
  \end{itemize}
  The notion of sectors can easily be extended to sectors on a $q$-sheet cover
  or sectors on the Universal cover of $\C^*$, but we \rewrite{will not use
  them}.
\end{defn}
\begin{figure}[h]
  \begin{center}
    \begin{tikzpicture}[scale=4]
      \clip (-1.1,-0.2) rectangle (1.3,1.3);
      \node (zero) at (0,0) {};
      \node[below left] at (zero) {$0$};
      \draw[blue,dashed] (zero) circle (1cm);
      \node[blue] at (-0.9,0) {$S^1$};

      \filldraw[fill=green!20!white
        ,draw=green!60!black
        ,thick
        ,path fading=west,fading angle=110] (0,0)
      -- ({cos( 70 )*.55},{sin( 70 )*.55}) arc (70:150:.55) -- cycle;
      \node[green!60!black] at (-.47,.57) {$\mathfrak{s}_{a,b} (r)$};

      \draw[thick,red!60!black] ({cos( 70 )},{sin( 70 )}) arc
        (70:150:1)
        node[midway,above] {$I$};

      \draw[path fading=west] ({cos( 70 )*.55},{sin( 70 )*.55}) --
        ({cos( 70 )},{sin( 70 )});
      \draw[path fading=east] ({cos( 150 )*.55},{sin( 150 )*.55}) --
        ({cos( 150 )},{sin( 150 )});

      \fill[red!60!black] ({cos( 70 )},{sin( 70 )}) circle(.7pt);
      \fill[white] ({cos( 70 )},{sin( 70 )}) circle (.3pt);
      \fill[red!60!black] ({cos( 150 )},{sin( 150 )}) circle(.7pt);
      \fill[white] ({cos( 150 )},{sin( 150 )}) circle (.3pt);

      \node at ({1.1 * cos(150)},{1.1 * sin(150)}) {$b$};
      \node at ({1.1 * cos(70)},{1.1 * sin(70)}) {$a$};

      \filldraw[fill=blue!20!white
        ,draw=blue!60!black
        ,thick
        ,path fading=west,fading angle=52] (0,0)
      -- ({cos( 10 )*.45},{sin( 10 )*.45}) arc (10:95:.45) -- cycle;
      \draw[path fading=west] ({cos( 52.25 )*.45},{sin( 52.25 )*.45}) --
        ({cos( 52.25 )},{sin( 52.25 )}) node[above right] {$\theta$};
      \fill[red!60!black] ({cos( 52.25 )},{sin( 52.25 )}) circle(.7pt);
      \fill[white] ({cos( 52.25 )},{sin( 52.25 )}) circle (.3pt);
      \node[blue!40!black] at (.4,-0.05)
        {$U(\theta,\epsilon)$};
      \draw[blue!60!black
           ,decorate
           ,decoration={brace,mirror,amplitude=10pt}]
        ({cos( 10 ) * .45},{sin( 10 ) * .45})
        to node[midway,above right,xshift=6pt,yshift=-3pt] {$\frac{\epsilon}{2}$}
        ({cos( 52.25 ) * .45},{sin( 52.25 ) * .45});

      \fill[white] (zero) circle (1.5pt);

      \draw[green!60!black
           ,decorate
           ,decoration={brace,amplitude=10pt}]
        (0,0) -- ({cos( 150 )*.55},{sin( 150 )*.55})
          node[midway,below,xshift=-6.5pt,yshift=-6.5pt]
          {$r$};

      \fill (zero) circle (.7pt);
    \end{tikzpicture}
  \end{center}
  \caption{An exemplary arc $\textcolor{green!40!black}{
      \mathfrak{s}_{a,b}(\textcolor{blue!60!black}{r})}$ spanning from $a$ to
    $b$ with radius \textcolor{green!60!black}{$r$} corresponding to the arc
    \textcolor{red!60!black}{$I$} and an arc
    \textcolor{blue!60!black}{$U(\theta,\epsilon)$} for some $\theta$ and
    $\epsilon$.}
\end{figure}
We will say that \emph{a sector $\mathfrak{s}_{I}(r)$ contains $\theta\in S^1$}
if $\theta\in I$. In the same way we will say that $\mathfrak{s}_I(r)$ contains
$U\subset S^1$ if $U\subset I$ or that $U\subset S^1$ contains a sector
$\mathfrak{s}_I(r)$ if $I\subset U$.
\begin{rem}\label{rem:arcsByWidthAndCenter}
  Since the explicit value $r$ of the radius does not matter as long as it is
  small enough, we shall simply speak of an \emph{arc} or \emph{interval} $I$
  as $\mathfrak{s}_I(r)$ for a small enough $r$.
  It will be sometimes convenient to talk of arcs
  \[
    U(\theta,\epsilon):=\left(\theta-\frac{\epsilon}{2}
    ,\theta+\frac{\epsilon}{2}\right)\subset S^1
  \]
  defined by the midpoint together with the opening.
\end{rem}
\begin{defn}
  A sector $\mathfrak{s}_{a',b'}(r')$ is said to be a proper sub-sector of the
  sector $\mathfrak{s}_{a,b}(r)$,
  $\mathfrak{s}_{a',b'}(r')\Subset\mathfrak{s}_{a,b}(r)$, if its closure
  $\bar{\mathfrak{s}}_{a',b'}(r')$ is included in $\mathfrak{s}_{a,b}(r)$.
  \begin{s-rem}
    $\mathfrak{s}_{a',b'}(r')\Subset\mathfrak{s}_{a,b}(r)$ if and only if
    $(a',b')$ is a real sub-interval of $(a,b)$ and $r'<r$.
  \end{s-rem}
\end{defn}
%%%%%%%%%%%%%%%%%%%%%%%%%%%%%%%%%%%%%%%%%%%%%%%%%%%%%%%%%%%%%%%%%%%%%%%%%%%%%%%
\section{Definition and basic properties}
Let us now give the definition of asymptotic expansions.
\begin{defn}
  \marginnote{\cite[Defn.7.1]{van2003galois}}
  A function $f\in\cO(\mathfrak{s})$ is said to have the
  formal\footnote{The term ``formal'' emphasizes that we do not restrict the
  coefficients $a_n\in\C$ in any way.} Laurent series
  $\hat g=\sum_{n\geq n_0}a_nt^n$ as asymptotic expansion (or \emph{to be
  asymptotic to the series}) on the sector $\mathfrak{s}$ if
  \begin{itemize}
    \item for all proper sub-sectors $\mathfrak{s}'\Subset\mathfrak{s}$ and
    \item all $N\in\N$,
  \end{itemize}
  there exists a constant $C(N,\mathfrak{s}')$ such that the following
  estimate\footnote{Sometimes, for example in~\cite{sabbah_cimpa90}, this
    condition is written as
    \[
      \lim_{z\to0,z\in{\mathfrak{s}'}}
      |t|^{-(N-1)}
      \bigl|
        f(x)-\sum_{n_0\leq n\leq N-1}a_nt^n
      \bigr|=0
      \qquad \text{ for all } t\in \mathfrak{s}' \,.
    \]} holds
  \[
    \bigl|
      f(x)-\sum_{n_0\leq n\leq N-1}a_nt^n
    \bigr|
    \leq C(N,\mathfrak{s}')|t|^{N} \qquad \text{ for all }t\in \mathfrak{s}' \,.
  \]
\end{defn}
We will say that \emph{$f$ has the formal Laurent series $\hat g$ as asymptotic
  expansion on the interval $I\subset S^1$} if there exists a radius
$r\in\R_{>0}$ such that $f$ has the formal Laurent series $\hat g$ as asymptotic
expansion on $\mathfrak{s}_I(r)$.

If $f$ has $\hat g$ as asymptotic expansion on the sector $\mathfrak{s}$ (resp.\
on the interval $I\subset S^1$), we denote that by $f\sim_{\mathfrak{s}}\hat g$
(resp.\ $f\sim_{I}\hat g$).\label{page:notationForAsymptoticExpansion}

\begin{thm}
  For $f\in\cO(\mathfrak{s})$, there is at most one formal Laurent series $\hat
  f$ which satisfies $f\sim_{\mathfrak{s}}\hat f$
  (cf.\ Theorem~\cite[Thm.XI-1-5]{hsieh2012basic}).
\end{thm}

The following proposition gives equivalent characterizations of asymptotic
expansions.
\begin{prop}
  Let $f\in\cO(\mathfrak{s})$ then are the statements
  \begin{enumerate}
  \item $f\sim_{\mathfrak{s}}\hat f=\sum_{n\geq n_0}a_nt^n$,
  \item the function $f$ is infinitely often differential at the origin and
    $f^{(n)}(0)=n!a_n$ for $n\geq0$ and
  \item all derivatives $f^{(n)}(t)$ are continuous at the origin and
    $\lim_{\substack{t\to0\\t\in\mathfrak{s}}}f^{(n)}(t)=n!a_n$
  \end{enumerate}
  equivalent (cf.~\cite[4.4.Prop.8]{Balser2000Formal}).
\end{prop}

\begin{defn}
  \begin{enumerate}
    \item Define $\cA(\mathfrak{s})$ as the set of holomorphic functions
      $f\in\cO(\mathfrak{s})$ which admit an asymptotic expansion at $0$ on
      $\mathfrak{s}$.
      \begin{s-rem}
        \begin{enumerate}
          \item
            \marginnote{\cite[II.2.2.4]{sabbah_cimpa90}}
            $\cA(\mathfrak{s})$ is a subring of $\cO(\mathfrak{s})$.
          \item
            \marginnote{\cite[II.2.2.4]{sabbah_cimpa90}}
            $\cA(\mathfrak{s})$ contains $\C(\!\{t\}\!)$ as a subfield.
        \end{enumerate}
      \end{s-rem}
    \item
      For $(a,b)=I$ define $\cA(a,b)=\cA(I)$ as the limit
      $\underset{r\to0}{\underrightarrow{\lim}}
      \cA(\mathfrak{s}_{a,b}(r))$\footnote{In more detail: the
        elements of $\cA(a,b)$ are pairs $(f,\mathfrak{s}_{a,b}(r))$ with
        $f\in\cA(\mathfrak{s}_{a,b}(r))$. The equivalence relation is given by
        $(f_1,\mathfrak{s}_{a,b}(r_1))\sim(f_2,\mathfrak{s}_{a,b}(r_2))$ if
      there is a pair $(f_3,\mathfrak{s}_{a,b}(r_3))$ such that
      $r_3<\min(r_1,r_2)$ and $f_3=f_1=f_2$ on $\mathfrak{s}_{a,b}(r_3)$.}.
  \end{enumerate}
\end{defn}
It can easily be seen that $\cA:U\to\cA(U)$ defines a sheaf on $S^1$.

One writes
\begin{align*}
  T_{I}:\cA(I)&\to\C(\!(t)\!)
  \\f&\mapsto\hat f
\end{align*}
for the so-called \emph{Taylor map} which associates to each function
$f\in\cA(I)$ its asymptotic expansion.
If $\hat f$ is not defined to be something else, we write $\hat f$ for the image
of $f$ under $T_{I}$.
\begin{defn}
  A function $f$ which is asymptotic to the identity, i.e.\ with
  $T_{I}(f)=\hat f=\id_{\mathfrak{s}_I}$, is called \emph{flat (on $I$)}.
\end{defn}

The kernel of $T_{I}$ is not zero in general.
For example is the function $e^{-\frac{1}{t}}$ asymptotic to zero on every
$I\subset\left(-\frac{\pi}{2},\frac{\pi}{2}\right)$.
\TODO[changed from $e^{-\frac{1}{\sqrt{t}}}$ after hint from giovanni]
\begin{defn}
  Denote by $\cA^{<0}(I)$ the kernel $\ker(T_{I})\subset\cA(I)$ of functions
  asymptotic to zero at $0$. This does also define a sheaf $\cA^{<0}$.
\end{defn}

\begin{exmp}[Trivial example]
  \marginnote{\cite[Exmp.2.2.3]{Loday2014}}
  If $f$ is a holomorphic function on $D=\{t\mid |t|<r\}$ then $f$ is asymptotic
  to its Taylor series at $0$ on every arc
  $\mathfrak{s}\subset\mathfrak{s}_{S^1}(r)=D\backslash\{0\}$.

  Reciprocally, if $f$ is holomorphic on $\mathfrak{s}_{S^1}(r)$ and has an
  asymptotic expansion at $0$ on $\mathfrak{s}_{S^1}(r)$ then, $f$ is bounded
  near $0$ and according to the removable singularity theorem, $f$ is
  meromorphic on $D$.
  \begin{s-rem}
    This implies that $\cA(S^1)\cong\C(\!\{t\}\!)$.
  \end{s-rem}
\end{exmp}
For many more extensively discussed examples, like the Euler function or the
exponential integral, see Loday-Richaud's book \cite[Sec.2.2]{Loday2014}.

\begin{prop}\label{prop:optionsOfAsymptoticExp}
  \marginnote{\cite[II.2.2.4]{sabbah_cimpa90}i
             ,\cite[XI-1-11]{hsieh2012basic}}
  For every arc $I\subset S^1$ is $\cA(I)$ is stable under derivation, i.e.\
  $f'(t)$ is on $I$ asymptotic to $\hat f'(t)$ for every function $f$
  asymptotic to $\hat f$ on $I$.
\end{prop}
\begin{proof}
  Let $f\in\cA(I)$ with asymptotic expansion $\hat f=\sum_{-n_0\leq n}a_nx^n$
  its asymptotic expansion.
  One has for each $m\geq0$
  \[
    f(t)=\sum_{n=0}^ma_nt^n+R_m(t)t^m
  \]
  with $\underset{\substack{t\to0\\t\in\mathfrak{s}_I(r)}}{\lim}R_m(t)=0$.
  This implies that $R_m$ is holomorphic in $\mathfrak{s}_I(r)$.
  Thus one has
  \[
    f'(t)=\sum_{n=0}^mna_nt^{n-1}+mR_m(t)t^{m-1}+R_m'(t)t^m
  \]
  Let $C_{x,\rho}$ be the circle at $x$ with radius $\rho$ contained in
  $\mathfrak{s}_I(r)$. The Cauchy theorem \rewrite{implies then that}
  \[
    |R_m'(t)|\leq\frac{1}{\rho}\max_{s\in\C_{x,\rho}}|R(s)| \,.
  \]
  Let $J$ be a relatively compact open set in $I$, then there exists a positive
  number $\alpha$ such that for all $x\in\mathfrak{s}_I(r)$ one has
  $C_{x,\alpha|x|}\subset\mathfrak{s}_I(r)$.
  \TODO{}
  This proves that $\hat{f'}$ is an asymptotic expansion for $f'$ in
  $\mathfrak{s}_I(r)$.
\end{proof}
Other properties can be found in the books~\cite{Balser2000Formal} from Balser
and~\cite{hsieh2012basic} from Hsieh and Sibuya. Some are stated in the
following remark.
\begin{rem}
  Let $f_1$ and $f_2$ be functions which satisfying $f_1\sim_I\hat f_1$ and
  $f_2\sim_I\hat f_2$.
  We have also the rules:
  \begin{enumerate}
    \item $f_1(t)+f_2(t)$ is on $I$ asymptotic to $\hat f_1(t)+\hat f_2(t)$
      (cf.\ Theorem~\cite[4.5.Thm.13]{Balser2000Formal}
      or~\cite[XI-1-6]{hsieh2012basic}).
    \item $f_1(t)f_2(t)$ is on $I$ asymptotic to $\hat f_1(t)\hat f_2(t)$
      (cf.\ Theorem~\cite[4.5.Thm.14]{Balser2000Formal}
      or~\cite[XI-1-6]{hsieh2012basic}).
    \item $\int_0^tf(s)\,ds$ is on $I$ asymptotic to $\int_0^t\hat f(t)\,ds$
      (cf.\ Theorem~\cite[4.5.Thm.20]{Balser2000Formal}).
    \item If $f\in\cA(I)\backslash\cA^{<0}(I)$ then $f^{-1}(t)$ is on $I$
      asymptotic to $\hat f^{-1}(t)$
      (cf.\ Theorem~\cite[4.5.Thm.21]{Balser2000Formal}
      or~\cite[XI-1-9]{hsieh2012basic}).
  \end{enumerate}
\end{rem}

\newpage%%%%%%%%%%%%%%%%%%%%%%%%%%%%%%%%%%%%%%%%%%%%%%%%%%%%%%%%%%%%%%%%%%%%%%
% \vspace{-3pt}
\section{Borel-Ritt Lemma}
\begin{comment}
  \begin{itemize}
    \item \cite[Lem.II.2.2.5]{sabbah_cimpa90}: $T_I$
    \item \textbf{\cite[Th.7.3]{van2003galois}}: $T_I=T_{(a,b)}$
    \item \cite[Th.2.4.1]{Loday2014}: $T_{\mathfrak{s}}$
    \item \cite[4.4.Thm.16]{Balser2000Formal}
    \item \cite[XI-1-13]{hsieh2012basic}
  \end{itemize}
\end{comment}
The most important theorem here is the Borel-Ritt Lemma.
\begin{thm}[Borel-Ritt]\label{thm:borel-ritt}
  \marginnote{\cite[Thm.2.4.1]{Loday2014}}
  Let $I\subsetneq S^1$ be an open sub-arc, i.e.\ an arc with a opening less
  than $2\pi$, then is the Taylor map
  $T_{I}:\cA(I)\twoheadrightarrow\C(\!(t)\!)$ onto.
\end{thm}
\marginnote{\cite[111]{Loday2014}
  \\,~\cite[XI-1-13]{hsieh2012basic}}
In other words does this theorem say, that to each formal Laurent series
 $\hat f\in\C(\!(t)\!)$ and every arc $I\subsetneq S^1$ there are functions $f$
asymptotic to $\hat f$ on $I$.
There are multiple approaches to obtain such a function $f\in\cA(I)$ asymptotic
to a given function $\hat f\in\C(\!(t)\!)$ in a canonical way.
For example (multi-)multisummability which is in depth described in
Loday-Richaud's book~\cite{Loday2014}.
This topic is also discussed in the paper~\cite{Martinet1991} from Martinet and
Ramis.

However, when $\hat f$ satisfies an equation, these asymptotic functions do not
necessarily satisfy the same equation in general.
This gap will be filled in Section~\ref{sec:MAET} by the main asymptotic
existence theorem for small enough sectors.

Let us now proof the Borel-Ritt Lemma.
\begin{proof}[Proof of Theorem~\ref{thm:borel-ritt}]
  Let $I:=(-\pi,\pi)$ and $R\in\R_{>0}$.
  We will prove this for the sector
  \[
    \mathfrak{s}:=\mathfrak{s}_{I}(R)=\left\{
      t\in\C\mid |\arg(t)|<\pi, 0<|t|<R
    \right\}
  \]
  in which, after rotation and for large enough $R$, every sector
  $\mathfrak{s}''\subsetneq\C^*$ lies.
  Let $\sum a_nz^n$ be a formal Laurent series. We look for a function
  $f\in\cA(I)$ with Taylor series $T_{I}f=\sum a_nz^n$.
  By subtracting the principal part\TODO[?] we may assume that this series has
  no terms with negative degree.
  Let $b_n$ be a sequence, which satisfies
  \begin{itemize}
    \item[] the series $\sum |a_n|b_nR^{n-\frac{1}{2}}$ is convergent.
  \end{itemize}
  For example\TODO[realy?], one may set
  \[
    b_n=\begin{cases}
      0                 & \text{when~} n=0
      \\\frac{1}{n!}|a_n| & \text{when~} n>0
    \end{cases}
    \,.
  \]
  Let $\sqrt{t}$ be the branch of the square root function that satisfies
  $|\arg(\sqrt{t})|<\frac{\pi}{2}$\footnote{For $t=re^{i\phi}$ is
  $\sqrt{t}=\sqrt{r}e^{i\frac{\phi}{2}}$ with $-\pi<\phi<\pi$.} for all
  $t\in\mathfrak{s}$.
  For any real number $b_n$, the function
  $\beta_n(t):=1-e^{-\frac{b_n}{\sqrt{t}}}$ satisfies
  \begin{enumerate}
    \item[(a)] $|\beta_n(t)|\leq\frac{b_n}{\sqrt{|t|}}$ since
      $1-e^t=-\int_0^te^s\,ds$ implies that $|1-e^t|<|t|$ for $\Re(t)<0$ and
    \item[(b)] $\beta_n$ has asymptotic expansion $1$ on $\mathfrak{s}$ (thus
      $\beta_n-1$ has asymptotic expansion $0$ on $\mathfrak{s}$).
  \end{enumerate}
  Define $f(t):=\sum a_n\beta_n(t)t^n$.
  Since
  \[
    |a_n\beta_n(t)t^n|
    \leq|a_n|b_n|z|^{n-\frac{1}{2}}
    \leq|a_n|b_nR^{n-\frac{1}{2}},
  \]
  the series
  $\sum a_n\beta_n(t)t^n$ converges and its sum $f(t)$ is in
  $\cO(\mathfrak{s})$.

  Consider a proper sub-sector $\mathfrak{s}'\Subset\mathfrak{s}$ and
  $t\in\mathfrak{s}'$. Then, for every $N>0$
  \begin{align*}
    \left|f(t)-\sum_{n=0}^{N-1} a_nt^n\right|
    &\leq \left| \sum_{n=0}^{N-1}a_n(\beta_n(t)-1)t^n \right|
      +|t|^N\sum_{n\geq N}\left| a_n\beta_n(t)t^{n-N} \right|
  \\&= \left| \sum_{n=0}^{N-1}a_ne^{-\frac{b_n}{\sqrt{t}}}t^n \right|
      +|t|^N\sum_{n\geq N}\left| a_n\beta_n(t)t^{n-N} \right|
  \end{align*}
  The first summand is a finite sum of terms, which are all asymptotic to $0$
  and thus is majorized by $C'|t|^N$ for a convenient constant $C'$.
  The second summand is majorized by
  \[
    |t|^N\left(2|a_n|+\sum_{n\geq N+1}|a_n|b_nR^{n-\frac{1}{2}-N}\right).
  \]
  By setting $C=C'+2|a_n|+\sum_{n\geq N+1}|a_n|b_nR^{n-\frac{1}{2}-N}$ we
  obtain a positive constant $C$, which depends on $N$ and the radius of
  $\mathfrak{s}'<R$, such that
  \[
    \left|f(t)-\sum_{n=0}^{N-1} a_nt^n\right| \leq C|t|^N
    \qquad\text{ for all } t\in\mathfrak{s}' \,.
  \]

  \TODO[The GENERAL CASE: see \cite{Loday2014} page 28]
\end{proof}

%%% Local Variables:
%%% TeX-master: "Maximilian_Huber-Masters_Thesis-with_notes.tex"
%%% End:

\fi

\ifnum\myHideVariable=0
\chapter{Systems and meromorphic connections}
Here we will start with the definition of meromorphic connections on holomorphic
bundles although we will only be interested in local information.
For local description we will use \textbf{germs of meromorphic connections}, the
coordinate dependent \textbf{systems} and \textbf{connection matrices}.
Other coordinate independent approaches arise for example from the theory of
(localized holonomic) $\cD$-modules.

Meromorphic connections are introduced and discussed in many resources.
A good starting point are Sabbah's lecture notes~\cite{sabbah_cimpa90}.
More advanced resources are for example
Sabbah's book~\cite{sabbah2007isomonodromic},
Varadarajan's book~\cite{Varadarajan96linearmeromorphic} or the
book~\cite{hotta2008} from Hotta et al.
The necessary facts about meromorphic connections are also stated in
Boalch's paper~\cite{boalch} (resp.\ his thesis~\cite{thboalch}),
and Loday-Richaud's paper~\cite{Loday1994}.

Although the language of meromorphic connections is often preferred, we will
use the language of systems most of the time.
Systems are, for example, discussed in the book~\cite{hotta2008} from Hotta et
al, Loday-Richaud's paper~\cite{Loday1994} and her book~\cite{Loday2014} and
Boalch's publications~\cite{boalch,thboalch}. Another resource \rewrite{might}
Remy's paper~\cite{Remy2014} be.

We will use \rewrite{all of the} above mentioned \rewrite{resources} in this
chapter.

\section{(Global) meromorphic connections}
Let $M$ be a riemanian surface and let $Z=k_1(a_1)+\cdots+k_m(a_m)>0$ be an
effective divisor\footnote{The $a_i\in M$ are distinct points and the $k_i$ are
positive integers.} on $M$.
It is sufficient to think $M=\P^1$ and $0\in|Z|$\footnote{If
$Z=k_1(a_1)+\cdots+k_m(a_m)$ then $|Z|:=\{a_1,\dots,a_m\}$.}, since we will
only be interested in local information (at $0$).

Let $\sM$ be a holomorphic bundle over $M$ i.e.\ a locally free $\cO_M$-module
of rank $n$.
A global meromorphic connection, i.e.\ a meromorphic connection on $\sM$, is
then defined as follows.
\begin{defn}\label{defn:mercon}
  A \emph{meromorphic connection $(\sM,\nabla)$ on $\sM$ with poles on $Z$} is
  is defined by a $\C$-linear morphism of sheaves
  \[
    \nabla:\sM\to\Omega_M^1(*Z)\otimes\sM
  \]
  satisfying, for each open subset $U\underset{\text{op.}}{\subset} M$, the
  \emph{Leibniz rule}
    \[
      \nabla(fs)=f\nabla s+(df)\otimes s
    \]
  for $s\in\Gamma(U,\sM)$ and $f\in\cO_M(U)$.
  The \emph{rank} of the meromorphic connection $(\sM,\nabla)$ is defined to be
  the rank of the Bundle $\sM$.
  It will usually denoted by the letter $n$.
  \begin{s-rem}
    Some authors use the factors $k_i$ of the divisor $Z$ to limit the pole
    orders at the points $a_i$. Since we do not need this restriction, we allow
    arbitrary pole orders. Denoted is this by the $*$ in $\Omega_M^1(*Z)$.

    The \emph{sheaf of meromorphic differential $1$-forms on $M$}
    $\Omega_M^1(*Z)$ is defined as
    \[
      \Omega_M^1(*Z):=\cO_M(*Z)\underset{\cO_M}\otimes\Omega_M^1 \,,
    \]
    where $\cO_M(*Z)$ is the sheaf of functions, which are meromorphic along
    $Z$.
    A detailed definition of these sheaves can be found in Sabbah's
    book~\cite{sabbah2007isomonodromic} in
    Section 0.8 and Section 0.9.b
  \end{s-rem}
\end{defn}
The word ``global'' used to emphasize that the connection is ``on $\sM$'', i.e.\
not only a germ at some point of $M$, \rewrite{like it will be used} later
(cf.\ Remark~\ref{rem:GlobalNothingGerm}).

In the definition above is the variant `holomorphic bundle with meromorphic
connection' chosen, like in Boalch's paper~\cite{boalch}.
There is also the interchanged description `meromorphic bundle with holomorphic
connection' which is for example used in Sabbah's
book~\cite{sabbah2007isomonodromic}.

By choosing a lattice of a meromorphic bundle
(cf.\ \cite[Def.0.8.3]{sabbah2007isomonodromic}), one gets a holomorphic bundle
but if the meromorphic bundle had a holomorphic connection, the induced
connection \rewrite{on the lattice} is no longer guaranteed to be holomorphic.
Thus we obtain a meromorphic connection on a holomorphic bundle in our sense.

We will only look at unramified connections. Some of the following statements
are only valid in the unramified case although we will not always mention this
restriction.

\begin{comment}
  \begin{defn}
    \marginnote{\cite[0.12.2]{sabbah2007isomonodromic}}
    The connection $\nabla:\sM\to \Omega_M^1\otimes_{\cO_M}\sM$ is said to be
    \emph{integrable} or \emph{flat}, if its curvature vanishes, i.e.\
    $R_\nabla\equiv0$
    where $R_\nabla:=\nabla\circ\nabla:\cE\to\Omega_M^2\otimes_{\cO_M}\cE$ is a
    $\cO_M$-linear morphism.
    \begin{s-rem}
      Here are all connections flat, since all connections will be of Dimension
      one.
    \end{s-rem}
  \end{defn}
\end{comment}

%%%%%%%%%%%%%%%%%%%%%%%%%%%%%%%%%%%%%%%%%%%%%%%%%%%%%%%%%%%%%%%%%%%%%%%%%%%%%%%
\section{Local expression of meromorphic connections and systems}
\marginnote{\cite[28]{sabbah2007isomonodromic},
  \\\cite[2]{thboalch} and
  \\\cite[11]{babbitt1989local}}
We will only be interested in local information of meromorphic connections.
This means that we look at a connection in a neighbourhood of $0$ which has its
unique singularity at $0$.
There are several ways of expressing the local information without the need of
an fixed neighbourhood.
We will either talk about germs (at $0$) of meromorphic connections or systems,
which depend on the choice of a trivialization.
\begin{prop}
  \marginnote{\textbf{\cite[Rem.5.2.4]{hotta2008}}\\\cite[Def.4.2.1]{Loday2014}}
  A germ of a meromorphic connection $(\sM,\nabla)$ is the sheaf-theoretic
  germ (at $t=0$) and thus is given by a tuple $(\cM,\nabla)$ where
  \begin{itemize}
    \item $\cM$ is the germ at $0$ of the holomorphic bundle $\sM$ and thus a
      $\C(\!\{t\}\!)$-vectorspace of dimension $n$, since the ring of germs of
      meromorphic functions with poles at $0$ is the ring $\C(\!\{t\}\!)$, and
    \item $\nabla:\cM\to \cM$ is a $\C$-linear map, which
      satisfies the \emph{Leibniz rule}
      \[
        \nabla(fm)=\frac{d}{dt} f\cdot m + f\nabla(m)
      \]
      for all $f\in\C(\!\{t\}\!)$ and $m\in \cM$.
  \end{itemize}
  \begin{s-rem}
    Loday-Richaud calls the germs of meromorphic connections in her
    book~\cite[Def.4.2.1]{Loday2014} \emph{differential modules}.
  \end{s-rem}
\end{prop}
\begin{rem}\label{rem:GlobalNothingGerm}
  From now on we will only talk about \textbf{germs of} meromorphic
  connections $(\cM,\nabla)$ and we will call them meromorphic connection. If
  we want to talk about meromorphic connections in the sense of
  Definition~\ref{defn:mercon} we will emphasize this by the word `global' or
  by talking about a meromorphic connection \textbf{on $M$}.
\end{rem}
It will occasionally be convenient to omit the $\nabla$ and simply call $\cM$
the meromorphic connection.
\begin{defn}
  \marginnote{\cite[Def.5.2.1]{hotta2008}}
  A \emph{(iso-)morphism of meromorphic connections}
  $\Phi:(\cM,\nabla)\overset{\sim}\to(\cM',\nabla')$ is a (iso-)morphism of
  $\C(\!\{t\}\!)$-vectorspaces $\Phi:\cM\overset{\sim}\to\cM'$ which commutes
  with the connections, i.e.\ which satisfies
  $\nabla\circ\phi = (\id\otimes\phi)\circ\nabla$.
\end{defn}
\marginnote{\cite[65]{Loday2014}, \cite[129]{hotta2008}}
Choose a $\C(\!\{t\}\!)$-basis $\underline{e}=(e_1,e_2,\dots,e_n)$ of $\cM$.
Let $A=(a_{jk})_{j,k\in\{1,\dots,n\}}$ be a $n\times n$ matrix with entries in
$\C(\!\{t\}\!)$ such that the it describes the action of $\nabla$ corresponding
to the chosen basis, i.e.\ satisfies
 $\nabla e_k=-\sum_{1\leq j\leq n} a_{jk}(t)e_j$ for every $k\in\{1,\dots,n\}$.
\begin{defn}
  The matrix $A$ is called a \emph{connection matrix} of $(\cM,\nabla)$.
\end{defn}
The connection $\nabla$ is fully determined by $A$. Indeed, let
$x=\sum_{0\leq j\leq n}x_je_j=\underline{e}\cdot X$ be an arbitrary element of
$\cM$, where $X={}^t\!(x_1,x_2 ,\dots,x_n)$ is a column matrix.
Applying $\nabla$ to $x$ using the Leibniz rule, yields
\begin{align*}
  \nabla x&=\nabla\left(\underline{e}\cdot X\right)
  \\&=\underline{e} \cdot dX + \nabla \underline{e} \cdot X
  \\&=\underline{e}\left(dX-AX\right).
\end{align*}
Such that horizontal sections of $(\cM,\nabla)$, i.e.\ sections which satisfy
$\nabla x=0$, correspond to solutions of
\begin{equation}\label{eq:ode}
  \frac{d}{dt}x=Ax \,.
\end{equation}
Thus, with the connection $\nabla$ and the chosen $\C(\!\{t\}\!)$-basis
$\underline{e}$, is naturally the differential operator
$\triangle=d-A$ associated, which has order one and dimension $n$.
\begin{defn}
  \marginnote{\cite[Sec.4]{Martinet1991}}
  We call (\ref{eq:ode}), determined by the differential operator
  $\triangle=d-A$, a \emph{germ of a meromorphic linear differential
  system\footnotemark} of rank $n$, or just a
  \emph{system}.
  \begin{s-prop}
    Thus, the set of systems is isomorphic to the set
    \[
      \End(E)\otimes\C(\!\{t\}\!)=\gl_n(\C(\!\{t\}\!))
    \]
    of all connection matrices.
  \end{s-prop}
  Such a system will be denoted by $[A]=d-A$ and we will call $A$ the
  connection matrix of the system $[A]$.
\end{defn}
  \footnotetext{Martinet and Ramis call them in~\cite{Martinet1991}
  \emph{germs of meromorphic differential operators.}}
\TODO[possibly multivalued solutions ($\tilde K$)?~\cite{hotta2008} on page
128]

If we start with a system $[A]$ and we want a meromorphic connection
$(\cM,\nabla)$ which has $A$ as connection matrix we can do this in the
following way.
\begin{prop}\label{prop:systToMeromConn}
  If we start with either a system $[A]$ of rank $n$, or a connection matrix
  $A\in\gl_n(\C(\!\{t\}\!))$, we get a germ of a meromorphic connection via
  \[
    (\cM_A,\nabla_A)=(\C(\!\{t\}\!)^n,d-A)
  \]
  which has $A$ as its connection matrix.
\end{prop}

\begin{prop}\label{prop:MatOfSumOfMerCon}
  Let $(\cM_1,\nabla_1)$ and $(\cM_2,\nabla_2)$ be two meromorphic
  connections
  with the connection matrices $A_1$ and $A_2$.
  A connection matrix of $(\cM_1,\nabla_1)\oplus(\cM_1,\nabla_1)$ is then
  given
  by the block-diagonal matrix $\diag(A_1,A_2)$.
  \begin{s-cor}\label{cor:MatOfSumOfMerCon}
    Let $(\cM,\nabla)$ be an meromorphic connection, whose connection matrix
    can be written as $\diag(A_1,A_2)$.
    Then one can decompose the meromorphic connection corresponding to the block
    structure of $A$, i.e.\ $(\cM,\nabla)$ is then isomorphic to the direct sum
    $(\cM_{A^1},\nabla_{A^1})\oplus(\cM_{A^2},\nabla_{A^2})$
  \end{s-cor}
\end{prop}
\begin{proof}
  Use Proposition~\ref{prop:systToMeromConn} to write the connection as
  \[
    (\C(\!\{t\}\!)^{n_1},d-A_1)\oplus(\C(\!\{t\}\!)^{n_2},d-A_2)
  \]
  and we want to show that it is isomorphic to
  \[
    \left(\C(\!\{t\}\!)^{n_1+n_2},d-
    \begin{pmatrix} A_1 & 0 \\0 & A_2 \end{pmatrix}\right) \,.
  \]
  Denote by $A:=\diag(A_1,A_2)$ the block-diagonal matrix build from $A_1$ and
  $A_2$.
  For every $j\in\{1,2\}$ we have the inclusion
  $i_j:\C(\!\{t\}\!)^{n_j}\hookrightarrow\C(\!\{t\}\!)^{n_1+n_2}$ and the
  diagram
  \[ \begin{tikzcd}
      \C(\!\{t\}\!)^{n_j} \rar{d-A_j}\dar{i_j} & \C(\!\{t\}\!)^{n_j} \dar{i_j}
    \\\C(\!\{t\}\!)^{n_1+n_2} \rar{d-A} & \C(\!\{t\}\!)^{n_1+n_2}
  \end{tikzcd} \]
  which commutes, since the derivation commutes with the inclusion and the
  matrix $A$ is build in the \rewrite{correct way}, to satisfy
  $i_j(A_jx)=A_j(i_j(x))$:
  \begin{align*}
    i_j(dx-A_jx) &= i_j(dx)-i_j(A_jx)
    \\&=d(i_j(x))-A_j(i_j(x))
    \\&=((d-A)\circ i_j)(x) \,.
  \end{align*}
\end{proof}

\begin{rem}
  Let $(\cM_1,\nabla_1)$ and $(\cM_2,\nabla_2)$ be meromorphic connections.
  Then $\cM_1\otimes\cM_2$ is endowed with the structure of a meromorphic
  connection by
  \[
  \nabla(u_1\otimes u_2)=\nabla_1u_1\otimes u_2+u_1\otimes\nabla_2u_2
  \]
  where $u_i\in\cM_i$ (cf.\ \cite[129]{hotta2008}).
  \begin{comment}
    $\Hom_{\C(\!\{t\}\!)}(\cM_1,\cM_2)$ is endowed with the structure of a
    meromorphic connection by
    \[
    (\nabla\phi)(u_1)=\nabla_2(\phi(u_1))-\phi(\nabla_1 u_1)
    \]
    where $\phi\in\Hom_{\C(\!\{t\}\!)}(\cM_1,\cM_2)$ and $u_i\in\cM_i$.
  \end{comment}
\end{rem}

\begin{comment}
  \begin{lem}
    \PROBLEM[how looks the connection matrix? Simply the sum?]
    Let $(\cM_1,\nabla_1)$ and $(\cM_2,\nabla_2)$ be meromorphic connections
    with connection matrices $A_1$ and $A_2$.
    The connection matrix of $(\cM_1\otimes\cM_2,\nabla)$ is then given by
    $A_1+A_2$.
  \end{lem}
  \begin{proof}
    Let $u_1\otimes u_2\in\cM_1\otimes\cM_2$.
    \begin{align*}
      \nabla(u_1\otimes u_2) &=\nabla_1u_1\otimes u_2+u_1\otimes\nabla_2u_2
      \\&=(\nabla_1u_1)\otimes u_2+u_1\otimes\nabla_2u_2
    \end{align*}
  \end{proof}
\end{comment}

\begin{comment}
%%%%%%%%%%%%%%%%%%%%%%%%%%%%%%%%%%%%%%%%%%%%%%%%%%%%%%%%%%%%%%%%%%%%%%%%%%%%%%%
  \subsubsection{Formalization}
  \begin{multicols}{2}
    Let $[A]$ be a system. We \rewrite{view it as a formal system, by} allowing
    formal solutions.
    \TODO{}

    \columnbreak

    Let $(\cM,\nabla)$ be a meromorphic connection. The connection $\nabla$
    naturally extends to $\hat\cM:=\cM\otimes\C(\!(t)\!)$ and
    $\tilde\cM_\theta:=\cM\otimes\cA_\theta$.
    \TODO{}
  \end{multicols}
\end{comment}


\begin{comment}
%%%%%%%%%%%%%%%%%%%%%%%%%%%%%%%%%%%%%%%%%%%%%%%%%%%%%%%%%%%%%%%%%%%%%%%%%%%%%%%
  \subsubsection{As differential operator}
  \marginnote{\cite[Sec.4.2]{Loday2014}}

  From the theory of ordinary differential equations we know that to
  (\ref{eq:ode}) there is a equivalent ordinary differential equation of order
  $n$ which can be written as
  \[
    \underset{=:P}{\underbrace{%
        (a_n\partial_t^n+a_{n-1}\partial_t^{n-1}+\cdots a_{1}\partial_t+a_{0})
    }} \cdot v=0
  \]
  where $a_i\in\C(\!\{t\}\!)$.
  This leads to the theory of $\cD$-modules.

  \marginnote{See \cite[Sec.1.4]{babbitt1983} for \textbf{ode of rank $n$} to
  \textbf{system}.}
\end{comment}

%%%%%%%%%%%%%%%%%%%%%%%%%%%%%%%%%%%%%%%%%%%%%%%%%%%%%%%%%%%%%%%%%%%%%%%%%%%%%%%
\subsection{Transformation of systems}
% \begin{comment}
%   \begin{notations}
%     We will use the following notations
%     \begin{itemize}
%     \item $G[t]=\Gl_n(\C[t])$;
%     \item $G\{t\}=\Gl_n(\C\{t\})$ the analytic transformations;
%     \item $G(\!\{t\}\!)=\Gl_n\left(\C\{t\}[t^{-1}]\right)$ the
%       meromorphic\footnotemark{} transformations;
%     \item $G\llbracket t\rrbracket=\Gl_n\left(\C\llbracket t\rrbracket\right)$
%       the maybe not applicable formal transformations;
%     \item $G(\!(t)\!)=\Gl_n\left(\C\llbracket t\rrbracket[t^{-1}]\right)$
%       the maybe not applicable formal meromorphic transformations.
%     \end{itemize}
%   \end{notations}
% \end{comment}
\marginnote{\cite[Sec.4.3.1]{Loday2014}}

By \emph{meromorphic\footnote{We use the term meromorphic in the sens of
convergent meromorphic. Otherwise we say formal meromorphic.} transformation},
or just \emph{transformation}, of a system we mean a $\C(\!\{t\}\!)$-linear
change of the trivialization.
Such a change is given by a matrix $F$ in
$G(\!\{t\}\!):=\Gl_n\left(\C\{t\}[t^{-1}]\right)$ and the transformed connection
matrix ${}^F\!A$ is obtained through the Gauge transformation
\begin{equation}\label{eq:gaugeTrafo}
  {}^F\!A=(dF)F^{-1} + FAF^{-1} \,.
\end{equation}
Let $B:={}^F\!A$ be the transformed matrix and rewrite the equation
(\ref{eq:gaugeTrafo}) to obtain the linear differential equation
\[
    \frac{dF}{dt}=BF-FA
\]
which will be denoted by $[A,B]$.
\begin{rem}
  The system matrix $B$ is obtained from $A$ by transformation $F$ if and only
  if $F$ solves the linear differential system $[A,B]$.
\end{rem}

If $F$ is formal, i.e.\
$F\in G(\!(t)\!):=\Gl_n\left(\C\llbracket t\rrbracket[t^{-1}]\right)$, it will
usually be denoted by $\hat F$.
The transformation of $A$ by $\hat F\in$ G(\!(t)\!) is not guaranteed to have
convergent entries and thus is no system.
We denote by $\hat G(A)$ the set of all \emph{(applicable) formal
transformations}
\[
  \hat G(A):=\bigl\{\hat F\in G(\!(t)\!)
    \mid {}^{\hat F}\!A
    \text{ has convergent entries, i.e.\ ${}^{\hat F}\!A\in G(\!\{t\}\!)$}
  \bigr\}\,.
\]
Let $\hat{F'}\in\hat G(A)$ and $A':={}^{\hat{F'}}\!A$, then are the sets
$\hat G(A)$ and $\hat G(A')$ related by
\[
  \hat G(A')=\hat G(A)\hat{F'}^{-1}=\bigl\{
    \hat F\in G(\!\{t\}\!) \mid \hat F\hat{F'}\in\hat G(A)
  \bigr\} \,.
\]

\begin{rem}
  The Gauge transformations arise naturally from an isomorphism of meromorphic
  connections in the following way.
  If we start with
  \begin{itemize}
  \item two base choices $\cM\overset{\sim}\longrightarrow\C(\!\{t\}\!)^n$
    and $\cM'\overset{\sim}\longrightarrow\C(\!\{t\}\!)^n$ and
  \item an isomorphism
    $\Phi:(\cM,\nabla)\overset{\sim}\longrightarrow(\cM',\nabla')$ together
    with the corresponding base change $F\in G(\!\{t\}\!)$
  \end{itemize}
  we have the following commutative diagram:
  \[ \begin{tikzcd}[column sep=.5cm,row sep=.5cm]
      \C(\!\{t\}\!)^n
      \arrow[green!60!black,thick]{rrrrr}{F}
      \arrow[green!60!black,thick]{dddd}{d-A}&&&&&
      \C(\!\{t\}\!)^n\arrow[green!60!black,thick]{dddd}{d-B}
      \\ & \cM \arrow{dd}{\nabla}\arrow{rrr}{\Phi}\ular&&&
      \cM'\arrow{dd}{\nabla'}\urar
      \\
      \\ & \cM \arrow{rrr}{\Phi}\dlar&&& \cM'\drar
      \\ \C(\!\{t\}\!)^n \arrow[green!60!black,thick]{rrrrr}{F}&&&&&
      \C(\!\{t\}\!)^n
    \end{tikzcd} \]
  Thus the commutation property for the outer rectangle reads
  \[
    (d-B)\circ F=F\circ(d-A)
  \]
  which is equivalent to
  \[
    \frac{dF}{dt}=BF-FA.
  \]
\end{rem}
\begin{rem}\label{rem:distributingTransforamtionRule}
  The simple but useful rule
  \[
    {}^{(F_2F_1)}\!A =
    {}^{F_2}\!\left({}^{F_1}\!A\right)
  \]
  can be seen by calculation:
  \begin{align*}
    {}^{(F_2F_1)}\!A
    &= d(F_2F_1)(F_2F_1)^{-1}+F_2F_1 A(F_2F_1)^{-1}
  \\&=\left(
      \left(dF_2\right)̂F_1
      +F_2\left(dF_1\right)̂
    \right) F_1^{̀-1} F_2^{̀-1}
    +F_2F_1 A F_1^{̀-1}F_2^{̀-1}
  \\&= \left(dF_2\right)̂F_2^{-1}
     +F_2\left(dF_1\right)̂ F_1^{̀-1} F_2^{̀-1}
     +F_2
     \left(
       {}^{F_1}\!A-\left(dF_1\right)F^{-1}
     \right)
     F_2^{̀-1}
  \\&= \left(dF_2\right)̂F_2^{-1} +F_2 {}^{F_1}\!A F_2^{̀-1}
  \\&= {}^{F_2}\!\left({}^{F_1}\!A\right) \,.
  \end{align*}
\end{rem}
Let us now introduce equivalence relations on systems, which will correspond to
the isomorphism relations on meromorphic connections.
\begin{defn}
  We define the \emph{(formal) equivalence relation on the connection matrices
  (resp.\ on the corresponding systems)} as
  \begin{einr}
    \textbf{\boldmath$A$ is (formally) equivalent to $B$}
  \end{einr}
  if and only if
  \begin{einr}
    \textbf{\boldmath$B$ is obtained from $A$ by (formal) transformation}.
  \end{einr}
  The \emph{class of a connection matrix} is the orbit under the Gauge
  transformation by $G(\!\{t\}\!)$. The \emph{formal class} is the orbit by
  $\hat G(A)$.
  \begin{s-rem}
    Thus $A$ is (formally) equivalent to $B$ if and only if there is a (formal)
    solution of $[A,B]$.
  \end{s-rem}
\end{defn}

The defined equivalence relations are compatible with the \rewrite{isomorphisms
relations} on meromorphic connections (cf.\ \cite[Lem.5.1.3]{hotta2008}),
i.e.\ the following proposition holds.
\begin{prop}
  Two germs of meromorphic connections are (formally) isomorphic if and only if
  their corresponding connection matrices are (formally) equivalent.
\end{prop}

\begin{defn}\label{defn:isotropies}
  \marginnote{\cite[853]{Loday1994}}
  An \emph{isotropy} of $A^0$ or of $[A^0]$ is a transformation
  $\hat F$ which satisfies ${}^{\hat F}\!A^0=A^0$.
  Thus, the isotropies are the solutions of the system
  $[\End A^0]:=[A^0,A^0]$.

  Let $G_0(A^0)$ denote the set of all isotropies of $A^0$.
  \begin{s-rem}
    The isotropies are, a priori, formal transformations.
    Loday-Richaud mentions in her paper~\cite[853]{Loday1994}, that $G_0(A^0)$
    is actually a subgroup of $\Gl_n(\C[1/t,t])$.
  \end{s-rem}
  \begin{s-lem}
    \marginnote{\cite[854]{Loday1994}}
    Two formal transformations $\hat F_1$ and $\hat F_2$ take $A^0$ into
    equivalent matrices ${}^{\hat F_1}\!A^0$ and  ${}^{\hat F_2}\!A^0$ if and
    only if there exists isotropy $f_0\in G_0(A^0)$ such that
    $\hat F_1=\hat F_2f_0$ (cf.~\cite[854]{Loday1994}).
  \end{s-lem}
\end{defn}

\marginnote{\cite[Defn.5.1.6]{hotta2008}}
The meromorphic connections are distinguished into regular and irregular
meromorphic connections.
\TODO[more filltext!]
\begin{defn}
  \marginnote{\cite[86]{sabbah2007isomonodromic}}
  A connection with connection matrix $A$ has \emph{regular singularity} at $0$
  if there exists a convergent transformation, by which $A$ is obtained from a
  matrix with at most a simple pole at $t=0$.
  Otherwise, the singularity is called \emph{irregular}.
  \begin{s-rem}
    This implies that, if $A$ has irregular (resp.\ regular) singularity, then
    also all meromorphic equivalent matrices ${}^{F}\!A$ have irregular
    (resp.\ regular) singularity.
  \end{s-rem}
\end{defn}
\begin{thm}\label{thm:regularConnConstMat}
  Let $(\cM,\nabla)$ be a regular singular meromorphic connection and $A$ its
  connection matrix.
  Then there exists a formal matrix $F\in G(\!\{t\}\!)$ such that after
  transformation by $F$ the matrix $B={}^F\!A$ is constant i.e.\
  ${}^F\!A\in\Gl_n(\C)$ (cf.~\cite[Thm.II.2.8]{sabbah2007isomonodromic}
  or~\cite[Sec.5.1.2]{hotta2008}).
\end{thm}

%%%%%%%%%%%%%%%%%%%%%%%%%%%%%%%%%%%%%%%%%%%%%%%%%%%%%%%%%%%%%%%%%%%%%%%%%%%%%%%
\subsection{Fundamental solutions and the monodromy of a system}
\marginnote{\cite[Sec.4.3.2]{Loday2014},
           \\\cite[130]{hotta2008},
           \\\textbf{\cite[50]{BJL1979Birkhoff}}
           \\,\cite[26]{sibuya1990Linear}}
It is well known in the theory of linear ODEs that the solutions to a system
like (\ref{eq:ode}) form a vector space of dimension $n$ over $\C$, i.e.\ if
$x_0(t)$ and $x_{00}(t)$ are two solution of (\ref{eq:ode}) and $c_1$,
$c_2\in\C$, then is also $c_1x_0(t)+c_2x_{00}(t)$ a solution of the same system.
\begin{defn}
  \marginnote{\cite[4f]{thboalch}}
  A \emph{fundamental matrix of (formal) solutions} or \emph{(formal)
  fundamental solution} $\cY$ on a sector $\mathfrak{s}$ of the
  system $[A]$ is an invertible $n\times n$ matrix (with formal entries), which
  solves $[0,A]$ on $\mathfrak{s}$.

  This means that the columns of $\cY$ have to be $n$ $\C$-linearly independent
  solutions of the system $[A]$ on $\mathfrak{s}$.
  \begin{comment}
    \begin{s-rem}
      \begin{itemize}
      \item Some authors\PROBLEM[?] introduce multi-valued solutions, to avoid
        the restriction to sectors.
        \begin{einr}
          \emph{Multi-valued} are functions, which are not single-valued and
          \emph{single-valued} is a function $f$, which satisfies
          \[
            f(t)=f(t\exp(2\pi i)) \qquad\text{whenever both sides are defined.}
          \]
        \end{einr}
        The function $t\to t^\alpha$, for example, is single-valued whenever
        $\alpha\in\Z$.
      \end{itemize}
    \end{s-rem}
  \end{comment}
\end{defn}
The Stokes phenomenon, which will be discussed in the
Chapter~\ref{chap:stokes}, results from the fact that there is always a formal
fundamental solution, which solves a system on the full arc $S^1$.
But holomorphic solutions, which are asymptotic to the formal solution, may
exist only on small sub-sectors of $S^1$. The Stokes structures are then
roughly the information about the relation between the holomorphic solutions on the overlap
of two sectors.
\begin{rem}\label{rem:fundChangeOnFund}
  \marginnote{\cite[2.1.3]{Zein2009}}
  If the trivialization is changed by $F$ (resp.\ $\hat F$) the fundamental
  solution $\cY\in G(\!\{t\}\!)$ changes to $F\cY$ (resp.\ $\hat F\cY$).
  (cf.\ \cite[Thm.4.3.1]{Loday2014} or \cite[2.1.3]{Zein2009})
\end{rem}

\marginnote{\cite[130]{hotta2008},\\\cite[6]{heu2010}}
Choose a fundamental solution $\cY$. Then analytic continuation along a closed
path $\gamma$ in $M\backslash Z$ provides another fundamental solution.
\begin{defn}
  Let $[A]$ be a system with fundamental solution $\cY$.
  The analytic continuation of $\cY$ along a small circle around $t=0$ yields
  the fundamental solution
  \[
    \lim_{s\to2\pi}\cY(e^{\sqrt{-1}s}t)=\cY(t)L
  \]
  where $L\in\GL_n(\C)$ is called the \emph{\TODO[(formal)?]monodromy matrix}
  of $[A]$.
  \PROBLEM[wie passt das zu normal-formen?]
\end{defn}

%%%%%%%%%%%%%%%%%%%%%%%%%%%%%%%%%%%%%%%%%%%%%%%%%%%%%%%%%%%%%%%%%%%%%%%%%%%%%%%
\section{Formal classification}\label{sec:formalClassification}
\marginnote{\cite[Thm.4.3.1]{Loday2014}}
In every formal equivalence class of meromorphic connections, there are some
meromorphic connections of special form, which we will call models. They are
not unique but all of them, which are formally isomorphic to a given
meromorphic connection, lie in the same convergent equivalence class.
In fact, every element of this convergent equivalence class will be a model in
our definition.

The models will be used to classify meromorphic connections up to formal
isomorphism. Two meromorphic connections are formally isomorphic, if they are
isomorphic to the (up to convergent isomorphism) same model. In the language of
systems will normal forms be the objects, which correspond to models.

This is essentially given by the Levelt-Turittin Theorem, which states that each
meromorphic connection is, after potentially needed ramification, formally
isomorphic to such a model. Thus the Levelt-Turittin Theorem solves the formal
classification problem.

\subsection{In the language of meromorphic connections: models}
We will start by defining the models in the language of meromorphic connections.
This approach is discussed by Sabbah in~\cite{sabbah_cimpa90} and in a more
general context in~\cite[Sec.II.5]{sabbah2007isomonodromic}.
\begin{defn}\label{defn:elemnMerConnBausteine}
  For $\alpha\in\C$, define the \emph{elementary regular meromorphic
    connection of rank one} $\cN_{\alpha,0}$ as the germ
  \[
    (\cN_{\alpha,0},\nabla):=\left(\C(\!\{t\}\!),d+\frac{\alpha}{t}\right)
    \,.
  \]
  This corresponds to the system satisfied by $t^{-\alpha}$.

  \marginnote{\cite[Defn.II.2.5]{sabbah2007isomonodromic}}
  An \emph{elementary regular model of arbitrary rank} is a meromorphic
  connection which has a basis, in which the connection Matrix can be
  written as
  \[
    \frac{1}{t} (\alpha\id+N)
  \]
  where $N\in\gl_{d+1}(\C)$ is a nilpotent matrix.
  If $\alpha\id + N$ is a single Jordan Block, we denote the corresponding
  connection by $\cN_{\alpha,d}$.
\end{defn}
The Corollary~\cite[Cor.II.2.9]{sabbah2007isomonodromic} in Sabbah's book
states, that every germ of a regular meromorphic connection $(\cR,\nabla)$ is
isomorphic to some direct sum
\[
  (\cR,\nabla)\cong\bigoplus_{\alpha,d}(\cN_{\alpha,d},\nabla)\,.
\]
of elementary regular meromorphic connections.
The rough idea to proof this is to use Theorem~\ref{thm:regularConnConstMat}
and assume that the connection matrix $A$ of $(\cR,\nabla)$ is constant.
Since the Gauge transformation by a constant matrix $F\in\Gl(\C)$ is given by
${}^F\!A=FAF^{-1}$, i.e.\ a similarity transformation, one can find a
transformation $F$ such that ${}^F\!A$ is in Jordan normal form.
One can then use Corollary~\ref{cor:MatOfSumOfMerCon} to decompose the
corresponding to the structure of the Jordan normal form and obtain the result.
This solves the classification problem, for regular meromorphic connections.
For a detailed analysis of regular meromorphic connections see Sabbah's
book~\cite[Sec.II.2]{sabbah2007isomonodromic} or the
book~\cite[Sec.5.2]{hotta2008} from Hotta et al.

The simplest irregular meromorphic connections will be called elementary and
are defined below.
\begin{defn}
  For a $\phi\in\C(\!\{t\}\!)$ we use $\cE^{\phi}$ to denote the meromorphic
  connection
  \[
    (\cE^{\phi},\nabla):=(\C(\!\{t\}\!),d-\phi')\,.
  \]
  This corresponds to the system satisfied by the function $e^\phi$, since
  $(d-\phi')e^\phi=0$.
  \begin{s-prop}\label{cor:uniqueAmbassadorForPhi}
    $\cE^\phi$ is determined by the class of $\phi$ in
    $\C(\!\{t\}\!)/\C\{t\}\cong t^{-1}\C[t^{-1}]$. In the following, we will
    only assume, that $\phi$ is always given by the unique ambassador which has
    no holomorphic part.
  \end{s-prop}
  \marginnote{\cite[Def.II.5.2]{sabbah2007isomonodromic}}
  A germ $(\cM,\nabla)$ is called \emph{elementary} if it is isomorphic to some
  germ $(\cE^\phi,\nabla)\otimes(\cR,\nabla)$ where $(\cR,\nabla)$ has regular
  singularity at $\{0\}$, i.e.\ is isomorphic to a direct sum of regular
  elementary meromorphic connections.
  \marginnote{\cite[II.2.f]{sabbah2007isomonodromic}}
\end{defn}
\begin{defn}\label{defn:model}
  \def\myPhi{\textcolor{red!60!black}{\phi}}
  \def\myE{\textcolor{green!40!black}{\cE^{\myPhi}}}
  A germ $(\cM',\nabla')$ is a \emph{model} if there exists after ramification
  $\cM=\pi^*\cM'$ by $\pi$ a isomorphism
  \[
    \lambda:(\cM,\nabla) \overset{\cong}{\longrightarrow}
    \bigoplus_{\phi} \cE^\phi \otimes \cR_{\phi}
  \]
  to a direct sum of elementary meromorphic connections where the
  $\phi\in t^{-1}\C[t^{-1}]$ are pairwise distinct and the $\cR_{\phi}$ have
  regular singularity.
  \begin{s-rem}
    Here it is not necessary to understand ramification, since we restrict to
    the unramified case.
  \end{s-rem}
\end{defn}
The important theorem here is the Levelt-Turittin Theorem, which solves the
formal classification problem.
\begin{thm}[Levelt-Turittin decomposition]\label{thm:leveltTurittin}
  To each germ $(\cM',\nabla')$ of a meromorphic connection there exists, after
  potentially needed pullback $\pi^{*}\textcolor{black}{\cM'}=:\cM$ by some
  suitable ramification $t=z^q$ of order $q\geq1$, a \textbf{formal} isomorphism
  \[
    \hat{\lambda}:\hat{\cM}\overset{\cong}\longrightarrow
    \hat{\textcolor{black}{\cM}}^{nf}:=\hat\cO_M\otimes\cM^{nf}
  \]
  to a model $\cM^{nf}$.
  We then call $\cM^{nf}$ a \emph{formal decomposition} or \emph{formal model}
  of $\cM$ or $\cM'$.
  \begin{s-rem}\label{rem:leveltTurittin}
    But there is in general \textbf{no} lift of the isomorphism $\hat\lambda$,
    i.e.\ there is no isomorphism making the diagram
    \[ \begin{tikzcd}
        \cM \dar \arrow[dotted]{r}[description]{?} & \cM^{nf} \dar
        \\\hat\cM \rar{\hat\lambda} & \hat\cM^{nf}
    \end{tikzcd} \]
    commutative.
    But sectorwise, there are lifts given by the main asymptotic existence
    theorem (cf.\ Theorem~\ref{thm:maet}).
  \end{s-rem}
\end{thm}
Proofs of this theorem can be found in multiple places, for example in
\cite[Thm.5.4.7]{sabbah_cimpa90} or in~\cite{myBa}.

%%%%%%%%%%%%%%%%%%%%%%%%%%%%%%%%%%%%%%%%%%%%%%%%%%%%%%%%%%%%%%%%%%%%%%%%%%%%%%%
\subsection{In the language of systems: normal forms}\label{sec:normalForms}
\marginnote{\cite{Remy2014}
           \\,\cite[146]{sibuya1990Linear}}
Here we will see that the fundamental solution of a model can be written in a
special form and that this property characterizes models. This will be used to
define normal forms which will be models on the level of systems.
\begin{defn}
  The fundamental solution corresponding to a model (resp.\ a normal form, see
  Definition~\ref{defn:normSol}) will be called a
  \emph{normal solution}.
\end{defn}

Let $(\cM,\nabla)$ be an unramified meromorphic connection which is via some
$\hat\lambda$ formally isomorphic to some model
$\cM^{nf}=\bigoplus_{j=1}^s\cE^{\phi_j}\otimes\cR_j$, where we assume that
$\phi_j\in t^{-1}\C[t^{-1}]$ (cf.~Corollary~\ref{cor:uniqueAmbassadorForPhi}).
For every $j$ is a connection matrix, corresponding to $\cR_j$ given by
$\frac{1}{t}L_j$, where $L$ is in Jordan normal form.
A connection matrix to $\cE^{\phi_j}\otimes\cR_j$ is then by
$q_j'(t^{-1})\cdot\id_{n_j}+\frac{1}{t}L_j$ given\TODO[(cf.~???)], where
$q_j(t^{-1})=\phi_j(t)$\PROBLEM[repl. all $\phi$ by $q$?] is a polynomial in
$t^{-1}$ without constant term and
$q_j'(t^{-1})=\frac{d}{dt}\phi_j(t)$\footnote{By abuse of notation we denote by
  $q_j'(t^{-1})$ the derivation $\frac{d}{dt}(q_j(t^{-1}))=
  \left(\left(\frac{d}{dt}q_j\right)\circ t^{-1}\right) \cdot\frac{d}{dt}t^{-1}$.
  The same does apply for $Q'(t^{-1})$.}.
A connection matrix for $\cM^{nf}$ is then obtained by
\begin{align*}
  A^0&= \bigoplus_{j=1}^s
       \left(q_j'(t^{-1})\cdot\id_{n_j}+\frac{1}{t}L_j\right)
  \\ &=\bigoplus_{j=1}^sq_j'(t^{-1})\cdot\id_{n_j}
       +\frac{1}{t}\bigoplus_{j=1}^sL_j
  \\ &=Q'(t^{-1})+\frac{1}{t}L \,,
\end{align*}
where
$Q(t^{-1}):=\bigoplus_{j=1}^sq_j(t^{-1})\cdot\id_{n_j}$
and $L:=\bigoplus_{j=1}^sL_j$.
\begin{defn}\label{defn:structureComparison}
  Let $L$ be a block diagonal matrix $L=\bigoplus_{j=1}^sL_j$, where the $L_j$
  are of size $n_j\times n_j$ and let $Q$ be a diagonal matrix.
  We will say that \emph{the block structure of $L$ is finer then the
  structure of $Q$} if there are $q_j$'s such that
  $Q=\bigoplus_{j=1}q_j\cdot\id_{n_j}$.
  \begin{s-rem}
    \begin{enumerate}\label{rem:structureComparison}
      \item The matrices $Q$ and $L$ defined above clearly satisfy this
        condition.
      \item When this condition is satisfied do $L$ and $Q$ commute.
    \end{enumerate}
  \end{s-rem}
\end{defn}
\begin{prop}\label{prop:fundSolBuilder}
  Let $L\in \GL_n(\C)$ be constant and
  $Q(t^{-1})=\diag(q_1(t^{-1}),\dots,q_n(t^{-1}))$ a diagonal matrix of
  polynomials in $t^{-1}$ where $L$ is in Jordan normal form and its block
  structure is finer then the structure of the derivative $Q(t^{-1})$.

  The matrix $\cY_0:=t^L e^{Q(t^{-1})}$ is then a fundamental solution of the
  system determined by the matrix
  \[
    A^0=Q'(t^{-1})+L\frac{1}{t} \,.
  \]
\end{prop}
\begin{proof}
  It is easy to see (or cf.~\cite[Appendix C.1]{Balser2000Formal}) that the
  block structure of $t^L$ is finer then the structure of $Q'(t^{-1})$ and thus
  $t^LQ'(t^{-1})t^{-L}=Q'(t^{-1})$, since the block structure is preserved under
  the exponential and derivation.
  Thus we can prove that $\cY_0$ is a matrix consisting of solutions:
  \begin{align*}
    \frac{d}{dt}\cY_0
    &=\frac{d}{dt}\left(t^Le^{Q(t^{-1})}\right)
  \\&= t^L\frac{d}{dt}e^{Q(t^{-1})} + \frac{d}{dt}t^Le^{Q(t^{-1})}
  \\&= t^LQ'(t^{-1})e^{Q(t^{-1})} + L\frac{1}{t}t^Le^{Q(t^{-1})}
  % \\&=L\frac{1}{t}t^Le^{Q(t^{-1})}+Q'(t^{-1})t^Le^{Q(t^{-1})}
  \\&=\left(t^LQ'(t^{-1})t^{-L}+L\frac{1}{t}\right)t^Le^{Q(t^{-1})}
  \\&=\left(Q'(t^{-1})+L\frac{1}{t}\right)t^Le^{Q(t^{-1})}
  \\&=A^0\cY_0 \,.
  \end{align*}
  The invertability condition is clear.
\end{proof}
\begin{rem}
  If one starts without the assumption on the structure, one only obtains
  \[
    A=t^LQ'(t^{-1})t^{-L}+L\frac{1}{t}
  \]
  as possible system matrix. But the obtained matrix may not define a system,
  since it could contain entries, which are not in $\C(\!\{t\}\!)$.

  A simple example where this happens is given by
  \begin{itemize}
  \item $L=\begin{pmatrix} \alpha & 1 \\ 0 & \alpha \end{pmatrix}$ and
  \item $Q=\begin{pmatrix} \phi_1 & 0 \\ 0 & \phi_2 \end{pmatrix}$,
  \end{itemize}
  with $\phi_1\neq\phi_2$.
  The matrix $A$ is then
  \begin{align*}
    A  &=t^LQ'(t^{-1})t^{-L}+L\frac{1}{t}
    \\ &=
    \begin{pmatrix}
      \phi_1' & (-\phi_1'+\phi_2')\ln(t)
      \\0     & \phi_2'
    \end{pmatrix}
    +
    \begin{pmatrix}
        \alpha & 1
        \\ 0   & \alpha
    \end{pmatrix} \frac{1}{t}
    \notin\gl_n(\C(\!\{t\}\!))
    \,.
  \end{align*}
\end{rem}
To $\hat\lambda$ and the (implicitly) used choices of bases corresponds a
formal transformation $\hat F$ and thus we have
\begin{itemize}
\item a fundamental solution $\hat Ft^L e^{Q(t^{-1})}$ (cf.\
  Remark~\ref{rem:fundChangeOnFund}) and
\item a connection matrix ${}^{\hat F}\!A^0$
\end{itemize}
of the meromorphic connection $(\cM,\nabla)$, which was defined at the beginning
of this subsection.
\begin{cor}
  From the Levelt-Turittin theorem we deduce that every meromorphic connection
  $(\cM,\nabla)$ and thus every system $[A]$ has a fundamental solution of the
  form
  \[
    \mathcal{Y}=\hat F t^L e^{Q(t^{-1})}
  \]
  where $\hat F$ is a formal transformation, solving the differential equation
  corresponding to the formal isomorphism
  $(\cM^{nf},\nabla^{nf})\to(\cM,\nabla)$.
  On the other hand is a system, which has $\cY$ as fundamental solution, given
  by the connection matrix
  \[
    \bar A={}^{\hat F}\!\left(Q'(t^{-1})+L\frac{1}{t}\right) \,.
  \]
  \begin{comment}
    \begin{s-rem}
      It is always possible to permutate the columns of a fundamental solution
      by
      \[
        P^{-1}\mathcal{Y}P=\hat F t^{P^{-1}LP} e^{P^{-1}Q(t^{-1})P}
      \]
      with a permutation matrix $P$ and \rewrite{obtain another fundamental
      solution for the same system} (cf.\ \cite[73]{Loday2014}).
    \end{s-rem}
  \end{comment}
\end{cor}
In the special case when $\hat\lambda$ is a convergent isomorphism, i.e.\ in
the case, when $A$ is already a model, we see that it has a fundamental
solution of the form $Ft^L e^{Q(t^{-1})}$, where $F$ is a convergent
transformation corresponding to $\hat\lambda$.
In fact are the models uniquely characterized as the meromorphic connections,
with a fundamental solution which can be written as $\cY_0=Ft^Le^{Q(t^{-1})}$
and we will use this fact to say when a system is something corresponding to a
model.
\begin{defn}\label{defn:normSol}
  Let $[A]$ be a system.
  We call $[A]$ (or $A$) a \emph{normal form} if a fundamental solution of
  $[A]$ can be written as
  \[
    \mathcal{Y}_0(t)=F t^L e^{Q(t^{-1})}
  \]
  with
  \begin{itemize}
    \item an \emph{irregular part} $e^{Q(t^{-1})}$ of $\mathcal{Y}_0$
      determined by
      \[
        Q(t^{-1})=\underset{j=1}{\overset{s}{\bigoplus}}q_j(t^{-1})\id_{n_j}
          =\diag(\underset{n_1\text{-times}}{\underbrace{%
          q_1,\dots,q_1}},q_2,\dots,q_s)
      \]
      where the $q_i(t^{-1})$ are polynomials in $\frac{1}{t}$ (or in a
      fractional power $\frac{1}{s}=\frac{1}{t^{1/p}}$ of $\frac{1}{t}$ for the
      ramified case) such that $q_j(0)=0$, i.e.\ without constant term,
    \item a constant matrix $L\in\gl_n(\C)$ called the \emph{matrix of formal
      monodromy}, where $t^L$ means $e^{L\ln t}$ and
    \item a (convergent) transformation $F\in G(\!\{t\}\!)$.
      \begin{s-rem}
        By changing the basis via $F^{-1}$ we obtain from $[A^0]$ the
        equivalent system $[{}^{F^{-1}}\!A^0]$ with the normal solution
        $F^{-1}\mathcal{Y}_0(t)=t^L e^{Q(t^{-1})}$.
        Thus it is always possible to assume that the transformation matrix $F$
        is trivial.
      \end{s-rem}
  \end{itemize}
  The normal forms will often be denoted $A^0$.
  If $A$ is formally equivalent to a normal form $A^0$ we say that $A^0$
  \emph{is a normal form for} $A$ and for $[A]$.
\end{defn}
\begin{cor}\label{cor:structuralAssumptions}
  Since we only look at unramified connections, we are able to assume that $L$
  is in Jordan normal form and that it has a block structure, which is finer
  then the structure of $Q=\bigoplus_{j=1}^sq_j(t^{-1})\cdot\id_{n_j}$
  (cf.~\cite[Sec.1]{Remy2014} or~\cite[Sec.4]{Martinet1991}).
\end{cor}
The following corollary enables us to use normal forms in place of models.
\begin{cor}
  A meromorphic connection $(\cM,\nabla)$ is a model if and only if its
  connection matrix $A$ is a normal form.
\end{cor}

% \begin{comment}
% %%%%%%%%%%%%%%%%%%%%%%%%%%%%%%%%%%%%%%%%%%%%%%%%%%%%%%%%%%%%%%%%%%%%%%%%%%%%%%%
% \subsection{In the language of systems: normal forms}
% In the language of system, the equivalent of models are normal forms, which are
% characterizes by the structure of their fundamental solutions.

% The normal forms are then the systems $[A^0]$, with a fundamental solution
% $\cY_0$ in a special form and from the Levelt-Turittin we then can deduce that
% every system $[{}^{\hat F}\!A^0]$, with $\hat F\in G(\!(t)\!)$, has a
% fundamental solution in the form $\hat F\cY_0$.
% \begin{defn}\label{defn:normSol}
%   Let $[A]$ be a system.
%   We call $[A]$ (or $A$) a \emph{normal form} if its fundamental solution can
%   be written as
%   \[
%     \mathcal{Y}_0(t)=F t^L e^{Q(t^{-1})}
%   \]
%   with
%   \begin{itemize}
%     \item \emph{irregular part} $e^{Q(t^{-1})}$ of $\mathcal{Y}_0$ determined
%       by
%       \[
%         Q(t^{-1})=\underset{j=1}{\overset{s}{\bigoplus}}q_j(t^{-1})\id_{n_j}
%           =\diag(\underset{n_1\text{-times}}{\underbrace{%
%           q_1,\dots,q_1}},q_2,\dots,q_s)
%       \]
%       where the $q_i(t^{-1})$ are polynomials in $\frac{1}{t}$ (or in a
%       fractional power $\frac{1}{s}=\frac{1}{t^{1/p}}$ of $\frac{1}{t}$ for the
%       ramified case) such that $q_j(0)=0$, i.e.\ without constant term,
%     \item $L\in\gl_n(\C)$ constant
%       \PROBLEM[block diagonal (corresponding to $Q$)]
%       matrix called the \emph{matrix of formal
%       monodromy}, where $t^L$ means $e^{L\ln t}$ and
%       \marginnote{in \cite[1]{Remy2014} $L$ is just a Jordan normal form.  Is
%       this generic enough?}
%     \item a transformation $F\in G(\!\{t\}\!)$.
%   \end{itemize}
%   The normal forms will often be denoted $A^0$.
%   If $A$ is formally equivalent to a normal form $A^0$ we say that $A^0$
%   \emph{is a normal form for} $A$ and for $[A]$.

%   The fundamental solution $F t^L e^{Q(t^{-1})}$ of a normal form will be
%   called \emph{normal solution} and will often be denoted by $\cY_0$.
%   \begin{s-rem}
%     \begin{enumerate}
%       \item \TODO[Cor from base change for fundamental solutions] By changing the
%         basis via $F^{-1}$ wo obtain from $[A^0]$ the equivalent system
%         $[{}^{F^{-1}}\!A^0]$ with the normal solution
%       \[
%         F^{-1}\mathcal{Y}_0(t)=t^L e^{Q(t^{-1})} \,.
%       \]
%       Thus it is always possible to assume that the transformation matrix $F$ is
%       trivial.
%       \item In the unramified case $Q$ and $L$ commute thus $L$ can be supposed
%         in Jordan form (cf.\ \cite[Sec.4]{Martinet1991}).
%     \end{enumerate}
%   \end{s-rem}
% \end{defn}
% \begin{thm}\label{thm:modelEqNormalForm}
%   A meromorphic connection $(\cM,\nabla)$ is a model if and only if its
%   connection matrix $A$ defines a normal form $[A]$.
%   \begin{s-rem}
%     \begin{itemize}
%       \item This allows us, to say that a model $(\cM^{nf},\nabla^{nf})$ is a
%         normal form and \rewrite{vice versa}.
%         \rewrite{This} explains, why we mark models with ${}^{nf}$.
%       \item In the unramified case are \rewrite{the $q_i(t^{-1})$ the
%         $\phi_i(t)$} from definition~\ref{defn:model}.
%         \TODO[Why/realy?]
%     \end{itemize}
%   \end{s-rem}
% \end{thm}
% \begin{proof}
%   \textbf{``\Rightarrow{}'':}
%   Let $\cM=\bigoplus_{\phi}\cE^\phi\otimes\cR_\phi$ be a model and let us
%   first fix a $\phi$ and look at the connection $\cE^\phi\otimes\cR_\phi$.
%   We know that $\cR_\phi$ is a direct sum of regular elementary meromorphic
%   connections whose fundamental matrices are of the form
%   $\frac{1}{t}(\alpha\id+N)$ where $\alpha\in\C$ and $N$ is nilpotent
%   (cf.\ Definition~\ref{defn:elemnMerConnBausteine}).
%   Thus the connection matrix for $\cR_\phi$ writes as
%   \[
%     \frac{1}{t} \bigoplus_{(\alpha,N)}\left( \alpha\id+N\right)
%     \,.
%   \]
%   A connection matrix of $\cE^\phi\otimes\cR_\phi$ is then given\PROBLEM[!]
%   by
%   \[
%     t^{\bigoplus_{(\alpha,N)}\left( \alpha\id+N\right)}
%     \left(-\phi'\id\right)
%     t^{-\bigoplus_{(\alpha,N)}\left( \alpha\id+N\right)}
%     + \frac{1}{t} \bigoplus_{(\alpha,N)}\left( \alpha\id+N\right) \,.
%   \]
%   Since $\cM$ is a sum of such elementary connections, The connection matrix
%   of $\cM$ writes as
%   \[
%     \bigoplus_\phi\left(
%       - t^{\bigoplus_{(\alpha,N)}\left( \alpha\id+N\right)}
%       \left(\phi'\id\right)
%       t^{-\bigoplus_{(\alpha,N)}\left( \alpha\id+N\right)}
%       +
%       \frac{1}{t} \bigoplus_{(\alpha,N)}\left( \alpha\id+N\right)
%     \right) \,,
%   \]
%   which can be rewritten as
%   \[
%     \underset{t^LQ'(t^{-1})t^{-L}}{
%       \underset{\text{\rotatebox[origin=c]{-90}{$=$}}}{%
%         \underbrace{-\bigoplus_\phi
%           \left(
%             t^{\bigoplus_{(\alpha,N)}\left( \alpha\id+N\right)}
%             \left(\phi'\id\right)
%             t^{-\bigoplus_{(\alpha,N)}\left( \alpha\id+N\right)}
%           \right)
%         }
%       }
%     }
%     +
%     \frac{1}{t}
%     \underset{L}{
%       \underset{\text{\rotatebox[origin=c]{-90}{$=$}}}{%
%         \underbrace{%
%           \bigoplus_\phi\left(
%             \bigoplus_{(\alpha,N)}\left( \alpha\id+N\right)
%           \right)
%         }
%       }
%     }
%     \,.
%   \]
%   Thus it has
%   \[
%     t^{\bigoplus_\phi\left(
%         \bigoplus_{(\alpha,N)}\left( \alpha\id+N\right)
%     \right)}
%     e^{-\bigoplus_\phi\phi'\id}
%   \]
%   as fundamental solution (cf.\ proposition~\ref{prop:fundSolBuilder}).
%   \PROBLEM[Construction only yields unipotent $N$]

%   \PROBLEM{}

%   \textbf{``\Leftarrow{}'':}
%   Let $A^0$ be a normal form with $\mathcal{Y}_0(t)=F t^L e^{Q(t^{-1})}$ as a
%   normal solution. We use Proposition~\ref{prop:fundSolBuilder} to assume that
%   the connection matrix $A^0$ has the form
%   \[
%     A^0=t^LQ'(t^{-1})t^{-L}+L\frac{1}{t}
%   \]
%   and the Proposition~\ref{prop:systToMeromConn} yields our meromorphic
%   connection $(\cM_{A^0},\nabla_{A^0})=(\C(\!\{t\}\!)^n,d-A^0)$ for
%   \rewrite{which we want to prove, to be} a model.
%   \TODO[We can use Proposition~\ref{prop:MatOfSumOfMerCon} to assume, that $L$ is
%     \textbf{only one Jordan block}.]

%   \PROBLEM[Jordan-NF]
% \end{proof}
% \begin{cor}
%   \TODO[Start with this, bevore defn of normal forms?]
%   From the Levelt-Turittin theorem\comm{, the fundamental solution
%   transformation rules} and Theorem~\ref{thm:modelEqNormalForm} we
%   deduce that every system $[A]$ with normal form $[A^0]$ and normal solution
%   $t^L e^{Q(t^{-1})}$, has the fundamental solution
%   \[
%     \mathcal{Y}=\hat F t^L e^{Q(t^{-1})}
%   \]
%   where $\hat F$ is a solution of $[A^0,A]$.
%   \begin{s-rem}
%     It is always possible to permutate the columns of a fundamental solution by
%     \[
%       P^{-1}\mathcal{Y}P=\hat F t^{P^{-1}LP} e^{P^{-1}Q(t^{-1})P}
%     \]
%     with a permutation matrix $P$ and \rewrite{obtain another fundamental
%     solution for the same system} (cf.\ \cite[73]{Loday2014}).
%   \end{s-rem}
% \end{cor}
% \end{comment}

%%%%%%%%%%%%%%%%%%%%%%%%%%%%%%%%%%%%%%%%%%%%%%%%%%%%%%%%%%%%%%%%%%%%%%%%%%%%%%%
\section{The main asymptotic existence theorem (M.A.E.T)}\label{sec:MAET}
\begin{comment}
  \begin{multicols}{2}
    \textbf{Classical:}
    \begin{itemize}
      \item \cite[Thm.4.4.1]{Loday2014}
      \item \cite[Thm.7.10]{van2003galois}{\tiny\cite[Thm.7.12]{van2003galois}}
      \item \cite[Thm.IV.12.1]{wasow2002asymptotic}
      \item \cite[5.3.Thm.1]{Varadarajan96linearmeromorphic}
      \item \cite[207]{Balser2000Formal}: Some historical remarks
      \item \cite[Thm.A]{BJL1979Birkhoff}
    \end{itemize}
  \columnbreak
    \textbf{Sheafical:}
    \begin{itemize}
      \item \cite[Thm.2.3.1]{sabbah_cimpa90}
      \item \cite[Sec.4.4]{Loday2014}
    \end{itemize}
  \end{multicols}
\end{comment}
Here we want to state the main asymptotic expansion theorem (or often M.A.E.T.)
which is essentially a deduction from the Borel-Ritt Lemma.
\marginnote{\cite[207]{Balser2000Formal}}
It states that to every formal solution of a system of meromorphic differential
equations and every sector with sufficiently small opening, one can find a
holomorphic solution of the system having the formal one as its asymptotic
expansion.

\begin{defn}\label{defn:lift}
  \marginnote{\cite[855]{Loday1994}}
  Let $A$ be via $\hat F$ formally equivalent to $A^0$.
  We call $F$ a \emph{lift of $\hat F$ on $I\subset S^1$} if
  \begin{itemize}
    \item $F\sim_I\hat F$
      (cf.\ Page~\pageref{page:notationForAsymptoticExpansion}) and
    \item $F$ satisfies the same system $[A^0,A]$ as $\hat F$.
  \end{itemize}
\end{defn}
The following theorem, often called the main asymptotic existence theorem, can
be for example found in as Theorem A in the paper~\cite{BJL1979Birkhoff} from
Balser, Jurkat and Lutz, Theorem 3.1 in Boalch's paper~\cite{boalch}
or Theorem 4.4.1 in Loday-Richaud's Book~\cite{Loday2014}.
\begin{thm}[M.A.E.T]\label{thm:maet}
  To every $\hat F\in G(\!(t)\!)$ and to every small enough arc
  $I\subsetneq S^1$ there exists a lift $F$ on $I$.
  \begin{s-rem}
    \marginnote{\cite[Thm.4.4.1]{Loday2014}}
    If we write the system $[A^0,A]$ as a differential operator $\triangle$.
    The theorem then states that $\triangle$ acts surjectively on the sheaf
    $\cA^{<0}$, i.e.\ the sequence
    \[
      \cA^{<0}\overset{\triangle}\longrightarrow\cA^{<0} \longrightarrow 0
    \]
    is an exact sequences of sheaves of $\C$-vector spaces.
    \begin{comment}
      (cf.~\cite[App.1;Thm.1]{malgrange1991})
    \end{comment}
  \end{s-rem}
\end{thm}
\begin{rem}
  \PROBLEM[Think about it?]
  In the language of meromorphic connections is this theorem sometimes called
  \emph{sectorial decomposition} and stated for example
  in~\cite[Thm.II.5.12]{sabbah2007isomonodromic}
  and~\cite[Sec.II.2.4]{sabbah_cimpa90}:
  \begin{s-thm}[Sectorial decomposition]\label{thm:sectorialDecompFromMAET}
    \PROBLEM[needs more defns]
    Let $(\cM,\nabla)$ be a meromorphic connection and let
    $\hat\lambda:\hat\cM\to\hat\cM^{nf}$ be the isomorphism given by
    Theorem~\ref{thm:leveltTurittin} together with the model $\cM^{nf}$.
    There exists then, for any $e^{i\theta^0}\in S^1$, an isomorphism
    $\tilde\lambda_{\theta^0}:
    \tilde\cM_{\theta^0}=\cA_{\theta^0}\otimes\cM\to\tilde\cM_{\theta^0}^{nf}$
    lift $\hat\lambda$, \rewrite{that is, such that} the following diagram
    \[ \begin{tikzcd}
        \tilde\cM_{\theta^o} \dar \rar{\tilde\lambda_{\theta^o}}
        & \tilde\cM^{nf}_{\theta^o} \dar
        \\\hat\cM \rar{\hat\lambda}
        & \hat\cM^{nf}
    \end{tikzcd} \]
    commutes
  \end{s-thm}
\end{rem}

%%%%%%%%%%%%%%%%%%%%%%%%%%%%%%%%%%%%%%%%%%%%%%%%%%%%%%%%%%%%%%%%%%%%%%%%%%%%%%%
\section{The classifying set}\label{sec:classifyingSet}
We want to understand the Set
$\bigl\{\big[(\cM,\nabla)\big]\bigr\}$ of the (convergent)
isomorphism classes of all meromorphic connections. Since we can use the formal
classification (cf.\ Section~\ref{sec:formalClassification}) and we know that
the convergent classification is \rewrite{finer} than the formal classification,
we can reduce the problem by fixing a model $(\cM^{nf},\nabla^{nf})$ with the
corresponding normal form $A^0$.
\rewrite{Thus we can} restrict to the subset
\begin{multline*}
  {}^0C(\cM^{nf},\nabla^{nf})=\bigl\{
    \bigl[(\cM,\nabla)\bigr]
    \mid \text{there exists a formal isomorphism }
  \\\qquad\hat f:(\hat\cM,\hat\nabla)
      \overset{\sim}\longrightarrow
      (\hat\cM^{nf},\hat\nabla^{nf})
  \bigr\}
\end{multline*}
of all isomorphism classes of meromorphic connections, which are formally
isomorphic to $(\cM^{nf},\nabla^{nf})$. This is the set that we will be
calling the \emph{classifying set (to $(\cM^{nf},\nabla^{nf})$)} and we will
also denote it also by ${}^0C(A^0)$, if we are using the language of systems.
\begin{rem}
  Note that we classify
  \begin{einr}
    meromorphic connections within fixed \textbf{formal meromorphic classes,
    modulo meromorphic equivalence}.
  \end{einr}
  Whereas for example Boalch in~\cite{boalch,thboalch} classifies
  \begin{einr}
    meromorphic connections within fixed \textbf{formal analytic classes,
    modulo analytic equivalence}
  \end{einr}
  as it was done in the older literature.
  This makes no difference, since the resulting classifying sets are isomorphic
  (cf.~\cite{thboalch} or~\cite{babbitt1989local}).

  This distinction relates to the difference between \textbf{‘regular
  singular’} connections and \textbf{‘logarithmic’} connections
  (cf.~\cite[Rem.1.41]{thboalch}).
\end{rem}
\marginnote{\cite[6]{thboalch}
  \\\tiny{(\cite[19]{boalch})}
  \\\cite[852]{Loday1994}
  \\\cite[111]{sabbah2007isomonodromic}
  \\\cite{babbitt1983}}
It is convenient to look at the slightly larger space of isomorphism classes of
\emph{marked (meromorphic) pairs}
\[
  \cH(\cM^{nf},\nabla^{nf})=\bigl\{
    \bigl[(\cM,\nabla,\hat f)\bigr]
      \mid
      \hat f:(\hat\cM,\hat\nabla)
        \overset{\sim}\longrightarrow
        (\hat\cM^{nf},\hat\nabla^{nf})
  \bigr\}
\]
in which we also handle the additional information of the formal isomorphism by
which the meromorphic connection is isomorphic to the model.
The isomorphisms of marked pairs are defined as expected:
\begin{defn}\label{defn:isomsOfPairs}
  Two germs $(\cM,\nabla,\hat f)$ and $(\cM',\nabla',\hat f')$ are
  isomorphic if there exists an isomorphism
  $g:(\cM,\nabla)\overset{\sim}\longrightarrow(\cM',\nabla')$ such that
  $\hat f=\hat f'\circ \hat g$.
  \begin{s-rem}
    \rewrite{Sabbah states in} \cite[111]{sabbah2007isomonodromic} that such
    an isomorphism is then unique.
  \end{s-rem}
\end{defn}

Equivalently, one can talk in terms of systems. We then denote by
\[
  \Syst_m(A^0):=\bigl\{[A]
    \mid A={}^{\hat F}\!A^0 \text{ for some } \hat F\in G(\!(t)\!)\bigr\}
\]
the set of systems formally meromorphic equivalent to $A^0$.
Since we use meromorphic equivalences, in contrast to \cite{boalch,thboalch},
we denote that in $\Syst_m$ by the subscript ${}_m$.
Thus ${}^0C(\cM^{nf},\nabla^{nf})$ corresponds to
the set ${}^0C(A^0):=\Syst_m(A^0)/G(\!\{t\}\!)$ of meromorphic classes which
are formally equivalent to $A^0$.
Analogous, $\cH(\cM^{nf},\nabla^{nf})$ corresponds to the set $\cH(A^0)$ of
equivalence classes, i.e.\ orbits of $G(\!\{t\}\!)$, in
\label{page:ofDefnOfIsomOfMarkedPairs}
\[
  \hat\Syst_m(A^0):=\bigl\{\bigl(A,\hat F\bigl)
    \mid A={}^{\hat F}\!A^0 \text{ for some } \hat F\in G(\!(t)\!)\bigl\} \,.
\]
The isomorphisms defined in Definition~\ref{defn:isomsOfPairs} translate into
the following equivalence relation:
\begin{einr}
  two marked pairs $(A,\hat F)$ and $(A',\hat F')$ are equivalent, if and only
  if there is a base change $H$ such that
  \PROBLEM[$H\in G(\!\{t\}\!)$? or asymptotic functions?]
  \begin{enumerate}
  \item $A'={}^{H}\!A$, i.e.\ $H$ is a solution of $[A,A']$, and
  \item $\hat F'=\hat H\hat F$ (cf.\ \cite[71]{babbitt1989local}).
  \end{enumerate}
\end{einr}
In the following diagram the first condition is equivalent to the commutativity
of the \textcolor{green!60!black}{green} square. The second property corresponds
to the commutation property of the top (resp.\ bottom) triangle on the level of
asymptotic expansions.
\[ \begin{tikzcd}[column sep=1cm,row sep=.5cm]
  &&&\C(\!\{t\}\!)^n \arrow[green!60!black, thick]{ddd}{d-A}
    \arrow[green!60!black, thick]{ddr}{H}
\\\C(\!\{t\}\!)^n \arrow{rrru}{\hat F}
  \arrow[ultra thick, white]{rrrrd}{\hat F'}
  \arrow{rrrrd}{\hat F'}
  \arrow{ddd}{d-A^0}
\\&&&&\C(\!\{t\}\!)^n \arrow[green!60!black, thick]{ddd}{d-A'}
\\&&&\C(\!\{t\}\!)^n \arrow[green!60!black, thick]{ddr}{H}
\\\C(\!\{t\}\!)^n \arrow{rrru}{\hat F} \arrow{rrrrd}{\hat F'}
\\&&&&\C(\!\{t\}\!)^n
\end{tikzcd} \]

\begin{lem}
  \marginnote{\url{http://mathworld.wolfram.com/Stabilizer.html}}
  Since $G_0(A^0)$ is by definition the stabilizer of $A^0$ (cf.\
  definition~\ref{defn:isotropies}) and $\Syst_m(A^0)$ is the corresponding
  orbit we can use the Proposition 3.1 of the
  book~\cite{wielandt1964finite} from Wielandt and deduce
  \[
    \Syst_m(A^0)\cong \hat G(A^0)/G_0(A^0) \,.
  \]
  \begin{s-cor}\label{cor:isomorphyOfClassfset}
    \marginnote{\cite[Eq.1.9b]{babbitt1983}}
    Thus the \emph{set of meromorphic classes of systems formally equivalent
    to $A^0$} are just the orbits of $G(\!\{t\}\!)$ in $\Syst_m(A^0)$ that is
    \[
      {}^0C(A^0)\cong G(\!\{t\}\!)\backslash\hat G(A^0)/G_0(A^0)
    \]
    whereas the \emph{set of meromorphic classes of marked pairs $\cH(A^0)$
    of $[A^0]$} is canonically isomorphic to the left quotient
    $G(\!\{t\}\!)\backslash\hat G(A^0)$ (cf.\ \cite[Lem.1.17]{thboalch}).
  \end{s-cor}
\end{lem}
The group $G_0(A^0)$ is easy to compute and is often trivial. In fact are the
elements of $G_0(A^0)$ block-diagonal corresponding to the structure of $Q$,
see~\cite[77]{Loday2014}.
Thus the structure of ${}^0C(A^0)$ is easily deduced from the structure of
$G(\!\{t\}\!)\backslash\hat G(A^0)$.

%%% Local Variables:
%%% TeX-master: "Maximilian_Huber-Masters_Thesis-with_notes.tex"
%%% End:

\fi

\ifnum\myHideVariable=0
\chapter{Stokes Structures}\label{chap:stokes}
\begin{comment}
  \begin{itemize}
  \item
  Characterized completely
  \end{itemize}
\end{comment}
Let $(\cM^{nf},\nabla^{nf})$ be a fixed model with the corresponding normal
form $A^0$ and let us also fix a normal solution $\cY_0$ of $A^0$.
Here we want to introduce the Stokes structures, which will characterize the
isomorphism classes of meromorphic connections uniquely, i.e.\ an space which is
isomorphic to the set $\cH(A^0)$ of isomorphism classes of marked pairs.
A great overview of this topic is given by Varadarajan
in~\cite{Varadarajan96linearmeromorphic}.

One of the most important theorems, which is fundamental for the whole chapter,
is the Malgrange-Sibuya Theorem.
It states that the classifying set $\cH(A^0)$ is via an map $\exp$ isomorphic
to the first non abelian cohomology space $H^1(S^1;\Lambda(A^0))=:\St(A^0)$ of
the Stokes sheaf $\Lambda(A^0)$ and will be proven in the first section.
In Section~\ref{sec:mainThm2} we will improve the Malgrange-Sibuya Theorem by
showing that each $1$-cohomology class in $\St(A^0)$ contains a unique
$1$-cocycle of a special form called \emph{the Stokes cocycle}.
We will further show that such cocycles can be identified with their germs at
some special directions, i.e.\ anti-Stokes directions. These germs are called
Stokes germs and for an anti-Stokes direction $\alpha$ do these germs form the
Stokes groups $\Sto_\alpha(A^0)$ (cf.\ Section~\ref{sec:StokesGroup}).
The morphism, which maps each product of Stokes germs to its corresponding
$1$-cocycle will be denoted by $h$.
This will be further improved in Section~\ref{sec:furtherImprovements}, where we
will collect multiple Stokes germs to their product to obtain a more robust
version of the Stokes space.

If one introduces the map $g$, which arises from the theory of
summation\footnote{We will not use the theory of summation and only think of it
  as a black box, although we will roughly discuss the theory of summation in
  Appendix~\ref{app:multisummability}.} and takes an equivalence
class\footnote{resp.\ an ambassador of such a class.} and returns the
corresponding Stokes cocycle in an canonically way one obtains the following
commutative diagram.
\begin{center}
  \begin{tikzpicture}[scale=3]
    % \node[] (modSpcMat) at (0,0.4) {$\cH(\cM^{nf},\nabla^{nf})$};
    \node[] (mat) at (0,-.8) {$\prod_{\theta\in\A}\Sto_\theta(A^0)$};
    \node[] (class) at (0,0) {$\cH(A^0)$};
    % \node[green!40!black] (sheaf) at (3,0.4) {$\St(\cM^{nf})$};
    \node[] (sheaf3) at (1.3,0) {$\St(A^0)$};

    % \draw[thick,double,blue] (modSpcMat) -- (class);
    % \draw[purple,thick,double] (sheaf) -- (sheaf3);

    % \draw[->,green!40!black] (modSpcMat) -- (sheaf)
    %   node[midway,above] {$\exp$};
    \draw[->] (class) -- (mat) node[midway,left] {g};
    \draw[->] (class) -- (sheaf3) node[midway,above] {$\exp$};
    \draw[->] (mat) -- (sheaf3) node[midway, below right] {$h$};
  \end{tikzpicture}
\end{center}\label{page:ofPreDiagram}
This diagram will be enhanced in Section~\ref{sec:theCompleteDiagram} by
adding a couple of isomorphisms.

In the first section of this chapter we will use
\rewrite{Sabbah's} book~\cite[section II]{sabbah2007isomonodromic} as its main
resource together with \cite{babbitt1989local} for the main proof.
In the Sections~\ref{sec:StokesGroup} and~\ref{sec:mainThm2} will
\rewrite{Loday-Richaud's} paper~\cite{Loday1994} and her book
\cite[Sec.4]{Loday2014} be useful.
Stokes groups are also discussed \rewrite{Boalch's} paper~\cite{boalch} (resp.\
his thesis~\cite{thboalch}) which looks only at the single leveled case or the
paper~\cite{Martinet1991} from Martinet and Ramis.
Another very important paper is~\cite{BJL1979Birkhoff} from Balser, Jurkat and
Lutz.

%%% Local Variables:
%%% TeX-master: "Maximilian_Huber-Masters_Thesis-with_notes.tex"
%%% End:

    \ifnum\myHideVariable=0
    \input{chp3sec1-thm1.tex}
    \fi

    \ifnum\myHideVariable=0
    %%%%%%%%%%%%%%%%%%%%%%%%%%%%%%%%%%%%%%%%%%%%%%%%%%%%%%%%%%%%%%%%%%%%%%%%%%%%%%%
\section{The Stokes groups}\label{sec:StokesGroup}
Here we want to introduce the notion of Stokes groups. They are for example
also introduced by Loday-Richaud in~\cite{Loday1994,Loday2014} or section 4
of~\cite{Martinet1991} by Martinet and Ramis.

Let us recall, that the normal form $A^0$ can be written as
$A^0=Q'(t^{-1})+L\frac{1}{t}$ and a normal solution is given by
$\cY_0(t)=t^Le^{Q(t^{-1})}$ (cf.\ Proposition~\ref{prop:fundSolBuilder}), where
\begin{itemize}
\item $Q(t^{-1})=\bigoplus_{j\in\{1,\dots,s\}}q_j(t^{-1})\cdot\id_{n_j}$ and
\item the block structure of $L$ is finer then the structure of $Q$
  (cf.\ Definition~\ref{defn:structureComparison}).
\end{itemize}
Let $\left\{q_1(t^{-1}),\dots,q_s(t^{-1})\right\}$ be the \emph{set of all
determining polynomials of $[A^0]$} and denote by
\[
  \cQ(A^0):=\left\{q_j-q_l
    \mid
    \text{$q_j$ and $q_l$ determining polynomials of $[A^0]$, } q_j \neq q_l
  \right\}
\]
the \emph{set of all determining polynomials of $[\End A^0]$}.
Instead of $q_j-q_l\in\cQ(A^0)$ we will sometimes talk of (ordered) pairs
$(q_j,q_l)\in\cQ(A^0)$.

\begin{defn}\label{defn:determiningPolysOfEndA}
  We call
  \begin{itemize}
    \item $a_{jl}\in\C\backslash\{0\}$ the \emph{leading factor},
    \item $\frac{a_{jl}}{t^{k_{jl}}}$ the \emph{leading
      coefficient} and
    \item $k_{jl}\in\Q$ the \emph{degree}
  \end{itemize}
  of $q_j-q_l\in\cQ(A^0)$ if
  \[
    q_j-q_l\in\left\{\frac{a_{jl}}{t^{k_{jl}}}+h \mid h \in o(t^{-k_{jl}}),
      a_{jl}\neq0
    \right\}\,.
  \]
  \begin{s-rem}
    \begin{enumerate}
      \item It is obvious that $k_{jl}=k_{lj}$ and
        $\frac{a_{jl}}{t^{k_{jl}}}=\frac{-a_{lj}}{t^{k_{lj}}}$.
      \item In Boalch's paper \cite{boalch} (and also in \cite{thboalch}) are
        the \rewrite{degrees of the pairs always incremented by one}.
        We will prefer the \rewrite{other notion}, which is also used in
        Loday-Richaud's paper \cite{Loday1994}.
      \item In Loday-Richaud's book \cite[Def.4.3.6]{Loday2014} $a_{jl}$ is
        negated to be consistent with calculations at $\infty$.
        Here this is not necessary, since we use the clockwise orientation on
        $S^1$ (cf.\ Definition~\ref{defn:antiStokesDir}).
    \end{enumerate}
  \end{s-rem}
  The degrees of the elements in $\cQ(A^0)$ are defined to be  the
  \emph{levels} of $A^0$.
  The set of all levels of $A^0$ will be denoted by
  \[
    \cK=\{k_1<\dots<k_r\} \subset \Q \,.
  \]
  \begin{s-rem}
    The system $[A^0]$ is unramified if and only if $\cK\subset\Z$.
    Since we only want to consider the unramified case, this will be always the
    case.
  \end{s-rem}
\end{defn}

%%%%%%%%%%%%%%%%%%%%%%%%%%%%%%%%%%%%%%%%%%%%%%%%%%%%%%%%%%%%%%%%%%%%%%%%%%%%%%%
\subsection{Anti-Stokes directions and the Stokes group}
\marginnote{\cite[I.4]{Loday1994}}
\begin{defn}
  \marginnote{\cite[130]{hotta2008}, \cite[79]{Loday2014}}
  Let $k\in\N$ and $a\in\C$.
  We say that an exponential $e^{q(t^{-1})}$, where
  $q(t^{-1})\in\frac{a}{t^{k}}+o(t^{-k})$, has \emph{maximal decay in a
  direction $\theta\in S^1$} if and only if $ae^{-ik\tilde\theta}$ is real
  negative. \rewrite{We say that an matrix has maximal decay, if every entry
  has maximal decay.}
  \begin{comment}
    $be^0\in\C$ corresponding to has maximal decay if and only if \PROBLEM[$1$
    has maximal decay?]
  \end{comment}
\end{defn}

On the determining polynomials of $[A^0]$ we define the following (partial)
order relations:
\begin{defn}\label{defn:definingRelations}
  Let $\tilde\theta$ be a determination of $\theta\in S^1$.
  \begin{itemize}
    \item We define the relation
      $\boldmath q_j \underset{\tilde\theta}{\prec} q_l$ to be equivalent to
      the condition
      \begin{einr}
        \rewrite{$e^{(q_j-g_l)(t^{-1})}$ is flat at $0$} in a neighbourhood of
        the direction $\tilde\theta$, i.e.\ if and only if
        $\Re(a_{jl}e^{-ik_{jl}\tilde\theta})<0$.
      \end{einr}
    \item Let us define another relation $\boldmath q_j\myrel{\tilde\theta}q_l$
      equivalent to
      \begin{einr}
        $e^{(q_j-g_l)(t^{-1})}$ is of maximal decay in the direction
        $\tilde\theta$,
      \end{einr}
      which by itself is equivalent to
      \begin{einr}
        $a_{jl}e^{-ik_{jl}\tilde\theta}$ is a real negative
        number, i.e.\ $q_j \underset{\tilde\theta}{\prec} q_l$ and
        $\Im(a_{jl}e^{-ik_{jl}\tilde\theta})=0$.
      \end{einr}
  \end{itemize}
  \begin{s-rem}
    In the unramified case do these relations not depend on the determination
    $\tilde\theta$ of $\theta$. As a consequence we will only write
    $\underset{\theta}{\prec}$ and $\myrel{\theta}$.
  \end{s-rem}
\end{defn}
To understand the previous definition better, it is convenient to look closer at
functions of the form $f:\theta\mapsto ae^{-ik\theta}$, $k\in\Z$, corresponding
to some pair $(q_j,q_l)$.
Write $a$ as $a=|a|e^{i\arg(a)}$, thus the function $f$ writes as
\begin{align*}
  f(\theta)&=|a|e^{i(\arg(a)-k\theta)}
  \\&=|a|(\cos(\arg(a)-k\theta) + i\sin(\arg(a)-k\theta)) \,.
\end{align*}
In the Figure~\ref{fig:functionF}, we illustrate the real and the imaginary
part of $f$.
\begin{figure}[h!] %{{{
  \begin{flushright}
    \tikzmarkc{n2}{purple} here is $\Im(ae^{-ik\theta})=0$
  \end{flushright}
  \begin{center}
    \begin{tikzpicture}[scale=4]
      \pgfmathsetmacro{\k}{3}
      \pgfmathsetmacro{\argA}{0.3}
      \pgfmathsetmacro{\absA}{0.4}

      \begin{scope}[thick]
        \clip (-0.6,{\absA+0.1}) rectangle (2.6,{-\absA-0.1});
        \foreach \x in {-2,-1,...,3}{
          \pgfmathsetmacro{\s}{{\argA + \x / \k * 2 - 1/2/\k}};
          \fill[blue!10!white]
            ({\s + 2/2/\k},0) sin ({\s + 3/2/\k},{-\absA})
                              cos ({\s + 4/2/\k},0);
          \draw[blue!40!black] (\s,0) sin ({\s + 1/2/\k},\absA)
                       cos ({\s + 2/2/\k},0)
                       sin ({\s + 3/2/\k},{-\absA})
                       cos ({\s + 4/2/\k},0);
          \fill[white] ({\s + 3/2/\k},{-\absA}) circle (1pt);
          \fill[blue!40!black] ({\s + 3/2/\k},{-\absA}) circle (.4pt);
          \pgfmathsetmacro{\s}{{\argA + \x / \k * 2}};
          \draw[purple] (\s,0) sin ({\s + 1/2/\k},\absA)
                       cos ({\s + 2/2/\k},0)
                       sin ({\s + 3/2/\k},{-\absA})
                       cos ({\s + 4/2/\k},0);
          \fill[white] (\s,0) circle (1pt);
          \fill[purple] (\s,0) circle (.4pt);
          \fill[white] ({\s + 2/2/\k},0) circle (1pt);
          \fill[purple] ({\s + 2/2/\k},0) circle (.4pt);
        }
      \end{scope}
      \draw [blue!40!black,dashed]
        ({\argA+4/\k + 1/2/\k},{-\absA}) -- ({\argA+4/\k + 1/2/\k},0);
      \draw [blue!40!black,dashed]
        ({\argA+4/\k + 3/2/\k},{-\absA}) -- ({\argA+4/\k + 3/2/\k},0);
      \draw [blue!40!black, thick
            ,decorate
            ,decoration={brace,mirror,amplitude=10pt}
            ,xshift=0pt
            ,yshift=0pt]
        ({\argA+4/\k + 1/2/\k},{-\absA}) -- ({\argA+4/\k + 3/2/\k},{-\absA})
        node [midway,yshift=-7pt] {$\tikzmark{g0}$};

      \node at ({\argA+2 - 3/\k},0) {$\tikzmark{f-1}$};
      \node at ({\argA+2 - 2/\k},0) {$\tikzmark{f0}$};
      \node at ({\argA+2 - 1/\k},0) {$\tikzmark{f1}$};

      \node at ({\argA-2/\k + 2/2/\k},{-\absA}) {$\tikzmark{e-1}$};
      \node at ({\argA + 2/2/\k},{-\absA}) {$\tikzmark{e0}$};
      \node at ({\argA+2/\k + 2/2/\k},{-\absA}) {$\tikzmark{e1}$};
      \node at ({\argA+4/\k + 2/2/\k},{-\absA}) {$\tikzmark{e2}$};

      \draw[-latex'] (-0.7,0) -- (2.7,0);
      \draw[dotted] (-0.6,\absA)node[left,font=\tiny] {$|a|$} -- (2.6,\absA);
      \draw[dotted] (-0.6,{-\absA})node[left,font=\tiny] {$-|a|$} -- (2.6,{-\absA});

      \draw[dotted] ({-0.5},{-\absA}) -- ({-0.5},{\absA + 0.1})
        node [above,font=\tiny,] {-0.5 \pi};
      \draw[-latex'] ({0},{-\absA-0.1}) -- ({0},{\absA + 0.2});
      \foreach \x in {0.5,1,...,2.5}{%
        \draw[dotted] ({\x},{-\absA}) -- ({\x},{\absA + 0.1})
          node [above,font=\tiny,] {\x \pi};
      }

      \draw[thick, dotted, green!50!black] (\argA,{-\absA}) -- (\argA,{\absA + 0.1})
        node [above,font=\tiny,] {$\frac{\arg(a)}{k}$};
    \end{tikzpicture}
  \end{center}
  \begin{flushright}
    \tikzmarkc{n3}{blue} here is $q_j\underset{\theta}{\prec}q_l$
  \end{flushright}
  \begin{flushright}
    \tikzmarkc{n1}{blue} here is $q_j\myrel{\theta}q_l$
  \end{flushright}
  \begin{tikzpicture}[remember picture,overlay]
    \draw[->,blue!50!white,thick] (n3) to[out=150,in=270] (g0);
    \draw[->,purple!50!white,thick] (n2) to[out=240,in=70] (f-1);
    \draw[->,purple!50!white,thick] (n2) to[out=253,in=120] (f0);
    \draw[->,purple!50!white,thick] (n2) to[out=266,in=70] (f1);
    \draw[->,blue!50!white,thick] (n1) to[out=180,in=270] (e-1);
    \draw[->,blue!50!white,thick] (n1) to[out=170,in=270] (e0);
    \draw[->,blue!50!white,thick] (n1) to[out=160,in=270] (e1);
    % \draw[->,blue!40!black,thick] (n1) to[out=150,in=270] (e2);
  \end{tikzpicture}
  \caption{In this plot is the real part of $f(\theta)=ae^{-ik\theta}$,
    corresponding to some pair $(q_j,q_l)$, in
    \textcolor{blue!60!white}{blue} and the imaginary part in
    \textcolor{purple}{purple} sketched.
  }\label{fig:functionF}
\end{figure} %}}}

The graphs, corresponding to the flipped pair $(q_l,q_j)$ are then obtained by
the transformation $\arg(a)\to\arg(-a)=\arg(a)+\pi$, i.e.\ the shift by
$\frac{\pi}{k}$ to the right. This $\frac{\pi}{k}$ is exactly a half period,
thus the new graphs are obtained by \rewrite{mirroring at the line $t=0$.}
\begin{rem}\label{rem:relationDistanceCondition}
  Let $k_{jl}$ be the degree of $q_j-q_l$.
  It is easy to see (cf.\ Figure~\ref{fig:functionF}), that the condition
  $q_j \underset{\theta}{\prec} q_l$ is equivalent to
  \begin{einr}
    there is a $\theta'\in U(\theta,\frac{\pi}{k_{jl}})$ such that
    $q_j\myrel{\theta}q_l$.
  \end{einr}
\end{rem}
Let us now use the defined relations to say, which are the interesting
directions of $S^1$.
\begin{defn}\label{defn:antiStokesDir}
  % \marginnote{See \cite[Def.I.4.5]{Loday1994}(for the ramified case)
  %   \cite[Def.3.2]{boalch}}
  Let $\theta\in S^1$ be an direction.
  \begin{enumerate}
    \item $\theta$ is an \emph{anti-Stokes direction} if there is at least one
      pair $(q_j,q_l)$ in $\cQ(A^0)$, which satisfies $q_j\myrel{\theta}q_l$.

      Let $\A=\{\alpha_1,\dots,\alpha_{\nu}\}$ denote the set of all
      anti-Stokes directions in a clockwise ordering. For a uniform
      notation later, \rewrite{define $\A$ to contain a single, arbitrary
        direction if $\cK=\{0\}$.}
      \begin{s-rem}
        The clockwise ordering is chosen, similar to Loday-Richaud's
        paper~\cite{Loday1994}, since the calculations are then compatible with
        the calculations, which look at $\infty$ and take a
        counterclockwise ordering.
        Boalch uses in~\cite{boalch} and~\cite{thboalch} the inverse
        ordering, but looks also at $0$, thus there \rewrite{might be some}
        incompatibilities.
        In Loday-Richaud's book~\cite{Loday2014} this problem is solved by
        an additional minus sign for some coefficient.
      \end{s-rem}
    \item The direction $\theta$ is a \emph{Stokes direction} if there is at
      least one pair $(q_j,q_l)$ in $\cQ(A^0)$, which satisfies neither
      $q_j\underset{\theta}{\prec} q_l$ nor $q_l\underset{\theta}{\prec} q_j$.
      \begin{comment}
        Let $\S=\{\sigma_1<\cdots<\sigma_\mu\}$ be the set of Stokes directions.
      \end{comment}
  \end{enumerate}
\end{defn}
We will use the Greek letter $\alpha$ whenever we want to emphasize that a
direction is an anti-Stokes direction.
For generic directions, we will use $\theta$.
In fact will most of the following definitions and
constructions work for every $\theta\in S^1$, but the Stokes group
(cf.\ Definition~\ref{defn:stokesGroup}) for example will be trivial
for every $\theta\notin\A$.
Thus the interesting directions are  only the anti-Stokes directions
$\alpha\in\A$.

\begin{lem}\label{lem:rotationalSym}%\label{rem:rotationalSymPrime}
  Let $\alpha\in\A$ together with a pair $(q_j-q_l)(t^{-1})\in\cQ(A^0)$ of
  degree $k_{jl}$, such that $q_j\myrel{\alpha}q_l$ be given.
  We then know for every $m\in\N$ that
  \[
  \underset{\alpha'}{
    \underset{\text{\rotatebox[origin=c]{-90}{$=:$}}}{%
      \underbrace{\alpha+m\frac{\pi}{k_{jl}}}}}
  \in \A \,.
  \]
  Especially is either $q_j\myrel{\alpha'}q_l$ (in the case, when $m$ is even)
  or $q_l\myrel{\alpha'}q_j$ (when $m$ is uneven) satisfied
  (see Figure~\ref{fig:functionF}).
  \begin{s-cor}
    It follows that in the case $\cK=\{k\}$, the set $\A$ has
    $\frac{\pi}{k}$-rotational symmetry.
  \end{s-cor}
\end{lem}
\begin{proof}
  \marginnote{\cite[8]{thboalch}}
  Let $(j,l)$ be a pair such that $q_j\myrel{\alpha}q_l$, i.e.\ such that
  $a_{jl}e^{-ik_{jl}\alpha}\in\R_{<0}$.
  Hence for $m\in\N$
  \begin{align*}
    a_{jl}e^{-ik_{jl}\left(\alpha+m\frac{\pi}{k_{jl}}\right)}
    &=a_{jl}e^{-ik_{jl}\alpha}e^{-im\pi}
    = \begin{cases}
      a_{jl}e^{-ik_{jl}\alpha}\in\R_{<0}
        & \text{, if $m$ is even}
    \\-a_{jl}e^{-ik_{jl}\alpha}\in\R_{>0}
        & \text{, if $m$ is uneven}
    \end{cases}
  \end{align*}
  is in the case when $m$ is even, also real and negative. In the other case,
  when $n$ is uneven, we use that $a_{jl}=-a_{lj}$ and $k_{jl}=k_{lj}$ to
  obtain
  $a_{lj}e^{-ik_{lj}\left(\alpha+m\frac{\pi}{k_{lj}}\right)} \in\R_{<0}$.

  Thus, for $\alpha':=\alpha+m\frac{\pi}{k_{jl}}$, we have $\alpha'\in\A$ since
  \begin{itemize}
    \item $q_j\myrel{\alpha'}q_l$ when $m$ is even or
    \item $q_l\myrel{\alpha'}q_j$ when $m$ is uneven.
  \end{itemize}
\end{proof}

As a \rewrite{subgroup of the stalk at $\theta$} of the in
Definition~\ref{defn:StokesSheaf} defined Stokes sheaf $\Lambda(A^0)$ we
define the Stokes group as follows.
\begin{defn}\label{defn:stokesGroup}
  Define the \emph{Stokes group}
  \[
    \Sto_\theta(A^0):=
    \left\{\phi_\theta\in\Lambda_\theta(A^0)
      \mid \phi_\theta \text{~has maximal decay at~} \theta
    \right\}
  \]
  whose elements are called \emph{Stokes germs}.
  \TODO[This is in fact a group, since\dots]
  \begin{s-rem}
    For $\theta\notin\A$ the group $\Sto_\theta(A^0)$ is trivial, since at
    $\theta$ no flat isotropy has maximal decay, but the identity.
  \end{s-rem}
\end{defn}

%%%%%%%%%%%%%%%%%%%%%%%%%%%%%%%%%%%%%%%%%%%%%%%%%%%%%%%%%%%%%%%%%%%%%%%%%%%%%%%
\subsection{Stokes matrices}\label{sec:matrixReps}
\marginnote{\cite[9f]{thboalch}, \cite[??]{Loday1994}}
Stokes matrices, which Wasow calls in his book~\cite{wasow2002asymptotic}
Stokes multipliers and Boalch calls Stokes factors in~\cite{boalch,thboalch},
arise either
\begin{einr}
  as faithful representations of Stokes germs
\end{einr}
or, if one starts by comparing the actual fundamental solutions on arcs, as
\begin{einr}
  the matrices describing the blending between two adjacent fundamental
  solutions, with some additional assumptions
  (cf.\ Definition~\cite[80]{Loday2014}).
\end{einr}
\begin{defn}\label{defn:groupOfFaithfullReps}
  Let us use
  \[
    \bdelta_{jl}:=
    \begin{cases}
      0 \in \C^{n_j\times n_l} & \text{,~if~} j\neq l
    \\\id \in \C^{n_j\times n_l} & \text{,~if~} j=l
    \end{cases}
  \]
  as a block version of Kronecker's delta corresponding to the structure of the
  normal solution $\cY_0$, which was fixed.
  Define the group
  \begin{align*}
    \SSto_\theta(A^0)= \Big\{K=(K_{jl})_{j,l\in\{1,\dots,s\}}\in\GL_n(\C) \mid
      K_{jl}=\bdelta_{jl} \text{~unless~} q_j\myrel{\theta}q_l \Big\}
  \end{align*}
  of all \emph{Stokes matrices} of $A^0$ in the direction $\theta$.
  They will arise as a faithful representation
  (cf.\ Section~\ref{sec:faithRepre}) of $\Sto_\theta(A^0)$.
  \begin{s-rem}\label{rem:groupOfFaithfullReps}
    There is obviously a bijection
    $\vartheta_\theta: \prod_{q_j\myrel{\theta}q_l}\C^{n_j \cdot n_l}
    \overset{\cong}{\longrightarrow} \SSto_\theta(A^0)$.
  \end{s-rem}
\end{defn}

\begin{prop}\label{prop:representation}
  \marginnote{\cite[Def.I.4.7]{Loday1994}\\\cite[78f]{Loday2014}}
  In this situation is
  \begin{align*}
    \rho_{\theta}:\Sto_\theta(A^0)&\longrightarrow\SSto_\theta(A^0)
    \\\phi_\theta
    &\longmapsto
    C_{\phi_\theta}:=\cY_{0}\phi_\theta\cY_{0}^{-1}
  \end{align*}
  \marginnote{Boalch uses $C_{\phi_\theta}:=\cY_{0}^{-1}\phi_\theta\cY_{0}$}
  an isomorphism which maps a germ of $\Sto_\theta(A^0)$ to the corresponding
  Stokes matrix $C_{\phi_\theta}$ such that
  \begin{equation}\label{eq:representation}
    \phi_\theta(t)\cY_{0}(t)=\cY_{0}(t)C_{\phi_\theta}
  \end{equation}
  near $\theta$.
  The matrix $C_{\phi_\theta}$ is then called a \emph{representation of
  $\phi_\theta$}.
  \begin{s-rem}\label{rem:representation}
    \begin{enumerate}
      \item In the ramified case does this morphism depend on the choice of the
        determination $\tilde\theta$ of $\theta$ and the corresponding choice of
        a realization of the fundamental solution with that determination of the
        argument near the direction $\theta$
        (cf.~\cite{Loday1994} or~\cite[78f]{Loday2014}).
      \item \marginnote{\cite[Defn.I.4.7]{Loday1994}}
        This construction defines also a morphism, which takes a germ
        $\phi_\theta\in\Lambda_\theta(A^0)\supset\Sto_\theta(A^0)$ into its
        unique representation matrix
        \[
          C_{\phi_\theta} \in
          \underset{\SSto_\theta(A^0)}{%
            \underset{\text{\rotatebox[origin=c]{90}{$\subset$}}}{%
              \hat\SSto_{\theta}(A^0)}}
          :=
          \left\{(K_{jl})_{j,l\in\{1,\dots,s\}}
            \in \GL_n(\C)\mid K_{jl}=\bdelta_{jl} \text{ unless }
            q_j \underset{\theta}{\prec} q_l \right\}
        \]
        and there is a bijection $\hat\vartheta_\theta:
        \prod_{q_j\underset{\theta}{\prec}q_l}\C^{n_j \cdot n_l}
        \overset{\cong}{\longrightarrow} \hat\SSto_\theta(A^0)$.
        \begin{comment}
          Does this define a local-constant sheaf
          \[
            I\mapsto \hat\SSto_{I}(A^0)
            :=
            \left\{(K_{jl})_{j,l\in\{1,\dots,s\}}
              \in \GL_n(\C)\mid K_{jl}=\bdelta_{jl} \text{ unless }
              q_j \underset{\theta}{\prec} q_l
              \text{ for some } \theta\in I\right\}
          \]
          and a skyscraper sheaf
          \[
            I\mapsto \SSto_{I}(A^0) \,.
          \]
          \PROBLEM
        \end{comment}
    \end{enumerate}
  \end{s-rem}
\end{prop}
\begin{proof}
  It is well known (cf.\ \cite[10]{thboalch}), that the morphism
  $\rho_{\theta}$, i.e.\ conjugation by the fundamental solution, relates
  solutions $\phi_\theta$ of $[\End(A^0)]=[A^0,A^0]$ to solutions of $[0,0]$
  which are the constant matrices $\GL_n(\C)$.
  Thus we have to show, that the image of $\Sto_\theta(A^0)$ under
  $\rho_{\theta}$ is $\SSto_\theta(A^0)$.

  To see that the obtained matrix has the necessary zeros, to lie in
  $\SSto_{\theta}(A^0)$ we look at Equation (\ref{eq:representation}) and
  deduce
  \begin{equation}\label{eq:repProof1}
    \phi_\theta(t)
    =t^L e^{Q(t^{-1})}C_{\phi_\theta}e^{-Q(t^{-1})}t^{-L}
  \end{equation}
  with the given choice of the argument near $\theta$.
  After decomposing $C_{\phi_\theta}$ into
  \begin{align*}
    C_{\phi_\theta}&=1_n+\begin{pmatrix}
      c_{(1,1)} & c_{(1,2)} & \cdots &\\
      c_{(2,1} & \ddots\\
      \vdots \\
      & & & c_{(s,s)}
    \end{pmatrix}
  \\&=1_n+
    \underset{C_{\phi_\theta}^{(1,1)}}{\underbrace{%
      \begin{pmatrix}
        c_{(1,1)} & 0 & \cdots &\\
        0\\
        \vdots&\\
        &
      \end{pmatrix}
    }}
    +
    \underset{C_{\phi_\theta}^{(1,2)}}{\underbrace{%
      \begin{pmatrix}
        0 & c_{(1,2)} & 0 & \cdots\\
        & 0 &\\
        &\vdots\\
        &
      \end{pmatrix}
    }}
    +\cdots+
    \underset{C_{\phi_\theta}^{(s,s)}}{\underbrace{%
      \begin{pmatrix}
        &\\
        & & & \vdots\\
        & & & 0\\
        & \cdots & 0 & c_{(s,s)}
      \end{pmatrix}
    }}
  \\&=1_n+\sum_{(l,j)}C_{\phi_\theta}^{(l,j)}
  \end{align*}
  where the $c_{(j,l)}$ are blocks\marginnote{One can ignore the block
  structure by using $1\times1$ sized blocks. But one looses the uniqueness of
  the $q_j$'s.} of size $n_j\times n_l$ which correspond to the structure of
  $Q$. After rewriting the Equation (\ref{eq:repProof1}) we get
  \[
    \phi_\theta=
      t^L\left(
        1_n+\sum_{(l,j)}C_{\phi_\theta}^{(l,j)}e^{(q_l-q_j)(t^{-1})}
      \right)t^{-L} \,.
  \]
  \begin{comment}
    \begin{align*}
      \phi_\theta(t)
      &=t^Le^{Q(t^{-1})}\left(
        1_n+C_{\phi_\theta}
      \right)e^{-Q(t^{-1})}t^{-L}
    \\&=t^Le^{Q(t^{-1})}\left(
        1_n+\sum_{(l,j)}C_{\phi_\theta}^{(l,j)}
      \right)e^{-Q(t^{-1})}t^{-L}
    \\&=t^L\left(
        1_n+\sum_{(l,j)}e^{Q(t^{-1})}C_{\phi_\theta}^{(l,j)}e^{-Q(t^{-1})}
      \right)t^{-L}
    \\&=t^L\left(
          1_n+\sum_{(l,j)}C_{\phi_\theta}^{(l,j)}e^{(q_l-q_j)(t^{-1})}
        \right)t^{-L} \,.
    \end{align*}
  \end{comment}
  Thus, for $\phi_{\theta}$ to be flat in direction $\theta$, it is
  necessary and sufficient that if $e^{(q_l-q_j)(t^{-1})}$ does not have
  maximal decay in direction $\theta$ the corresponding
  block $C_{\phi_\theta}^{(l,j)}$ vanishes.
  Thus we have seen, that $C_{\phi_\theta}$ is an
  element of $\SSto_\theta(A^0)$.

  The \textbf{surjectivity} can \rewrite{now be seen easily} since every
  constant matrix with zeros at the necessary positions characterizes a unique
  element of $\Sto_\theta(A^0)$:
  \begin{einr}
    Let $C=1_n + \sum_{(l,j)\mid q_j\myrel{\theta}q_l} C^{(l,j)}$ be an element
    of $\SSto_\theta(A^0)$.
    Then is a pre-image of $C$ given by
    $t^Le^{Q(t^{-1})}Ce^{-Q(t^{-1})}t^{-L}$ which lies in $\Sto_\theta(A^0)$,
    since it satisfies the condition discussed above.
  \end{einr}

  The map $\rho_{\tilde\theta}$ is also \textbf{injective}, since it is the
  conjugation by an invertible matrix.
\end{proof}

From the calculations in the proof it is clear that
\begin{enumerate}
  \item for $j=l$ the (diagonal) blocks $C_{\phi_\theta}^{(l,j)}$ vanish since
    $q_l-q_j=0$ does not have maximal decay and
  \item if $e^{q_j-q_l}$ has has maximal decay, then $e^{q_l-q_j}$ has not.
    Thus if $C_{\phi_\theta}^{(l,j)}$ is not equal to zero, the block
    $C_{\phi_\theta}^{(j,l)}$ is necessarily zero.
\end{enumerate}
This implies that the matrix $C_{\phi_\theta}$ is unipotent, and \rewrite{hence
is $\Sto_\theta(A^0)$ is a unipotent Lie group.}
\begin{prop}
  Some of the above results are also true for the groups
  $\hat\SSto_\theta(A^0)$. One could use these, to see that the sheaf
  $\Lambda(A^0)$ is a piecewise-constant sheaf of non-Abelian unipotent Lie
  groups (cf.~\cite[Prop.2.1]{Loday2004}).
\end{prop}

One can use the Stokes matrices to give an alternative characterization of
Stokes germs:
\begin{einr}
  a germ $\phi_\theta\in\Lambda_\theta(A^0)$ is in $\Sto_\theta(A^0)$ if and
  only if there exists a $C\in\SSto_\theta(A^0)$ such that
  $\phi_\theta=\cY_{0}C\cY_{0}^{-1}$.
\end{einr}
Formulated is this in the following corollary.
\begin{cor}
  \marginnote{\cite[Def.I.4.12]{Loday1994}}
  A germ $\phi_\theta\in\Lambda_\theta(A^0)$ is a Stokes germ, i.e.\ an element
  in $\Sto_\theta(A^0)$, if and only if it has a representation
  $C_{\phi_\theta}$ where
  \[
    C_{\phi_\theta}=
        1_n + \sum_{(l,j)\mid q_j\myrel{\theta}q_l}C_{\phi_\theta}^{(l,j)}
  \]
  and the $C_{\phi_\theta}^{(l,j)}$ have the necessary block structure, i.e.\ it
  is in $\SSto_\theta(A^0)$.
  \begin{s-rem}
    In Loday-Richaud's book~\cite[78]{Loday2014} are the elements of
    $\Sto_\theta(A^0)$ actually characterized as the flat transformations, such
    that Equation (\ref{eq:representation}) is satisfied for some unique
    constant invertible matrix $C\in\SSto_\theta(A^0)$.
  \end{s-rem}
\end{cor}

\begin{defn}
  We denote the set of \emph{levels of the germ}
  $\phi_{\theta}\in\Lambda_\theta(A^0)$ by
  \[
    \cK(\phi_\theta):= \left\{\deg(q_j-q_l)\mid C_{\phi_\theta}^{(l,j)}\neq0
      \text{ in some representation of }\phi_\theta\right\} \subset \cK \,.
  \]
  A germ $\phi_\theta$ is called a \emph{$k$-germ} when
  $\cK(\phi_{\theta})\subset\{k\}$, i.e.\ it has at most the level $k$.
\end{defn}

\begin{comment}
  \PROBLEM[This would be good]
  \begin{lem}
    Every $k$-germ in direction $\alpha$ can be extended to the Stokes arc
    $U(\frac{\pi}{k},\alpha)$.
    \begin{cor}
        Every germ $\phi_\alpha\in\Sto_\alpha(A^0)$ can be extended to the arc
        $U(\frac{\pi}{\max\cK(\phi_\alpha)},\alpha)$, i.e.\ there is a section
        $\phi\in\Gamma\left(U(\frac{\pi}{\max\cK(\phi_\alpha)},\alpha),
          \Lambda(A^0)\right)$ which has $\phi_\alpha$ as its germ at $\alpha$.
    \end{cor}
  \end{lem}
  Let $\phi_\alpha$ be a \textbf{simple} $k$-germ in the sense that it is build
  from a single block, i.e.\
  \[
    \phi_\theta=
      t^L\left(
        1_n+C_{\phi_\theta}^{(l,j)}e^{(q_l-q_j)(t^{-1})}
      \right)t^{-L} \,.
  \]
  for some pair $(j,l)$. Assume also that the block has size $1\times1$.

  Let $\phi$ be the extension of $\phi_\alpha$ around alpha, i.e.\ the matrix
  which has germ $\phi_\alpha$ at $\alpha$ and which solves $[\End A^0]$ and is
  multiplicatively flat.

  \begin{einr}
    The system $[\End A^0]$ is
    \[
    \frac{dF}{dt}=A^0F-FA^0 \,.
    \]
  \end{einr}

  \textbf{Question:} Which form has the extension $\phi$ around $\alpha$ of
  the germ $\phi_\alpha$, does it retain the structure?
  \begin{einr}
    \begin{enumerate}
    \item Look at a diagonal element $\phi^{j,j}$ an $(j,j)$:
      \begin{itemize}
      \item it satisfies $\phi_\alpha^{j,j}=1$ and
      \item is satisfies some complicated equation
        \textbf{Question:} Is it constantly $1$
        \begin{einr}
          hopefully yes
        \end{einr}
      \end{itemize}
    \item Look at an off-diagonal position at $(j,l)$:
      \begin{itemize}
      \item[case a:] $\phi_\alpha^{j,l}\neq 0$
      \item[case b:] $\phi_\alpha^{j,l}=0$
      \end{itemize}
    \end{enumerate}
  \end{einr}
\end{comment}

%%%%%%%%%%%%%%%%%%%%%%%%%%%%%%%%%%%%%%%%%%%%%%%%%%%%%%%%%%%%%%%%%%%%%%%%%%%%%%%
\subsection{Decomposition of the Stokes group by levels}
\marginnote{\cite{Loday1994}, \cite[362ff]{Martinet1991}}
\rewrite{The goal of this section is, to introduce a filtration of
$\Lambda(A^0)$, which will be restricted to $\Sto_\theta(A^0)$ and defines there
a filtration.} This leads to a decomposition of $\Sto_\theta(A^0)$ into a
semidirect product (cf.\ Proposition~\ref{prop:filtrationOfStokesGroup}).

Let us introduce a couple of notations and definitions, which coincide with the
notations used in Loday-Richaud's paper \cite{Loday1994}.
Another good resource, which uses similar notations, is for example the
paper~\cite[362f]{Martinet1991} from Martinet and Ramis.
\begin{notations}
  \marginnote{\cite[Not.I.4.15]{Loday1994},\\\cite[362]{Martinet1991}}
  For every level $k\in\cK$ and direction $\theta\in S^1$ we set
  \begin{itemize}
    \item $\Lambda^{k}(A^0)$ as the subsheaf of $\Lambda(A^0)$ of all germs,
      which are generated by $k$-germs;
    \item $\Lambda^{\leq k}(A^0)$ (resp. $\Lambda^{<k}(A^0)$ or
      $\Lambda^{\geq k}(A^0)$) as the subsheaf of $\Lambda(A^0)$ generated by
      $k'$-germs for all $k'\leq k$ (resp. $k'<k$ or $k'\geq k$).
  \end{itemize}
  Let $\star\in\{k,<k,\leq k,\dots\}$.
  The restrictions to $\Sto_\theta$ yield the groups
  \[
    \Sto_\theta^\star(A^0):=\Sto_\theta(A^0)\cap\Lambda_\theta^{\star}(A^0)
  \]
  and let us also define $\SSto^\star_\theta(A^0)$ as the groups of
  representations, which correspond to elements of  $\Sto^\star_\theta(A^0)$.
\end{notations}
Corresponding to the definitions above, one can define
$\A^\star:=\left\{\alpha\in\A\mid\Sto_\alpha^\star(A^0)\neq\{\id\}\right\}$
for $\star\in\{k,<k,\leq k,\dots\}$ and we say that \emph{$\alpha$ is bearing
the level $k$} if $\alpha\in\A^k$.
\begin{rem}
  It is clear that for every $k\in\cK$ we have the canonical inclusions
  $\A^k\hookrightarrow\A^{\leq k}$ and $\A^{<k}\hookrightarrow\A^{\leq k}$.
\end{rem}
Sometimes it is also useful to talk about the \emph{set of levels beared by an
direction $\alpha\in\A$}:
\[
  \cK_\alpha:=\left\{k\in\cK\mid\Sto_\alpha^k(A^0)\neq\{\id\}\right\} \,.
\]
\begin{cor}
  The Lemma~\ref{lem:rotationalSym} implies that from $k\in\cK_{\alpha}$ follows
  that $k\in\cK_{\alpha+m\frac{\pi}{k}}$ for $m\in\N$.
\end{cor}
Let us now study the sheaves $\Lambda^\star(A^0)$ and discuss how they
correlate \rewrite{and how they can be composed from the others}.

The following proposition can be found as~\cite[Prop.I.5.1]{Loday1994} and the
key-statement is also given in~\cite[Prop.4.10]{Martinet1991}.
\begin{prop}\label{prop:PropertiesOfStokesSheafSplitting}
  \marginnote{\cite[Prop.I.5.1]{Loday1994}}
  For any level $k\in\cK$ one has that $\Lambda^{k}(A^0)$,
  $\Lambda^{\leq k}(A^0)$ and $\Lambda^{<k}(A^0)$ are sheaves of subgroups of
  $\Lambda(A^0)$ and the sheaf $\Lambda^k(A^0)$ is normal in
  $\Lambda^{\leq k}(A^0)$.
  \begin{comment}
    A subgroup $N$ is normal in $G$ ($N\vartriangleleft G$) if it is stable
    under conjugation, i.e.
    \[
      N\vartriangleleft G \Leftrightarrow \forall n\in N \forall g\in G,
      gng^{-1}\in N ,.
    \]
  \end{comment}

  \marginnote{\cite[Proposition 10]{Martinet1991}}
  \rewrite{We even know more,} let
  \begin{itemize}
    \item $i:\Lambda^k(A^0)\hookrightarrow\Lambda^{\leq k}(A^0)$ be the
      canonical inclusion and
    \item $p:\Lambda^{\leq k}(A^0)\twoheadrightarrow\Lambda^{<k}(A^0)$ be the
      truncation to terms of levels $<k$.
  \end{itemize}
  Then does the exact sequence of sheaves
  \[
    1\longrightarrow\Lambda^k(A^0)
    \overset{i}\longrightarrow\Lambda^{\leq k}(A^0)
    \overset{p}\longrightarrow\Lambda^{<k}(A^0)
    \longrightarrow 1 \,,
  \]
  split.
\end{prop}
From the splitting of the sequence, we obtain immediately the following
decomposition into a semidirect product.
\begin{cor}\label{cor:factorStokesGerms}
  \marginnote{\cite[Cor.I.5.2]{Loday1994}}
  For any $k\in\cK$, there are the two following ways of factoring
  $\Lambda^{\leq k}(A^0)$ in a semidirect product:
  \begin{align*}
    \Lambda^{\leq k}(A^0)&\cong \Lambda^{<k}(A^0)\ltimes\Lambda^{k}(A^0)
  \\                    &\cong \Lambda^{k}(A^0)\ltimes\Lambda^{<k}(A^0)\,.
  \end{align*}
  This means that any germ $f^{\leq k}\in\Lambda^{\leq k}(A^0)$ can be uniquely
  written as
  \begin{itemize}
    \item $f^{\leq k}=f^{<k}g^k$, where $f^{<k}\in\Lambda^{<k}$ and
      $g^k\in\Lambda^k$, or
    \item $f^{\leq k}=f^kf^{<k}$, where $f^k\in\Lambda^k$ and
      $f^{<k}\in\Lambda^{<k}$.
  \end{itemize}
  \begin{s-rem}\label{rem:algFactorization}
    \marginnote{\cite[Cor.I.5.2(ii)]{Loday1994}}
    We can get the factor $f^{<k}$ common to both factorizations by truncation
    of $f^{\leq k}$ to terms of level $<k$, i.e.\ by applying the map $p$ from
    Proposition~\ref{prop:PropertiesOfStokesSheafSplitting}.
    This truncation can explicitly be achieved, in terms of Stokes matrices, by
    keeping in representations $1+\sum C^{(j,l)}$ of $f^{\leq k}$ only the
    blocks $C^{(j,l)}$ such that $\deg(q_j-q_l)<k$.

    A factorization algorithm could then be:
    \begin{einr}
      get the factor $f^{<k}$ common to both factorizations by truncation of
      $f^{\leq k}$ to terms of level $<k$ and set $g^k:=(f^{<k})^{-1}f^{\leq k}$
      and $f^k:=f^{\leq k}(f^{<k})^{-1}$.
    \end{einr}
  \end{s-rem}
  This decomposition in a semidirect product can be extended to all levels,
  since $\Lambda^{<k}(A^0)=\Lambda^{\leq\max\{k'\in\cK\mid k' < k\}}$.
  Thus
  \[
    \Lambda(A^0)\cong\underset{k\in\cK}\bigltimes\Lambda^k(A^0) \,,
  \]
  where the semidirect product is taken in an ascending or descending order of
  levels.
\end{cor}
\begin{rem}
  Loday-Richaud states in her paper \cite[Prop.I.5.3]{Loday1994} the following
  proposition, which is a more general version of
  Proposition~\ref{prop:PropertiesOfStokesSheafSplitting}.
  \begin{s-prop}
    \marginnote{\cite[Prop.I.5.3]{Loday1994}}
    For any levels $k$,$k'\in\cK$ with $k'<k$ one has:
    \begin{enumerate}
      \item the sheaf $\Lambda^{\geq k'}(A^0)\cap\Lambda^{\leq k'}(A^0)$ is
        normal in $\Lambda^{\leq k}(A^0)$;
      \item the exact sequence of sheaves
        \[
          1\longrightarrow\Lambda^{\geq k'}(A^0)\cap\Lambda^{\leq k}(A^0)
          \overset{i}\longrightarrow\Lambda^{\leq k}(A^0)
          \overset{p}\longrightarrow\Lambda^{<k'}(A^0)
          \longrightarrow 1 \,,
        \]
        where
        \begin{itemize}
          \item $i$ is the canonical inclusion and
          \item $p$ is the truncation to terms of levels $<k'$,
        \end{itemize}
        splits.
    \end{enumerate}
    \TODO[is $\Lambda^{\geq k'}(A^0)\cap\Lambda^{\leq k}(A^0)=\Lambda^k(A^0)$
    and thus the first proposition a corollary of this?]
  \end{s-prop}
  We can use this proposition to follow (cf.\ \cite[Cor.I.5.4]{Loday1994}) that
  \begin{enumerate}
    \item
      the filtration
      \[
        \Lambda^{k_r}(A^0)
        =
        \Lambda^{\geq k_r}(A^0)
        \subset
        \Lambda^{\geq k_{r-1}}(A^0)
        \subset
        \cdots
        \subset
        \Lambda^{\geq k_{1}}(A^0)
        =
        \Lambda(A^0)
      \]
      is normal and
    \item we can use this to achieve the decomposition
      \[
        \Lambda(A^0)\cong\underset{k\in\cK}\bigltimes\Lambda^k(A^0)
      \]
      taken in an arbitrary order. In fact, one can also extend the algorithm
      from Remark~\ref{rem:algFactorization} to an arbitrary order of levels.
  \end{enumerate}
\end{rem}
The important statement, which we will use later, is then the following
Proposition. It is stated by Loday-Richaud in her Paper~\cite{Loday1994} as
Proposition I.5.5 or in the Paper~\cite[Thm.4.8]{Martinet1991} by Martinet and
Ramis.
\begin{prop}\label{prop:filtrationOfStokesGroup}
  The results from above can be restricted to the Stokes groups.
  Thus, for $\alpha\in\A$, one has
  \[
    \Sto_\alpha(A^0)\cong\underset{k\in\cK_\alpha}\bigltimes\Sto_\alpha^k(A^0)
  \]
  the semidirect product being taken in an arbitrary order\comm{~(we will only
  be interested in the ascending order)}.
  \begin{s-defn}
    We will denote the map which gives the factors of this factorization by
    \[
      i_\alpha:
      \Sto_\alpha(A^0)
      \overset{\cong}\longrightarrow
      \prod_{k\in\cK_\alpha}\Sto_\alpha^k(A^0)\,,
    \]
    where the factorization is taken in ascending order.
  \end{s-defn}
  \begin{s-rem}\label{rem:filtrationOfStokesMats}
    Write $\rho_{\alpha}^k:\Sto^k_\alpha(A^0)\to\SSto^k_\alpha(A^0)$ for
    the \rewrite{restriction} of the map $\rho_{\alpha}$
    (cf.\ Proposition~\ref{prop:representation}) to the level $k$.
    \rewrite{Then, one can} denote, by ab the induced decomposition also by
    \[
      i_\alpha:
      \SSto_\alpha(A^0)
      \overset{\cong}\longrightarrow
      \prod_{k\in\cK_\alpha}\SSto_\alpha^k(A^0)
    \]
    and the corresponding diagram
    \[ \begin{tikzcd}
        \Sto_\alpha(A^0)
        \rar{i_\alpha}
        \dar{\rho_{\alpha}}
        & \prod_{k\in\cK_\alpha}\Sto_\alpha^k(A^0)\,,
        \dar{\prod_{k\in\cK}\rho_{\alpha}^k}
      \\\SSto_\alpha(A^0)
        \rar{i_\alpha}
        & \prod_{k\in\cK_\alpha}\SSto_\alpha^k(A^0)\,,
    \end{tikzcd} \]
    commutes.
  \end{s-rem}
\end{prop}

%%% Local Variables:
%%% TeX-master: "Maximilian_Huber-Masters_Thesis-with_notes.tex"
%%% End:

    \fi

    \ifnum\myHideVariable=0
    %%%%%%%%%%%%%%%%%%%%%%%%%%%%%%%%%%%%%%%%%%%%%%%%%%%%%%%%%%%%%%%%%%%%%%%%%%%%%%%
\section{Stokes structures: using Stokes groups}\label{sec:mainThm2}
\marginnote{\cite{Loday1994},~\cite[Thm.4.3.11]{Loday2014}
  \\and~\cite{boalch,thboalch}
  \\and~\cite{babbitt1989local}
  \\and~\cite{BJL1979Birkhoff}
  \\and~\cite[Chapter 4]{Martinet1991}}

The goal in this section is to prove that there is a bijective and natural map
\[
  h:\prod_{\alpha\in\A}\Sto_\alpha(A^0)\longrightarrow\St(A^0)
\]
which endows $\St(A^0)$ with the structure of a unipotent Lie group.
And since $\Sto_\alpha(A^0)$ has $\SSto_\alpha(A^0)$ as a faithful
representation, we also get the isomorphism
$\prod_{\alpha\in\A}\SSto_\alpha(A^0)\cong\St(A^0)$ as a corollary.
\TODO[This goes back to \cite{BJL1979Birkhoff}?]

Let us recall, that $\St(A^0)$ is defined to be $H^1(S^1;\Lambda(A^0))$
(cf.\ Section~\ref{sec:mainThm1}).
The elements of $\prod_{\alpha\in\A}\Sto_\alpha(A^0)$ define in a canonical way
cocycles of the sheaf $\Lambda(A^0)$ (cf.\ Equation (\ref{eq:mapStoToCocy})),
called Stokes cocycles (cf.\ Definition~\ref{defn:stokesCocycle}).
In fact, will $h$ map such cocycles to the cohomology class, to which they
correspond.
Thus the statement, that $h$ is a bijection, is equivalent to the statement
that
\begin{einr}
  in each cohomology class of $\St(A^0)$ is a unique $1$-cocycle, which is a
  Stokes cocycle.
\end{einr}

%%%%%%%%%%%%%%%%%%%%%%%%%%%%%%%%%%%%%%%%%%%%%%%%%%%%%%%%%%%%%%%%%%%%%%%%%%%%%%%
\subsubsection{Cyclic coverings}
To formulate the theorem in the next section, we use the notion of cyclic
coverings and nerves of such coverings, which are defined as follows.

\begin{defn}
  \marginnote{\cite[Sec.II.1]{Loday1994} and \cite[Sec.II.3.1]{Loday1994}}
  Let $J$ be a finite set, identified to $\{1,\dots,p\}\subset\Z$.
  \begin{enumerate}
    \item A \emph{cyclic covering} of $S^1$ is a finite covering
      $\cU=\big(U_j=U(\theta_j,\epsilon_j)\big)_{j\in J}$ consisting of
      arcs, which satisfies that
      \begin{enumerate}
        \item $\tilde\theta_j \geq \tilde\theta_{j+1}$ for
          $j\in\{1,\dots,p-1\}$, i.e.\ the center points are ordered and
        \item $\tilde\theta_j+\frac{\epsilon_j}{2}\geq
          \tilde\theta_{j+1}+\frac{\epsilon_{j+1}}{2}$ for
          $j\in\{1,\dots,p-1\}$ and
          $\tilde\theta_p+\frac{\epsilon_p}{2}\geq
          \tilde\theta_{1}-2\pi+\frac{\epsilon_{1}}{2}$, i.e.\ the arcs are not
          encased by another arc,
      \end{enumerate}
      where the $\tilde\theta_j\in [0,2\pi[$ are determinations of the
      $\theta_j\in S^1$.
      \begin{comment}
        \begin{enumerate}
          \item the $\theta_j$ are in ascending order with respect to the
            clockwise orientation of $S^1$;
          \item the $U_j\cap U_{j+1}$ have only one connected component when
            $\#J>2$;
          \item the $U_j$ are not encased by another arc, this means that the
            open sets $U_j\backslash U_l$ are connected for all $j,l\in J$.
        \end{enumerate}
      \end{comment}
    \item The \emph{nerve} of a cyclic covering $\cU=\{U_j;j\in J\}$ is the
      family $\dot\cU=\{\dot U_j;j\in J\}$ defined by:
      \begin{itemize}
        \item $\dot U_j=U_j\cap U_{j+1}$ when $\#J>2$,
        \item $\dot U_1$ and $\dot U_2$ the connected components of
          $U_1\cap U_2$ when $\#J=2$.
      \end{itemize}
      \begin{s-rem}
        The nerve of the cyclic covering
        $\cU=\left(U(\theta_j,\epsilon_j)\right)_{j\in J}$ is explicitly given
        by
        \[
          \dot\cU=\left(
            \bigl(\theta_{j+1}-\frac{\epsilon_{j+1}}{2},
            \theta_{j}+\frac{\epsilon_{j}}{2}\bigr)
          \right)_{j\in J} \,.
        \]
      \end{s-rem}
  \end{enumerate}
\end{defn}
The cyclic coverings correspond one-to-one to nerves of cyclic coverings. If
one starts with a nerve $\{\dot U_j \mid j\in J\}$, one obtains a cyclic
covering as $\cU=\{U_j \mid j\in J\}$ where the arc $U_j$ is the connected
clockwise hull from $\dot U_{j-1}$ to $\dot U_j$.

\begin{defn}
  A covering $\cV$ is said to \emph{refine} a covering $\cU$ if, to each open
  set $V\in\cV$ there is at least one $U\in\cU$ with $V\subset U$.
\end{defn}
\begin{prop}
  \marginnote{\cite[Prop.II.1.3]{Loday1994}}
  The covering $\cV$ refines $\cU$ if and only if the corresponding nerves
  $\dot\cU=\{\dot U_j\}$ and $\dot\cV=\{\dot V_l\}$ satisfy
  \begin{einr}
    each $\dot U_j$ contains at least one $\dot V_l$.
  \end{einr}
\end{prop}

%%%%%%%%%%%%%%%%%%%%%%%%%%%%%%%%%%%%%%%%%%%%%%%%%%%%%%%%%%%%%%%%%%%%%%%%%%%%%%%
\subsection{The theorem}
\marginnote{\cite[868]{Loday1994}}
Let $\{\theta_j\mid j\in J\}\subset S^1$ be a finite set and
$\dot\phi=(\dot\phi_{\theta_j})_{j\in J}
\in\prod_{j\in J}\Lambda_{\theta_j}(A^0)$ be a finite family of germs.
Let $\dot\phi_j$ be the function representing the germ $\dot\phi_{\theta_j}$
on its (maximal) arc of definition $\Omega_j$ around $\theta_j$.
In the following way, one can associate a cohomology class in $\St(A^0)$ to
$\dot\phi$:
\begin{einr}
  for every cyclic covering $\cU=(U_j)_{j\in J}$ which satisfies
  $\dot U_j\subset\Omega_j$ for all $j\in J$, one can define the $1$-cocycle
  $(\dot\phi_{j|\dot U_j})_{j\in J}\in\Gamma(\dot\cU;\Lambda(A^0))$.
\end{einr}
To a different cyclic covering, satisfying the condition above, this
construction yields a cohomologous $1$-cocycle, thus the induced map
\begin{equation}\label{eq:mapStoToCocy}
  \prod_{j\in J}\Lambda_{\theta_j}(A^0)
  \longrightarrow
  H^1(S^1;\Lambda(A^0))=\St(A^0)
\end{equation}
is welldefined (cf.\ \cite[868]{Loday1994}).
\begin{defn}\label{defn:stokesCocycle}
  \marginnote{\cite[Def.II.1.8]{Loday1994}
    \\, \cite[4.3.10]{Loday2014}
    \\, \cite[Defn 6 on p 374]{Martinet1991}}
  Let $\nu=\#\A$ the number of all anti-Stokes directions and write
  $\A=\{\alpha_1,\dots,\alpha_\nu\}$.

  A \emph{Stokes cocycle} is a 1-cocycle $(\phi_j)_{j\in\{1,\dots,\nu\}}\in
  \prod_{j\in\{1,\dots,\nu\}}\Gamma(U_j;\Lambda(A^0))$ corresponding to some
  cyclic covering with nerve $\dot\cU=(\dot U_j)_{j\in\{1,\dots,\nu\}}$,
  which satisfies for every $j\in\{1,\dots,\nu\}$
  \begin{itemize}
    \item $\alpha_j\in\dot U_j$ and
    \item the germ $\phi_{\alpha_j}:=\phi_{j,\alpha_j}$ of $\phi_j$ at
      $\alpha_j$ is an element of $\Sto_{\alpha_j}(A^0)$.
  \end{itemize}
  \begin{s-rem}\label{rem:inclusionGermRemark}
    \PROBLEM[refactor!]
    The sections in $\Gamma(\dot U_j;\Lambda(A^0))$ are uniquely
    determined as the extension of the germ at $\alpha_j$, since the sheaf
    $\Lambda(A^0)$ defined via the system $[A^0,A^0]$
    (cf.\ Definition~\ref{defn:StokesSheaf}).
    We thus have an injective map
    \[
      \prod_{j\in\{1,\dots,\nu\}}\Gamma(\dot U_j;\Lambda(A^0))
      \hookrightarrow
      \prod_{j\in\{1,\dots,\nu\}}\Sto_{\alpha_j}(A^0) \,,
    \]
    which takes an Stokes cocycle and yields the corresponding Stokes germs.
    For a fine enough covering $\cU$, i.e.\ a covering $\cU$ with a nerve
    $\dot\cU$ which consists of small enough arcs satisfying the conditions
    above, is this map a bijection.

    We will use this fact implicitly and assume that the covering is always
    fine enough to call elements of $\prod_{\alpha\in\A}\Sto_\alpha(A^0)$
    Stokes cocycles.
  \end{s-rem}
\end{defn}
We can use the Equation (\ref{eq:mapStoToCocy}) to obtain for Stokes cocycles
a mapping
\[ \begin{tikzcd}
    h:\prod_{\alpha\in\A}\Sto_\alpha(A^0)
    \rar[hook]&
    \prod_{\alpha\in\A}\Lambda_\alpha(A^0)
    \rar{\text{(\ref{eq:mapStoToCocy})}}&
    \St(A^0),
\end{tikzcd} \]
which takes a Stokes cocycle to its corresponding cohomology class.
\begin{center}
  \begin{minipage}[t]{0.8\textwidth}
    \begin{tthm}\label{thm:mainThm2}
      The map
      \[ \begin{tikzcd}
          h:\prod_{\alpha\in\A}\Sto_\alpha(A^0) \rar& \St(A^0)
      \end{tikzcd} \]
      is a bijection and natural.
      \begin{s-rem}
        \marginnote{\cite[869]{Loday1994},\cite[Sec.III.3.3]{Loday1994}}
        Natural means that $h$ commutes to isomorphisms and constructions over
        systems or connections they represent.
      \end{s-rem}
    \end{tthm}
  \end{minipage}
\end{center}
From Theorem~\ref{thm:mainThm2} and Proposition~\ref{prop:representation} we
get the following corollary.
\begin{cor}\label{cor:isomToChochN}
  \PROBLEM[mentioned twice]
  Using the isomorphisms $\Sto_\theta(A^0)\cong\SSto_\theta(A^0)$ from
  Proposition~\ref{prop:representation} we obtain
  \[
    \St(A^0) \cong \prod_{\alpha\in\A}\SSto_\alpha(A^0)
  \]
  which endows $\St(A^0)$ with the structure of an unipotent Lie group with the
  finite complex dimension $N:=\dim_\C\St(A^0)$
  (cf.~\cite[Sec.III.1]{Loday1994}).
  This can be rewritten in the following way:
  \[
    N=\sum_{\alpha\in\A}\dim_\C\SSto_\alpha(A^0)
      =\sum_{\alpha\in\A}\sum_{q_j\myrel{\alpha}q_l}n_j\cdot n_l
      =\sum_{\substack{1\leq j,l\leq n\\j<l}}2\cdot\deg(q_j-q_l)\cdot
        n_j\cdot n_l \,.
  \]
  \begin{s-rem}
    This number $N$ is known to be the \emph{irregularity} of $[\End A^0]$.
  \end{s-rem}
\end{cor}
\begin{comment}
  \marginnote{\cite[880f]{Loday1994}}
  We have also two structures of a linear affine variety on the set $\St(A^0)$.
  \begin{rem}
    Let $\sto_\alpha(A^0)$ be the Lie algebra corresponding to
    $\Sto_\alpha(A^0)$. The exponential map\footnote{This is not the map $\exp$
    from Section~\ref{sec:mainThm1}.} induces an homomorphism
    $\exp:\sto_\alpha(A^0)\to\Sto_\alpha(A^0)$ and denote by $\ln=\exp^{-1}$
    the inverse map.
    \begin{enumerate}
      \item The \emph{tangent linear structure} is defined as \TODO
      \item Using the map
        \begin{align*}
          \sto_{\alpha}(A^0) &\overset{\id+\cdot}\longrightarrow
          \Sto_{\alpha}(A^0)
        \\\dot{f}_\alpha & \longmapsto \id+\dot{f}_\alpha \,.
        \end{align*}
    \end{enumerate}
  \end{rem}
\end{comment}

\begin{rem}
  To define the inverse map of $h$, one has to find in each cocycle in
  $\St(A^0)$ the Stokes cocycle. Loday-Richaud gives an algorithm in section
  II.3.4 of her paper~\cite{Loday1994}, which takes a cocycle over an arbitrary
  cyclic covering and outputs cohomologous Stokes cocycle and thus solves this
  problem.
\end{rem}

%%%%%%%%%%%%%%%%%%%%%%%%%%%%%%%%%%%%%%%%%%%%%%%%%%%%%%%%%%%%%%%%%%%%%%%%%%%%%%%
\subsection{Proof of Theorem~\ref{thm:mainThm2}}\label{sec:proofOfMatrixThm}
We will only look at the unramified case, for which we refer to
\cite[Sec.II.3]{Loday1994}.
The proof in the ramified case can be found in \cite[Sec.II.4]{Loday1994}.
We first have to introduce adequate coverings, which will be used in the proof.

%%%%%%%%%%%%%%%%%%%%%%%%%%%%%%%%%%%%%%%%%%%%%%%%%%%%%%%%%%%%%%%%%%%%%%%%%%%%%%%
\subsubsection{Adequate coverings}
\begin{defn}
  \marginnote{\cite[371]{Martinet1991} defines adapted coverings}
  Let $\star\in\{k,<k,\leq k,\dots\}$.
  A covering $\cU$ beyond which the inductive limit
  $\underset{\cU}{\underrightarrow{\lim}}H^1(\cU;\Lambda^\star(A^0))$ is
  stationary is said to be \emph{adequate} to describe
  $H^1(S^1;\Lambda^\star(A^0))$ or \emph{adequate} to $\Lambda^\star(A^0)$.
  \begin{comment}
    A covering $\cU$ is said to be \emph{adequate} to describe
    $H^1(S^1;\Lambda^\star(A^0))$ or \emph{adequate} to $\Lambda^\star(A^0)$ if
    for every element in
    $\underset{\cU}{\underrightarrow{\lim}}H^1(\cU;\Lambda^\star(A^0))$
    given by some covering $\cU'$ and an element of
    $\Gamma(\cU';\Lambda^\star(A^0))$
    there exists
    \begin{itemize}
      \item an element in $\Gamma(\cU;\Lambda^\star(A^0))$ and
      \item an common refinement of $\cU$ and $\cU'$
    \end{itemize}
    such that \PROBLEM[the elements are~?? on the refined covering.]
  \end{comment}

  In other words is a covering $\cU$ adequate, if and only if the quotient map
  \[
      \Gamma(\cU;\Lambda^\star(A^0))\longrightarrow H^1(S^1;\Lambda^\star(A^0))
  \]
  is surjective.\TODO[Proof? check??]
  \begin{comment}
    \cite[371]{Martinet1991} introduces the following definition
    \begin{s-defn}
      A covering $\cU$ is \emph{adapted} if every anti-Stokes direction is
      contained in exactly one element of the nerve $\dot\cU$.
    \end{s-defn}
  \end{comment}
\end{defn}
The following proposition is in Loday-Richaud's paper~\cite{Loday1994} given as
Proposition II.1.7. It contains a simple characterization, which will be used
to see, that our defined coverings are adequate.
\begin{prop}\label{prop:adeqCovCondition}
  \marginnote{\cite[Prop.II.1.7]{Loday1994}}
  Let $k\in\cK_\alpha$.
  \begin{s-defn}
    Let $\alpha\in\A^k$.
    An arc $U(\alpha,\frac{\pi}{k})$ is called a \emph{Stokes arc of level $k$
    at $\alpha$}.
  \end{s-defn}
  A cyclic covering $\cU=(U_j)_{j\in J}$, which satisfies
  \begin{einr}
    for every $\alpha\in\A^k$ contains the Stokes arc $U(\alpha,\frac{\pi}{k})$
    at least one arc $\dot U_j$ from the nerve $\dot\cU$ of $\cU$
  \end{einr}
  is adequate to $\Lambda^k(A^0)$.

  The covering $\cU$ is adequate to $\Lambda^{\leq k}(A^0)$ (resp.\
  $\Lambda^{<k}(A^0)$) if it is adequate to $\Lambda^{k'}(A^0)$ for every
  $k'\leq k$ (resp.\ $k'<k$).
\end{prop}

Let $k\in\cK$.
We want to define the three cyclic coverings $\cU^{k}$, $\cU^{\leq k}$ and
$\cU^{<k}$ which will be adequate to $\Lambda^k(A^0)$, $\Lambda^{\leq k}(A^0)$
and $\Lambda^{<k}(A^0)$. Furthermore will the coverings be comparable at the
different levels.

\begin{enumerate}
  \item The first covering $\cU^{k}=\{\dot U_\alpha^k\mid\alpha\in\A^k\}$
    is the cyclic covering with nerve
    \[
      \dot\cU^k:=
      \left\{\dot U_\alpha^k=U(\alpha,\frac{\pi}{k})\mid\alpha\in\A^k\right\}
    \]
    consisting of all Stokes arcs of level $k$ for anti-Stokes directions
    bearing the level $k$.
\end{enumerate}

\begin{rem}\label{rem:superSectors}
  Boalch introduces in his publications~\cite[19]{boalch}
  and~\cite[Def.1.23]{thboalch} the notion of \emph{supersectors}, they are in
  the case of a single level $k$, defined as follows:
  \begin{einr}
    write the anti-Stokes directions as $\A=\{\alpha_1,\dots,\alpha_\nu\}$
    arranged according to the clockwise ordering, then is the $i$-th
    supersector defined as the arc
    \[
      \hat\Sect^k_i:=
        \left(\alpha_i-\frac{\pi}{2k},\alpha_{i+1}+\frac{\pi}{2k}\right) \,.
    \]
  \end{einr}
  This yields a cyclic covering $(\hat\Sect^k_i)_{i\in\{1,\dots,\nu\}}$ whose
  nerve is exactly $\dot\cU^k$ defined above.
  \begin{comment}
    \begin{center}
      \begin{tikzpicture}[scale=3]
        \node[] (zero) at (0,0) {};
        \draw[blue] (zero) circle (1cm);

        \fill[fill=green!20!white] (0,0) -- ({cos( 15 )},{sin( 15 )}) arc
          (15:85:1) -- cycle;

        \fill[fill=red!60!black] (0,0) -- ({cos( 15 )*0.5},{sin( 15 )*0.5}) arc
          (15:45:0.5) -- cycle;
        \fill[fill=red!60!black] (0,0) -- ({cos( 55 )*0.5},{sin( 55 )*0.5}) arc
          (55:85:0.5) -- cycle;

        \fill[fill=green!20!white] (0,0) -- ({cos( 15 )*0.4},{sin( 15 )*0.4}) arc
          (15:85:0.4) -- cycle;

        \node[green!40!black] at (1,1) {$\widehat\Sect_i$};

        \node[] (lft) at ({cos( 85 )},{sin( 85 )}) {};
        \node[] (rgt) at ({cos( 15 )},{sin( 15 )}) {};

        \draw[->,red!40!black] ({cos( 45 )},{sin( 45 )})
          to [out=35, in=25] (rgt);

        \draw[->,red!40!black] ({cos( 55 )},{sin( 55 )})
          to [out=65, in=75] (lft);

        \draw[thick,red!40!black,path fading=west] (0,0) -- +({cos( 85 )},{sin( 85 )});
        \draw[thick,red!40!black,path fading=west] (0,0) -- +({cos( 15 )},{sin( 15 )});
        \fill[red!40!black] ({cos( 85 )},{sin( 85 )}) circle (1pt);
        \fill[red!40!black] ({cos( 15 )},{sin( 15 )}) circle (1pt);

        \foreach \w/\str in {10/,
                             20/,
                             45/$\alpha_{i+1}$,
                             55/$\alpha_{i}$}
        {\draw (0,0) -- +({cos( \w )},{sin( \w )}) node[right] {\str};
         \fill[blue!20!white] ({cos( \w )},{sin( \w )}) circle (.7pt);
         \foreach \sep in {60,120,180,240,300}
         {\draw (0,0) -- +({cos( \w + \sep )},{sin( \w + \sep )});
          \fill[blue!20!white]
            ({cos( \w + \sep )},{sin( \w + \sep )}) circle (.7pt);
         }
        };


        \fill (zero) circle (1pt);
      \end{tikzpicture}
    \end{center}
  \end{comment}
\end{rem}
If we extend to more then one level level, $\#\cK>1$, the set
$\bigcup_{k\in\cK}\left\{U(\alpha,\frac{\pi}{k})\mid\alpha\in\A^{k}\right\}$ is
no longer a nerve.
Hence we have to define the coverings $\cU^{\leq k}$ and $\cU^{<k}$ in a
different way.
Denote by
\[
  \left\{K_1<\cdots<K_s=k\right\}
  =\left\{\max\left(\cK_\alpha\cap[0,k]\right)\mid\alpha\in\A^{\leq k}\right\}
\]
the set of all \emph{$k$-maximum levels} for $\alpha\in\A^{\leq k}$.
\begin{enumerate}
  \setcounter{enumi}{1}
  \item The cyclic covering
    $\cU^{\leq k}=\left\{U_\alpha^{\leq k}\mid\alpha\in\A^{\leq k}\right\}$
    will be defined by induction.
    Let us assume that
    \begin{einr}
      the $\dot U_\alpha^{\leq k}$ are defined for all $\alpha\in\A^{\leq k}$
      with $k$-maximum level greater than $K_i$ such that their complete family
      is a nerve.
    \end{einr}
    Let
    \begin{itemize}
      \item $\alpha$ be a anti-Stokes direction with $k$-maximum level $K_i$
        and
      \item $\alpha^-$ (resp.\ $\alpha^+$) be the next anti-Stokes direction
        with $k$-maximum level greater then $K_i$ on the left (resp.\ on the
        right) and define $\dot U_{\alpha^-,\alpha^+}$ as the clockwise hull of
        the arcs $\dot U_{\alpha^-}^{\leq k}$ and $\dot U_{\alpha^+}^{\leq k}$
        already defined by induction.
        If there are no anti-Stokes directions with $k$-maximum level greater
        then $K_i$ we set $\dot U_{\alpha^-,\alpha^+}=S^1$.
    \end{itemize}
    We then set
    \[
      \dot U_\alpha^{\leq k}
        :=U\left(\alpha,\frac{\pi}{K_i}\right)\cap\dot U_{\alpha^-,\alpha^+}
    \]
    and the family of all $\dot U_\alpha^{\leq k}$ is a nerve.
    \begin{rem}
      If $\alpha$ has a $k$-maximum level equal to $k$ then is
      $\dot U_\alpha^{\leq k}$ equal to the Stokes arc
      $U\left(\alpha,\frac{\pi}{k}\right)=\dot U_\alpha^k$.
      \comm{\dots{}and then no $0$-cochain with level $k$ or $\geq k$ can
      exists on the covering $\cU^{\leq{k}}$.}
    \end{rem}
\end{enumerate}

\begin{enumerate}
  \setcounter{enumi}{2}
\item The last cyclic covering,
  $\cU^{<k}=\left\{U_{\alpha}^{<k}\mid\alpha\in\A^{<k}\right\}$, is defined as
  $\cU^{<k}:=\cU^{\leq k'}$ where $k':=\max\{k''\in\cK\mid k''<k\}$.
\end{enumerate}

\begin{figure} %{{{
  \centering

  \def\kOne{7}
  \def\kTwo{10}
  \begin{tikzpicture}[scale=5] %{{{

    \newcommand{\myDrawArcA}[4]{%{{{
      % Parameter: radius , center , width , color
      \pgfmathsetmacro\hwdth{#3 / 2}
      \draw[#4]
        ({cos( #2 )},{sin( #2 )})
        --
        ({cos( #2 ) * #1 },{sin( #2 ) * #1 });
      \draw[ultra thick,#4]
        ({cos(#2 - \hwdth) * #1},{sin(#2 - \hwdth) * #1})
        arc
        ({#2 - \hwdth}:{#2 + \hwdth}:#1);
      \draw[dotted,#4]
        ({cos(#2 - \hwdth) * #1},{sin(#2 - \hwdth) * #1})
        --
        ({cos(#2 - \hwdth)},{sin(#2 - \hwdth)});
      \draw[dotted,#4]
        ({cos(#2 + \hwdth) * #1},{sin(#2 + \hwdth) * #1})
        --
        ({cos(#2 + \hwdth)},{sin(#2 + \hwdth)});
      \filldraw[white] ({cos(#2)},{sin(#2)}) circle (0.5pt);
      \filldraw[red] ({cos(#2)},{sin(#2)}) circle (0.2pt);
    }%}}}
    \newcommand{\myDrawArcB}[5]{%{{{
      % Parameter: radius , center , start , stop , color
      \draw[#5]
        ({cos( #2 )},{sin( #2 )})
        --
        ({cos( #2 ) * #1 },{sin( #2 ) * #1 });
      \draw[ultra thick,#5]
        ({cos( #3 ) * #1},{sin( #3 ) * #1})
        arc
        ({ #3 }:{ #4 }:#1);
      \draw[dotted,#5]
        ({cos( #3 ) * #1},{sin( #3 ) * #1})
        --
        ({cos( #3 )},{sin( #3 )});
      \draw[dotted,#5]
        ({cos( #4 ) * #1},{sin( #4 ) * #1})
        --
        ({cos( #4 )},{sin( #4 )});
      \filldraw[white] ({cos(#2)},{sin(#2)}) circle (0.5pt);
      \filldraw[red] ({cos(#2)},{sin(#2)}) circle (0.2pt);
    }%}}}

    \node (zero) at (0,0) {};
    \draw (zero) circle (1cm);

    %%%%%%%%%%%%%%%%%%%%%%%%%%%%%%%%%%%%%%%%%%%%%%%%%%%%%%%%%%%%%%%%%%%%%%%%%
    %% Inner:
    \foreach \n/\mycol/\baseR in {\kOne/orange/0.8
                                 ,\kTwo/purple/0.9}
    {%{{{
      \pgfmathsetmacro\wdth{180/\n}
      \foreach \i in {1,2,...,\n}{%
        \pgfmathsetmacro\r{{\baseR + mod(\i+1,2)*0.05}}
        \pgfmathsetmacro\angl{\i* \wdth}
        \myDrawArcA{\r}{\angl}{\wdth}{\mycol}

        \pgfmathsetmacro\r{{\baseR + mod(\i + mod(\n+1,2),2)*0.05}}
        \pgfmathsetmacro\angl{\i* \wdth+180}
        \myDrawArcA{\r}{\angl}{\wdth}{\mycol}
      }
    }%}}}

    %%%%%%%%%%%%%%%%%%%%%%%%%%%%%%%%%%%%%%%%%%%%%%%%%%%%%%%%%%%%%%%%%%%%%%%%%
    %% Outer:
    \def\mycol{brown}
    \pgfmathsetmacro{\wdth}{180/\kTwo}
    \foreach \i in {1,2,...,\kTwo}{%
      \pgfmathsetmacro\r{{1.05 + mod(\i+1,2)*0.05}}
      \pgfmathsetmacro\angl{\i* \wdth}
      \myDrawArcA{\r}{\angl}{\wdth}{\mycol}

      \pgfmathsetmacro\r{{1.05 + mod(\i + mod(\kTwo+1,2),2)*0.05}}
      \pgfmathsetmacro\angl{\i* \wdth+180}
      \myDrawArcA{\r}{\angl}{\wdth}{\mycol}
    }
    \foreach \i in {1,2,...,\kOne}{%
      \foreach \j in {1,...,\kTwo}{%
        \pgfmathparse{\i/\kOne <= \j/\kTwo ? 0 : 1}
        \ifnum\pgfmathresult=0{%
          \pgfmathparse{\i/\kOne < \j/\kTwo ? 0 : 1}
          \ifnum\pgfmathresult=0{%
            \pgfmathsetmacro\r{{1.15 + mod(\i+1,2)*0.05}}
            \pgfmathsetmacro\center{\i/\kOne*180}
            \pgfmathsetmacro\start{(\i-0.5)*180/\kOne > (\j-1.5)*180/\kTwo
              ? (\i-0.5)*180/\kOne : (\j-1.5)*180/\kTwo}
            \pgfmathsetmacro\stop{(\i+0.5)*180/\kOne > (\j+0.5)*180/\kTwo
              ? (\j+0.5)*180/\kTwo : (\i+0.5)*180/\kOne}
            \myDrawArcB{\r}{\center}{\start}{\stop}{\mycol}

            \pgfmathsetmacro\r{{1.15 + mod(\i,2)*0.05}}
            \pgfmathsetmacro\center{\center+180}
            \pgfmathsetmacro\start{\start+180}
            \pgfmathsetmacro\stop{\stop+180}
            \myDrawArcB{\r}{\center}{\start}{\stop}{\mycol}
          }\fi
          \breakforeach
        }\fi
      }
    }

    \fill[white] (zero) circle (1.5pt);
    \fill (zero) circle (.7pt);
  \end{tikzpicture} %}}}
  \caption{The adequate coverings for an example with $\cK=\{\kOne,\kTwo\}$ and
    $\A=\left\{ \frac{j\cdot\pi}{k}\mid k\in\cK\text{, } j\in\N \right\}$.
    The anti-Stokes directions are marked by the \textcolor{red}{red} dots.
    The arcs of $\dot\cU^{\kOne}=\dot\cU^{\leq\kOne}$ are
    \textcolor{orange}{orange}, the arcs of $\dot\cU^{\kTwo}$ are
    \textcolor{purple}{purple} and the arcs of $\dot\cU^{\leq\kTwo}=\dot\cU$
    are \textcolor{brown}{brown}.
  }\label{fig:adequateCovering}
\end{figure}%}}}

\begin{rem}
  The coverings $\cU^{k}$, $\cU^{\leq k}$ and $\cU^{<k}$ depend only on
  $\cQ(A^0)$. Hence they depend only on the determining polynomials.
\end{rem}
It is obvious, that for every $k\in\cK$ the covering $\cU^{\leq k}$ refines
$\cU^{k}$ and $\cU^{<k}$.
Furthermore are the coverings defined, such that they satisfy the
condition in Proposition~\ref{prop:adeqCovCondition}.
Thus the first property in the following proposition is satisfied. The other
two can be found at \cite[Prop.II.3.1 (iv)]{Loday1994}.
\begin{prop}\label{prop:adequateProperties}
  \marginnote{\cite[Prop.II.3.1]{Loday1994}}
  \PROBLEM[Was bedeutet das?]
  Let $k\in\cK$, then
  \begin{enumerate}
    \item the coverings $\cU^{k}$, $\cU^{\leq k}$ and $\cU^{<k}$ are adequate
      to $\Lambda^k(A^0)$, $\Lambda^{\leq k}(A^0)$ and $\Lambda^{<k}(A^0)$, respectively;
    \item there exists no $0$-cochain in $\Lambda^k(A^0)$ on $\cU^k$;
    \item on $\cU^{\leq k}$ there is no $0$-cochain in $\Lambda^{\leq k}(A^0)$
      of level $k$, i.e.\ all $0$-cochains of $\Lambda^{\leq k}(A^0)$ belong to
      $\Lambda^{<k}(A^0)$.
  \end{enumerate}
\end{prop}
To have a shorter notation, we denote the product
$\prod_{\alpha\in\A^\star}\Gamma(\dot U_\alpha^\star;\Lambda^\star(A^0))$ by
$\Gamma(\dot U^\star;\Lambda^\star(A^0))$ for every
$\star\in\{k,<k,\leq k,\dots\}$.

%%%%%%%%%%%%%%%%%%%%%%%%%%%%%%%%%%%%%%%%%%%%%%%%%%%%%%%%%%%%%%%%%%%%%%%%%%%%%%%
\subsubsection{The case of a unique level}
\marginnote{\cite[II.3.2]{Loday1994}}
First we will proof Theorem~\ref{thm:mainThm2} in the case of a unique level.
This means that
\begin{itemize}
  \item either $\Lambda(A^0)$ has only one level $k$, thus
    \begin{itemize}
      \item $\Lambda(A^0)=\Lambda^k(A^0)$ and
      \item $\Sto_\theta(A^0)=\Sto_\theta^k(A^0)$ for every $\theta$,
    \end{itemize}
  \item or we restrict to a given level $k\in\cK$.
\end{itemize}
\begin{lem}
  Let $k\in\cK$.
  The morphism $h$ from Theorem~\ref{thm:mainThm2} is in the case of a unique
  level build as
  \[ \begin{tikzcd}[row sep=0cm]
    \underset{\alpha\in\A}\prod\Sto_\alpha^k(A^0)
    \rar{i^k}
    & \Gamma(\dot\cU^k;\Lambda^k(A^0))
    \rar{s^k}
    & H^1(S^1;\Lambda^k(A^0))
    \\
    \text{\rotatebox[origin=c]{-90}{$=$}}
    \tikzmark{e1}
    &&
    \tikzmark{e2}
    \text{\rotatebox[origin=c]{-90}{$=$}}
    \\
    \underset{\alpha\in\A}\prod\Sto_\alpha(A^0)
    \arrow{rr}{h}
    && H^1(S^1;\Lambda(A^0))
  \end{tikzcd} \]
  \begin{flushright}
    \tikzmarkc{n1}{blue} only in the single leveled case
    \begin{tikzpicture}[remember picture,overlay]
      \draw[->,blue!50!white,thick] (n1) to[out=180,in=0] (e1);
      \draw[->,blue!50!white,thick] (n1) to[out=180,in=180,distance=3cm] (e2);
    \end{tikzpicture}
  \end{flushright}
  from
  \begin{itemize}
    \item the canonical injective map
      \[ \begin{tikzcd}
        i^k:\underset{\alpha\in\A}\prod\Sto_\alpha^k(A^0) \rar&
        \Gamma(\dot\cU^k;\Lambda^k(A^0)) \,,
      \end{tikzcd} \]
      i.e.\ the map which is the canonical extension of germs to their natural
      arc of definition, and
    \item the quotient map
      \[ \begin{tikzcd}
        s^k:\Gamma(\dot\cU^k;\Lambda^k(A^0)) \rar& H^1(S^1;\Lambda^k(A^0))
      \end{tikzcd} \]
  \end{itemize}
  which are both isomorphisms.
\end{lem}
\begin{proof}
  \begin{enumerate}
  \item The map
    \[ \begin{tikzcd}
      i^k: \underset{\alpha\in\A}\prod\Sto_\alpha^k(A^0)
      \rar &
      \underset{\Gamma(\dot\cU^k;\Lambda^k(A^0))}{%
        \underset{\text{\rotatebox[origin=c]{-90}{$=$}}}{%
          \underbrace{%
            \prod_{\alpha\in\A^k}\Gamma(\dot\cU_\alpha^k;\Lambda^k(A^0))
          }
        }}
    \end{tikzcd} \]
    is welldefined, since the sections of $\Lambda^k(A^0)$ are solution of
    the system $[A^0,A^0]$ and it is very well known from the theory of
    differential equations\TODO[source?] that an element
    $f_\alpha\in\Gamma(\dot\cU_\alpha^k;\Lambda^k(A^0))$ is uniquely
    determined as the extension of its germ at some point $\alpha$.

    It is also a isomorphism groups (cf.~\cite{Loday1994}).
    \begin{comment}
      Problems:
      \begin{itemize}
      \item Show that a element of $\Lambda_\alpha^k(A^0)$ is extensionable
        to the arc $U_\alpha\in\dot\cU_\alpha^k$ if and only if it has
        maximal decay in direction $\alpha$.
      \item \PROBLEM[$\Sto_\alpha^k(A^0)\subsetneq\Lambda_\alpha^k(A^0)$]
        and thus
        \[
        i^k:
        \underset{\alpha\in\A}\prod\Sto_\alpha^k(A^0)
        \subsetneq
        \underset{\alpha\in\A}\prod\Lambda_\alpha^k(A^0)
        \to
        \prod_{\alpha\in\A^k}\Gamma(\dot\cU_\alpha^k;\Lambda^k(A^0))
        \]
      \end{itemize}
      Facts:
      \begin{itemize}
      \item Every arc $U\in\dot\cU_\alpha^k$ has the width $\frac{\pi}{k}$
        and is delimited by Stokes directions, i.e.\ is $U_\alpha^k$ the largest
        arc, to contain no \textbf{corresponding} Stokes ray
        \begin{einr}
          Then might \cite[Lemma 1]{BJL1979Birkhoff} on page 72-73 help? \textbf{NO!}
        \end{einr}
      \end{itemize}
      IDEAS:
      \begin{itemize}
      \item See \cite[375]{Martinet1991} and
        \textbf{\cite[Def.5 on 372]{Martinet1991}}
      \item \textbf{See \cite{babbitt1989local}} and
        \textbf{\textcolor{green!60!black}{\cite[72]{babbitt1989local}}}
      \item or maybe \cite{wasow2002asymptotic}
      \end{itemize}
    \end{comment}
  \item The second map
    \[ \begin{tikzcd}
      s^k:
      \overset{\Gamma(\dot\cU^k;\Lambda^k(A^0))}{%
        \overset{\text{\rotatebox[origin=c]{-90}{$=$}}}{%
          \overbrace{%
            \prod_{\alpha\in\A^k}\Gamma(\dot\cU_\alpha^k;\Lambda^k(A^0))}}}
      \rar &
      H^1(S^1;\Lambda^k(A^0))
    \end{tikzcd} \]
    is a bijection, \rewrite{since} from
    Proposition~\ref{prop:adequateProperties} we know that it is
    \begin{itemize}
    \item \textbf{surjective}, \rewrite{since} $\cU^k$ is adequate to
      $\Lambda^k(A^0)$ and
    \item \textbf{injective}, \rewrite{since} on $\cU^k$ there is no
      $0$-cochain in $\Lambda^k(A^0)$.
    \end{itemize}
  \end{enumerate}
  \PROBLEM[Show that this is the correct $h$]
  \PROBLEM[Naturality?]
\end{proof}

%%%%%%%%%%%%%%%%%%%%%%%%%%%%%%%%%%%%%%%%%%%%%%%%%%%%%%%%%%%%%%%%%%%%%%%%%%%%%%%
\subsubsection{The case of several levels}
\rewrite{In the proof of the case of several levels, we will still use}
Loday-Richaud's paper~\cite{Loday1994} as reference.
\begin{defn}\label{defn:firstSetOfInlusions}
  Here we want to define a \emph{product map of cocycles}
  $\mathfrak{S}^{\leq k}$.
  This map will be composed from the following injective maps:
  \begin{enumerate}
    \item The first map is defined as
      \begin{align*}
        \sigma^k:\Gamma(\dot\cU^k;\Lambda^k(A^0))
        &\longrightarrow \Gamma(\dot\cU^{\leq k};\Lambda^{\leq k}(A^0))
      \\\dot f=(\dot f_\alpha)_{\alpha\in\A^k}
        &\longmapsto (\dot G_\alpha)_{\alpha\in\A^{\leq k}}
      \end{align*}
      where
      \[
        \dot G_\alpha=\begin{cases}
          \dot f_\alpha \text{~restricted to } \dot U_\alpha^{\leq k}
          \text{~and seen as being in } \Lambda^{\leq k}(A^0)
          & \text{~when } \alpha\in\A^k
        \\\id \text{~(the identity) }
          & \text{~when } \alpha\notin\A^k
        \end{cases}
      \]
    \item and the second map
      \begin{align*}
        \sigma^{<k}:\Gamma(\dot\cU^{<k};\Lambda^{<k}(A^0))
        &\longrightarrow \Gamma(\dot\cU^{\leq k};\Lambda^{\leq k}(A^0))
      \\\dot f=(\dot f_\alpha)_{\alpha\in\A^{<k}}
        &\longmapsto (\dot F_\alpha)_{\alpha\in\A^{\leq k}}
      \end{align*}
      is defined, in a similar way, as
      \[
        \dot F_\alpha=\begin{cases}
          \dot f_\alpha \text{~restricted to } \dot U_\alpha^{\leq k}
          \text{~and seen as being in } \Lambda^{\leq k}(A^0)
          & \text{~when } \alpha\in\A^{<k}
        \\\id \text{~(the identity) }
          & \text{~when } \alpha\notin\A^{<k}
        \end{cases}
      \]
  \end{enumerate}
  Thus we can define
  \begin{align*}
    \mathfrak{S}^{\leq k}:
    \Gamma(\dot\cU^{<k};\Lambda^{<k}(A^0))
    \times
    \Gamma(\dot\cU^{k};\Lambda^{k}(A^0))
    &\longrightarrow
    \Gamma(\dot\cU^{\leq k};\Lambda^{\leq k}(A^0))
  \\(\dot f, \dot g)
    &\longmapsto
    (\dot F_\alpha\dot G_\alpha)_{\alpha\in\A^{\leq k}}
  \end{align*}
  where $(\dot F_\alpha)_{\alpha\in\A^{\leq k}}=\sigma^{<k}(\dot f)$ and
  $(\dot G_\alpha)_{\alpha\in\A^{\leq k}}=\sigma^k(\dot g)$ are defined as
  above.
  \begin{s-rem}
    This map $\mathfrak{S}^{\leq k}$ is injective, since injectivity for germs
    implies injectivity for sections.
  \end{s-rem}
\end{defn}

\begin{lem}
  \marginnote{\cite[Lem.II.3.3]{Loday1994}}
  Let $k\in\cK$.
  \begin{enumerate}
    \item If the cocycles $\mathfrak{S}^{\leq k}(\dot f,\dot g)$ and
      $\mathfrak{S}^{\leq k}(\dot f',\dot g')$ are cohomologous in
      $\Gamma(\dot\cU^{\leq k};\Lambda^{\leq k}(A^0))$
      then $\dot f$ and $\dot f'$ are cohomologous in
      $\Gamma(\dot\cU^{<k};\Lambda^{<k}(A^0))$.
    \item Any cocycle in $\Gamma(\dot\cU^{\leq k},\Lambda^{\leq k}(A^0))$ is
      cohomologous to a cocycle in the range of $\mathfrak{S}^{\leq k}$.
  \end{enumerate}
\end{lem}
\begin{proof}
  \begin{enumerate}
    \item Denote by $\alpha^+$ the nearest anti-Stokes direction in
      $\A^{\leq k}$ on the right\footnote{In clockwise direction.} of $\alpha$.
      The cocycles $\mathfrak{S}^{\leq k}(\dot f,\dot g)$ and
      $\mathfrak{S}^{\leq k}(\dot f',\dot g')$ are cohomologous if and only
      if there is a $0$-cochain $c=(c_\alpha)_{\alpha\in\A^{\leq k}}
      \in\Gamma(\cU^{\leq k},\Lambda^{\leq k}(A^0)$ such that
      \begin{equation}\label{eq:UniqueString:urdtindfgupndtcn}
        \dot F_\alpha\dot G_\alpha =
        c_\alpha^{-1}\dot F_\alpha'\dot G_\alpha'c_{\alpha^+}
      \end{equation}
      for every $\alpha\in\A$. From Proposition~\ref{prop:adequateProperties}
      follows, that $c$ is with values in $\Lambda^{<k}(A^0)$.
      The fact that $\Lambda^k(A^0)$ is normal in $\Lambda^{\leq k}(A^0)$ in
      Proposition~\ref{prop:PropertiesOfStokesSheafSplitting}, can be used to
      see that $ c_{\alpha^+}^{-1}G_\alpha'c_{\alpha^+}
      \in\Gamma(\cU^{\leq k};\Lambda^{k}(A^0))$.
      Thus, we rewrite the relation (\ref{eq:UniqueString:urdtindfgupndtcn}) to
      \[
        \dot F_\alpha\dot G_\alpha =
        (c_\alpha^{-1}\dot F_\alpha'c_{\alpha^+})
        (c_{\alpha^+}^{-1}G_\alpha'c_{\alpha^+})
        \,,\qquad\text{~for~} \alpha\in\A^{\leq k} \,.
      \]
      Since Corollary~\ref{cor:factorStokesGerms} tells us, that the
      factorization into the factors of the semidirect product are unique, we
      get for all $\alpha\in\A^k$
      \[
        \dot F_\alpha=c_\alpha^{-1}\dot F_\alpha'c_{\alpha^+}
        \qquad \text{~and~} \qquad
        \dot G_\alpha=c_{\alpha^+}^{-1}\dot G_\alpha'c_{\alpha^+}.
      \]
      The former relation implies that $(\dot F_\alpha)$ and $(\dot F_\alpha')$
      are cohomologous with values in $\Lambda^{<k}(A^0)$ on $\cU^{\leq k}$.
      Since $\cU^{<k}$ is already adequate to $\Lambda^{<k}(A^0)$, \rewrite{are
      $(\dot F_\alpha)$ and $(\dot F_\alpha')$} already on $\cU^{<k}$,
      i.e.\ in $\Gamma(\dot\cU^{<k};\Lambda^{<k}(A^0))$, cohomologous.
    \item The proof of part 2.\ (together with a proof of part 1.) can be
      found in Loday-Richaud's paper \cite[Proof of Lem.II.3.3]{Loday1994}.
  \end{enumerate}
\end{proof}
Let $k\in\cK$ and $k'=\max\{k'\in\cK\mid k'<k\}$. We then know by definition
that $\cU^{<k}=\cU^{\leq k'}$ as well as
$\Lambda^{< k}(A^0)=\Lambda^{\leq k'}(A^0)$ and thus
$\Gamma(\dot\cU^{< k};\Lambda^{< k}(A^0))=
\Gamma(\dot\cU^{\leq k'};\Lambda^{\leq k'}(A^0))$ and obtain the following
proposition.
\begin{prop}\label{prop:theMapTau}
  By applying $\mathfrak{S}^{\leq k}$ successively for different $k$'s
  in decending order, one obtains the \emph{product map of single leveled
  cocycles $\tau$} in the following way
  \[ \begin{tikzcd}[column sep=1.4cm,row sep=.9cm]
      \underset{\tikzmark{e1}}{\underbrace{%
        \Gamma(\dot\cU^{<k_r};\Lambda^{<k_r}(A^0))}}
      \times
      \Gamma(\dot\cU^{k_r};\Lambda^{k_r}(A^0))
      \rar{\mathfrak{S}^{\leq k_r}}&
      \overset{\Gamma(\dot\cU;\Lambda(A^0))}{%
        \overset{\text{\rotatebox[origin=c]{-90}{$=$}}}{%
          \Gamma(\dot\cU^{\leq k_r};\Lambda^{\leq k_r}(A^0))
      }}
      \\\underset{\tikzmark{e2}}{\underbrace{%i
        \Gamma(\dot\cU^{<k_{r-1}};\Lambda^{<k_{r-1}}(A^0))}}
      \times
      \Gamma(\dot\cU^{k_{r-1}};\Lambda^{k_{r-1}}(A^0))
      \rar{\mathfrak{S}^{\leq k_{r-1}}}&
      \overset{\tikzmark{n1}}{%
        \Gamma(\dot\cU^{\leq k_{r-1}};\Lambda^{\leq k_{r-1}}(A^0))}
      \\\hspace{6cm}\cdots \rar{\mathfrak{S}^{\leq k_{r-2}}}&
      \overset{\tikzmark{n2}}{%
        \Gamma(\dot\cU^{\leq k_{r-2}};\Lambda^{\leq k_{r-2}}(A^0))}
      \\\underset{\tikzmark{eEND}}{\underbrace{%
        \Gamma(\dot\cU^{<k_3};\Lambda^{<k_3}(A^0))}}
      \times
      \Gamma(\dot\cU^{k_3};\Lambda^{k_3}(A^0))
      \rar{\mathfrak{S}^{\leq k_3}}&
      \cdots\hspace{3cm}
      \\
      \Gamma(\dot\cU^{k_1};\Lambda^{k_1}(A^0))
      \times
      \Gamma(\dot\cU^{k_2};\Lambda^{k_2}(A^0))
      \rar{\mathfrak{S}^{\leq k_2}}&
      \overset{\tikzmark{nEND}}{%
        \Gamma(\dot\cU^{\leq k_2};\Lambda^{\leq k_2}(A^0))}
  \end{tikzcd} \]
  \begin{tikzpicture}[remember picture,overlay]
    \draw[<-] (n1) to[out=90,in=270] node[midway,fill=white]{$=$}
      ([yshift=.3em]e1);
    \draw[<-] (n2) to[out=90,in=270] node[midway,fill=white]{$=$}
      ([yshift=.3em]e2);
    \draw[<-] (nEND) to[out=90,in=270] node[midway,fill=white]{$=$}
      ([yshift=.3em]eEND);
  \end{tikzpicture}
  which can be written in the following compact form
  \begin{align*}
    \tau:\prod_{k\in\cK}\Gamma(\dot \cU^k;\Lambda^k(A^0))
    &\longrightarrow
    \Gamma(\dot\cU;\Lambda(A^0))
  \\(\dot f^k)_{k\in\cK}
    &\longmapsto
    \prod_{k\in\cK}\tau^k(\dot f^k)
  \end{align*}
  where the product is following an ascending order of levels and the maps
  $\tau_k$ are defined as
  \[ \begin{tikzcd}[row sep=0cm]
    \tau^k:\Gamma(\dot\cU^k;\Lambda^k(A^0))
    \arrow{r}{\sigma^k}&
    \Gamma(\dot\cU^{\leq k};\Lambda^{\leq k}(A^0))
    \rar &
    \Gamma(\dot\cU;\Lambda(A^0))
  \\~~~(\dot f_\alpha)_{\alpha\in\A^k}
    \arrow[|->]{rr}
    &&
    (\dot G_\alpha)_{\alpha\in\A}
  \end{tikzcd} \]
  with
  \[
    \dot G_\alpha=\begin{cases}
      \dot f_\alpha \text{~restricted to } \dot U_\alpha
      \text{~and seen as being in } \Lambda(A^0)
      & \text{~when } \alpha\in\A^k
    \\\id \text{~(the identity on $\dot U_\alpha$) }
      & \text{~when } \alpha\notin\A^k
    \end{cases}
  \]
  The defined map $\tau$ is clearly injective and it can be extended to an
  arbitrary order of levels (cf.\ Remark~\cite[Rem.II.3.5]{Loday1994}).
\end{prop}
\begin{comment}
  \begin{defn}\label{defn:theMapTau}
    Define the injective map
    \begin{align*}
      \tau^k:\Gamma(\dot\cU^k;\Lambda^k(A^0))
      &\longrightarrow \Gamma(\dot\cU;\Lambda(A^0))
    \\\dot f=(\dot f_\alpha)_{\alpha\in\A^k}
      &\longmapsto
      (\dot F_\alpha)_{\alpha\in\A}
    \end{align*}
    where
    \[
      \dot F_\alpha=\begin{cases}
        \dot f_\alpha \text{~restricted to } \dot U_\alpha
        \text{~and seen as being in } \Lambda(A^0)
        & \text{~when } \alpha\in\A^k
      \\\id \text{~(the identity on $\dot U_\alpha$) }
        & \text{~when } \alpha\notin\A^k
      \end{cases}
    \]
    \marginnote{\cite[Prop.II.3.4]{Loday1994}}
    The \emph{product map of single-leveled cocycles} is then defined as
    \begin{align*}
      \tau:\prod_{k\in\cK}\Gamma(\dot \cU^k;\Lambda^k(A^0))
      &\longrightarrow
      \Gamma(\dot\cU;\Lambda(A^0))
    \\(\dot f^k)_{k\in\cK}
      &\longmapsto
      \prod_{k\in\cK}\tau^k(\dot f^k)
    \end{align*}
    following an ascending order of levels.
    \begin{s-rem}
      The map $\tau$
      \begin{enumerate}
        \item is injective since it is composed from $\sigma^k$ and the clearly
          injetive mapping
          \begin{align*}
            \Gamma(\dot\cU^{\leq k};\Lambda^{\leq k}(A^0))
            &\longrightarrow
            \Gamma(\dot\cU;\Lambda(A^0))
          \\(\dot F_\alpha)_{\alpha\in\A^{\leq k}}
            &\longmapsto
            (\dot F'_\alpha)_{\alpha\in\A}
          \end{align*}
          where
          \[
            \dot F'_\alpha=\begin{cases}
              \dot F_\alpha \text{~restricted to } \dot U_\alpha
              \text{~and seen as being in } \Lambda(A^0)
              & \text{~when } \alpha\in\A^{\leq k}
              \\\id \text{~(the identity on $\dot U_\alpha$) }
              & \text{~when } \alpha\notin\A^{\leq k}
            \end{cases}
          \]
          and
        \item it can be extended to an arbitrary order of levels
          (cf.\ Remark~\cite[Rem.II.3.5]{Loday1994}).
      \end{enumerate}
    \end{s-rem}
  \end{defn}
\end{comment}
\begin{cor}
  The product map of single-leveled cocycles $\tau$ induces on the cohomology
  a bijective and natural map
  \[ \begin{tikzcd}
    \cT:
    \underset{\underset{k\in\cK}\prod H^1(S^1;\Lambda^k(A^0))}{%
      \underset{\text{\rotatebox[origin=c]{-90}{$\cong$}}}{%
        \prod_{k\in\cK}\Gamma(\dot\cU^k;\Lambda^k(A^0))}}
    \rar&
    \underset{H^1(S^1;\Lambda(A^0))}{%
      \underset{\text{\rotatebox[origin=c]{-90}{$\cong$}}}{%
        H^1(\cU;\Lambda(A^0))}}\,.
  \end{tikzcd} \]
\end{cor}

%%%%%%%%%%%%%%%%%%%%%%%%%%%%%%%%%%%%%%%%%%%%%%%%%%%%%%%%%%%%%%%%%%%%%%%%%%%%%%%
\paragraph{Composing functions to obtain $h$}
We have the ingredients to define the function $h$ from
Theorem~\ref{thm:mainThm2} by composition of already bijective maps.
\begin{proof}[Proof of Theorem~\ref{thm:mainThm2}]
  Let $i_\alpha:\Sto_\alpha(A^0)\to\prod_{k\in\cK}\Sto_\alpha^k(A^0)$ be the
  map which corresponds to the filtration from
  Proposition~\ref{prop:filtrationOfStokesGroup} and
  denote the composition
  \[ \begin{tikzcd}[column sep=1.8cm,row sep=0]
      \displaystyle \prod_{\alpha\in\A}\Sto_\alpha(A^0)
      \rar{\prod_{\alpha\in\A}i_\alpha}
      \arrow[ddrr, out=270,in=200,"\mathfrak{T}"]
      &
      \displaystyle \prod_{\alpha\in\A}\prod_{k\in\cK}\Sto_\alpha^k(A^0)
    \\&\text{\rotatebox[origin=c]{-90}{$\equiv$}}
    \\&\displaystyle \prod_{k\in\cK}\prod_{\alpha\in\A}\Sto_\alpha^k(A^0)
      \rar{\prod_{k\in\cK}i^k}&
      \displaystyle \prod_{k\in\cK}\Gamma(\dot\cU^k;\Lambda^k(A^0))
  \end{tikzcd} \]
  by $\mathfrak{T}$. The bijection $h$ is then obtained as
  \[
    \cT\circ\mathfrak{T}: \prod_{\alpha\in\A}\Sto_\alpha(A^0)
    \longrightarrow H^1(\cU;\Lambda(A^0)) \,.
  \]
  \PROBLEM[naturality (is obvious?)]
  \PROBLEM[Show that this is the correct $h$]
\end{proof}

%%%%%%%%%%%%%%%%%%%%%%%%%%%%%%%%%%%%%%%%%%%%%%%%%%%%%%%%%%%%%%%%%%%%%%%%%%%%%%%
\subsection{Some exemplary calculations}\label{sec:WhichInformationIsNeeded}
Here we want to discuss which information is required to describe the Stokes
cocycle corresponding to a multileveled system in more depth.
We will look at a sigle-leveled system corresponding to a normal form
$A^0\in \GL_3(\C(\!\{t\}\!))$ with exactly two levels and will apply the
techniques developed in the previous sections in an rather explicit way.

Let $A^0$ be a normal form with dimension $n=3$ and two levels
$\cK=\{k_1<k_2\}$ which satisfies that there is at least one anti-Stokes
direction $\theta$ which is beared by both levels.
Let $q_j(t^{-1})$ be the determining polynomials and let $k_{jl}$ be the
degrees of $(q_j-q_l)(t^{-1})$.
Up to permutation, we know that in our case are the leading terms of
$(q_1-q_2)(t^{-1})$ and $(q_1-q_3)(t^{-1})$ equal and thus
\begin{itemize}
  \item up to permutation is $k_2=k_{1,2}=k_{1,3}$ and $k_1=k_{2,3}$, i.e.\ the
    larger degree appears twice, and
    \begin{comment}
      let $q_1,q_2,q_3$ be polynomials, such that
      \[
        \deg(q_1-q_2) =: k_2 > k_1 := \deg(q_2-q_3)
      \]
      then is the degree of $q_1-q_3$ given by
      \begin{itemize}
        \item[\textbf{case 1}] $\deg(q_1)<k_2$: then is $\deg(q_2)=k_2$ and
          thus $\deg(q_3)=k_2$.
        \item[\textbf{case 2}] $\deg(q_2)<k_2$: then is $\deg(q_1)=k_2$ and
          $\deg(q_3)\leq\deg(q_2)$.
        \item[\textbf{case 3}] $\deg(q_1)=\deg(q_2)=k_2$: thus follows that the
          leading term of $q_1$ and $q_2$ are different.
          \begin{itemize}
            \item[\textbf{subcase 3.a}] $\deg(q_3)<k_2$: everything is clear
            \item[\textbf{subcase 3.b}] $\deg(q_3)=k_2$: here has the leading
              term of $q_3$ be equal to them from $q_1$ to satisfy that
              $k_2>\deg(q_2-q_3)$.
          \end{itemize}
      \end{itemize}
      From this follows that $\deg(q_1-q_3)=k_2$.
    \end{comment}
  % \item $q_1 \underset{\alpha,\max}{\prec} q_2$
  %   \Leftrightarrow{}
  %   $q_1 \underset{\alpha,\max}{\prec} q_3$
  %   \rewrite{respectively}
  %   $q_2 \underset{\alpha,\max}{\prec} q_1$
  %   \Leftrightarrow{}
  %   $q_3 \underset{\alpha,\max}{\prec} q_1$
  \item $q_0\myrel{\alpha}q_2$ (resp.~$q_2\myrel{\alpha}q_1$) if and only if
    $q_1\myrel{\alpha}q_3$ (resp.~$q_3\myrel{\alpha}q_1$) and thus do they
    determine the same anti-Stokes directions.
\end{itemize}
The set of all anti-Stokes directions is then given as
\[
  \A=\left\{\theta+\frac{\pi}{k}\cdot j\mid k\in\cK\text{, }j\in\N\right\} \,.
\]
Denote by $\cY_0(t)$ a normal solution of $[A^0]$.

Let us start by looking at a single germ in depth.
The Proposition~\ref{prop:representation} states that every Stokes germ
$\phi_\alpha$ can be written as its matrix representation conjugated by the
normal solution, i.e.\ as $\phi_\alpha=\cY_{0}C_{\phi_\alpha}\cY_{0}^{-1}
=\rho_{\alpha}^{-1}(C_{\phi_\alpha})$.

Look at an example in which we will demonstrate, from which relations on the
determining polynomials which restriction on the form of the Stokes matrices
arise.
\begin{exmp}
  Let $\alpha\in\A$ be an anti-Stokes direction.
  From the definition of $\SSto_\alpha(A^0)$ (cf.\
  Definition~\ref{defn:groupOfFaithfullReps}) we know that, if one has
  $q_1\myrel{\alpha}q_2$, the Stokes matrix has the form
  \[
    \begin{pmatrix}
      1 & \text{\boldmath$c_1$} & \star
    \\\text{\boldmath$0$} & 1 & \star
    \\\star & \star & 1
    \end{pmatrix}
  \]
  where $c_j\in\C$ and $\star\in\C$.

  We have seen that $q_1\myrel{\alpha}q_2$ \Rightarrow{}
  $q_1\myrel{\alpha}q_3$ thus the representation has the
  \rewrite{form}
  \[
    \begin{pmatrix}
      1 & c_1 & \text{\boldmath$c_2$}
    \\0 & 1 & \star
    \\\text{\boldmath$0$} & \star & 1
    \end{pmatrix}
  \]
  and if we also know that neither $q_2\myrel{\alpha}q_3$ nor
  $q_3\myrel{\alpha}q_2$ it has the \rewrite{form}
  \[
    \begin{pmatrix}
      1 & c_1 & c_2
    \\0 & 1 & \text{\boldmath$0$}
    \\0 & \text{\boldmath$0$} & 1
    \end{pmatrix}\,.
  \]
  \begin{comment}
    We also know that every matrix of this \rewrite{form} is a representation
    to some Stokes germ.
    Thus we have an isomorphism
    \begin{align*}
      \vartheta_\alpha:\C^2 &\longrightarrow \SSto_\alpha(A^0)
    \\(c_1,c_2)&\longmapsto
      \begin{pmatrix}
        1 & c_1 & c_2
      \\0 & 1 & 0
      \\0 & 0 & 1
      \end{pmatrix}
    \end{align*}
  \end{comment}
\end{exmp}
In fact, the following $9$ cases of Stokes matrices can arise:
\begin{center}
  \def\arraystretch{1.3}
  \setlength\tabcolsep{4mm}
  \begin{tabular}{r|c|c|c}
    & $q_2\myrel{\alpha}q_3$ & $q_3\myrel{\alpha}q_2$ & else
    \tabularnewline
    \hline
    $\substack{q_1\myrel{\alpha}q_2\\\text{and} \\q_1\myrel{\alpha}q_3}$
    & $\begin{pmatrix} 1 & c_2 & c_3 \\0 & 1 & c_1 \\0 & 0 & 1 \end{pmatrix}$
   \cellcolor{blue!15}
    & $\begin{pmatrix} 1 & c_2 & c_3 \\0 & 1 & 0 \\0 & c_1 & 1 \end{pmatrix}$
   \cellcolor{blue!15}
    & $\begin{pmatrix} 1 & c_2 & c_3 \\0 & 1 & 0 \\0 & 0 & 1 \end{pmatrix}$
   \cellcolor{green!15}
    \tabularnewline
    \hline
    $\substack{q_2\myrel{\alpha}q_1\\\text{and} \\q_3\myrel{\alpha}q_1}$
    & $\begin{pmatrix} 1 & 0 & 0 \\c_2' & 1 & c_1 \\c_3 & 0 & 1 \end{pmatrix}$
   \cellcolor{blue!15}
    & $\begin{pmatrix} 1 & 0 & 0 \\c_2 & 1 & 0 \\c_3' & c_1 & 1 \end{pmatrix}$
   \cellcolor{blue!15}
    & $\begin{pmatrix} 1 & 0 & 0 \\c_2 & 1 & 0 \\c_3 & 0 & 1 \end{pmatrix}$
   \cellcolor{green!15}
    \tabularnewline
    \hline
    else
    & $\begin{pmatrix} 1 & 0 & 0 \\0 & 1 & c_1 \\0 & 0 & 1 \end{pmatrix}$
   \cellcolor{purple!15}
    & $\begin{pmatrix} 1 & 0 & 0 \\0 & 1 & 0 \\0 & c_1 & 1 \end{pmatrix}$
   \cellcolor{purple!15}
    & $\begin{pmatrix} 1 & 0 & 0 \\0 & 1 & 0 \\0 & 0 & 1 \end{pmatrix}$
  \end{tabular}
\end{center}
In the \textcolor{blue!75!black}{blue} cases we have $\cK_\alpha=\cK$ and
$\C^3\overset{\vartheta_\alpha}{\underset{\cong}{\longrightarrow}}\SSto_\alpha(A^0)$.
In the \textcolor{green!50!black}{green} cases $\cK_\alpha=\{k_2\}$ and
$\C^2\overset{\vartheta_\alpha}{\underset{\cong}{\longrightarrow}}\SSto_\alpha(A^0)$
as well as in the \textcolor{purple!75!black}{purple} cases
$\cK_\alpha=\{k_1\}$ and
$\C^1\overset{\vartheta_\alpha}{\underset{\cong}{\longrightarrow}}\SSto_\alpha(A^0)$.
\comm{Thus, for every $\alpha\in\A$, we have an isomorphism
$\rho_{\alpha}^{-1}\circ\vartheta_\alpha$.}
We will replace $c_2'$ by $c_2+c_1c_3$ and $c_3'$ by $c_1c_2+c_3$ to be
consistent with the decomposition in the next part
(cf.\ Example~\ref{exmp:decompositionHere}).
\begin{cor}
  The morphism $\prod_{\alpha\in\A}\vartheta_\alpha$ is an isomorphism of
  pointed sets, which maps the element only containing zeros to
  \[
    (\id,\id,\dots,\id)\in\prod_{\alpha\in\A}\SSto_\alpha(A^0),
  \]
  which gets by $\left(\prod_{\alpha\in\A}\right)^{-1}\circ h$ mapped to the
  trivial cohomology class in $\St(A^0)$.
\end{cor}

%%%%%%%%%%%%%%%%%%%%%%%%%%%%%%%%%%%%%%%%%%%%%%%%%%%%%%%%%%%%%%%%%%%%%%%%%%%%%%%
In proposition~\ref{prop:filtrationOfStokesGroup} and especially
Remark~\ref{rem:filtrationOfStokesMats} we have defined a decomposition of the
Stokes group $\Sto_\alpha(A^0)$ in subgroups generated by $k$-germs for
$k\in\cK$.
In our case, we have at most two nontrivial factors. Especially is this
decomposition given by
\[
  \phi_\alpha=\phi_\alpha^{k_1} \phi_\alpha^{k_2}
  \overset{i_\alpha}\longmapsto
    \left(\phi_\alpha^{k_1},\phi_\alpha^{k_2}\right)
      \in\Sto_\alpha^{k_1}(A^0)\times\Sto_\alpha^{k_2}(A^0) \,,
\]
and $i_\alpha$ is the map, wich gives the factors of this factorization in
ascending order.
This decomposition, of a germ $\phi_\alpha$, is trivial if
$\#\cK(\phi_\alpha)\leq1$, thus the interesting cases are the
\textcolor{blue!75!black}{blue} cases.

\begin{exmp}\label{exmp:decompositionHere}
  Look at the example
  \[
    \vartheta_\alpha(c_1,c_2,c_3)=
    \cY_{0}
    \begin{pmatrix} 1 & 0 & 0 \\c_2 & 1 & 0 \\c_1c_2+c_3 & c_1 & 1 \end{pmatrix}
    \cY_{0}^{-1}
    =\phi_\alpha
    \,.
  \]
  According to Remark~\ref{rem:algFactorization} the factor
  $\phi_\alpha^{k_1}\in\Sto_\alpha^{k_1}(A^0)$, is given by
  \[
    \phi_\alpha^{k_1}=
    \cY_{0}
    \begin{pmatrix}
      1 & 0 & 0
    \\\text{\boldmath $0$} & 1 & 0
    \\\text{\boldmath $0$} & c_1 & 1
    \end{pmatrix}
    \cY_{0}^{-1}
    \,.
  \]
  The other factor $\phi_\alpha^{k_2}$ is then obtained as
  \begin{align*}
    \phi_\alpha^{k_2}&=
    \left(\phi_\alpha^{k_1}\right)^{-1}
    \phi_\alpha^{k_2}
  \\&=\cY_{0}
    \begin{pmatrix}
      1     & 0    & 0
    \\0     & 1    & 0
    \\0     & -c_1 & 1
    \end{pmatrix}
    \underset{=\id}{\underbrace{%
        \cY_{0}^{-1}
        \cY_{0}
    }}
    \begin{pmatrix} 1 & 0 & 0 \\c_2 & 1 & 0 \\c_1c_2+c_3 & c_1 & 1 \end{pmatrix}
    \cY_{0}^{-1}
  \\&=\cY_{0}
    \begin{pmatrix}
      1     & 0 & 0
    \\c_2     & 1          & 0
    \\c_3     & 0          & 1
    \end{pmatrix}
    \cY_{0}^{-1}
    \,.
  \end{align*}
\end{exmp}
The four nontrivial decomposition in our situation, are given by:
\begin{enumerate}
  \item $\begin{pmatrix} 1 & 0 & 0 \\0 & 1 & c_1 \\0 & 0 & 1 \end{pmatrix}
  \cdot\begin{pmatrix} 1 & c_2 & c_3 \\0 & 1 & 0 \\0 & 0 & 1 \end{pmatrix}=
  \begin{pmatrix} 1 & c_2 & c_3 \\0 & 1 & c_1 \\0 & 0 & 1 \end{pmatrix}$
  \item $\begin{pmatrix} 1 & 0 & 0 \\0 & 1 & 0 \\0 & c_1 & 1 \end{pmatrix}
  \cdot\begin{pmatrix} 1 & c_2 & c_3 \\0 & 1 & 0 \\0 & 0 & 1 \end{pmatrix}=
  \begin{pmatrix} 1 & c_2 & c_3 \\0 & 1 & 0 \\0 & c_1 & 1 \end{pmatrix}$
  \item $\begin{pmatrix} 1 & 0 & 0 \\0 & 1 & c_1 \\0 & 0 & 1 \end{pmatrix}
  \cdot\begin{pmatrix} 1 & 0 & 0 \\c_2 & 1 & 0 \\c_3 & 0 & 1 \end{pmatrix}=
  \begin{pmatrix} 1 & 0 & 0 \\c_2+c_1c_3 & 1 & c_1 \\c_3 & 0 & 1 \end{pmatrix}$
  \item $\begin{pmatrix} 1 & 0 & 0 \\0 & 1 & 0 \\0 & c_1 & 1 \end{pmatrix}
  \cdot\begin{pmatrix} 1 & 0 & 0 \\c_2 & 1 & 0 \\c_3 & 0 & 1 \end{pmatrix}=
  \begin{pmatrix} 1 & 0 & 0 \\c_2 & 1 & 0 \\c_1c_2+c_3 & c_1 & 1 \end{pmatrix}$
\end{enumerate}

%%%%%%%%%%%%%%%%%%%%%%%%%%%%%%%%%%%%%%%%%%%%%%%%%%%%%%%%%%%%%%%%%%%%%%%%%%%%%%%
\subsubsection{Explicit example}
\def\kOne{1}
\def\kTwo{3}
\def\zkOnepzKtwo{14} % 2\cdot(\kOne+2\cdot\kTwo
\def\zkOne{2} % 2*\kOne
\def\zkTwo{6} % 2*\kTwo

Even more explicit, we can fix the levels $k_1=\kOne$ and $k_2=\kTwo$ together
with $\theta=0$.
Assume without any restriction that $q_1\myrel{\theta}q_2$ and
$q_1\myrel{\theta}q_3$ as well as $q_2\myrel{\theta}q_3$.
Other choices would result \rewrite{in reordering of the tuples below.}
\comm{Let the matrix $L$ be given as $L=\diag(l_1,l_2,l_3)\in\Gl_n(\C)$.}

The classification space is in this case isomorphic to
$\C^{2\cdot(\kOne+2\cdot\kTwo)}=\C^{\zkOnepzKtwo}$.
The element
\[
  ({}^1c_1,{}^2c_1,
  {}^1c_2,{}^1c_3,{}^2c_2,{}^2c_3,\dots,{}^{\zkTwo}c_2,{}^{\zkTwo}c_3)
  \in\C^{\zkOnepzKtwo}
\]
gets, via the isomorphism $\prod_{\alpha\in\A}j_\alpha$, mapped to
\begin{align*}
  &\left(
  \left(
    \begin{pmatrix} 1 & 0 & 0 \\0 & 1 & {}^1c_1 \\0 & 0 & 1 \end{pmatrix},
    \begin{pmatrix} 1 & 0 & 0 \\0 & 1 & 0 \\0 & {}^2c_1 & 1 \end{pmatrix}
  \right),
  \right.
\\&\qquad\left(
  \left.
    \begin{pmatrix} 1 & {}^1c_2 & {}^1c_3 \\0 & 1 & 0 \\0 & 0 & 1 \end{pmatrix},
    \begin{pmatrix} 1 & 0 & 0 \\{}^2c_2 & 1 & 0 \\{}^2c_3 & 0 & 1 \end{pmatrix},
    \dots,
    \begin{pmatrix} 1 & 0 & 0 \\{}^{\zkTwo}c_2 & 1 & 0 \\{}^{\zkTwo}c_3 & 0 & 1 \end{pmatrix}
  \right)
  \right)
\end{align*}
in $\prod_{\alpha\in\A^{\kOne}}\SSto_{\alpha}^{\kOne}(A^0) \times
\prod_{\alpha\in\A^{\kTwo}}\SSto_{\alpha}^{\kTwo}(A^0)$ and thus the element
\begin{align*}
  &\left(
  \left(
    \begin{pmatrix} 1 & 0 & 0 \\0 & 1 & {}^1c_1 \\0 & 0 & 1 \end{pmatrix},
    \id,\id,
    \begin{pmatrix} 1 & 0 & 0 \\0 & 1 & 0 \\0 & {}^2c_1 & 1 \end{pmatrix},
    \id,\id
  \right),
  \right.
\\&\qquad\left(
  \left.
    \begin{pmatrix} 1 & {}^1c_2 & {}^1c_3 \\0 & 1 & 0 \\0 & 0 & 1 \end{pmatrix},
    \begin{pmatrix} 1 & 0 & 0 \\{}^2c_2 & 1 & 0 \\{}^2c_3 & 0 & 1 \end{pmatrix},
    \dots,
    \begin{pmatrix} 1 & 0 & 0 \\{}^{\zkTwo}c_2 & 1 & 0 \\{}^{\zkTwo}c_3 & 0 & 1 \end{pmatrix}
  \right)
  \right)
\end{align*}
in
$\prod_{\alpha\in\A}\SSto_{\alpha}^{\kOne}(A^0) \times
\prod_{\alpha\in\A}\SSto_{\alpha}^{\kTwo}(A^0)$.
Using the morphism $\prod_{\alpha\in\A}i_\alpha^{-1}$ we get a complete set of
Stokes matrices as
\begin{align*}
  &\left(
    \begin{pmatrix} 1 & {}^1c_2 & {}^1c_3 \\0 & 1 & {}^1c_1 \\0 & 0 & 1 \end{pmatrix},
    \begin{pmatrix} 1 & 0 & 0 \\{}^2c_2 & 1 & 0 \\{}^2c_3 & 0 & 1 \end{pmatrix},
    \begin{pmatrix} 1 & {}^3c_2 & {}^3c_3 \\0 & 1 & 0 \\0 & 0 & 1 \end{pmatrix},
  \right.
\\&\qquad
  \left.
    \begin{pmatrix} 1 & 0 & 0 \\{}^4c_2 & 1 & 0 \\{}^2c_1{}^4c_2+{}^4c_3 & {}^2c_1 & 1 \end{pmatrix},
    \begin{pmatrix} 1 & {}^5c_2 & {}^5c_3 \\0 & 1 & 0 \\0 & 0 & 1 \end{pmatrix},
    \begin{pmatrix} 1 & 0 & 0 \\{}^{\zkTwo}c_2 & 1 & 0 \\{}^{\zkTwo}c_3 & 0 & 1 \end{pmatrix}
  \right)
  \in
  \prod_{\alpha\in\A}\SSto_{\alpha}(A^0) \,.
\end{align*}
Applying the isomorphism $\prod_{\alpha\in\A}\rho_\alpha^{-1}$, i.e.\
conjugation by the fundamental solution $\cY_0(t)=t^Le^{Q(t^{-1})}$
(cf.\ Proposition~\ref{prop:representation}), yields then the corresponding
Stokes cocycle in $\prod_{\alpha\in\A}\Sto_{\alpha}(A^0)$ and thus an element
in $\St(A^0)$.
\begin{comment}
  This element is explicitly given as
  \begin{align*}
    &\left(
      \begin{pmatrix}
          1 & {}^1c_2 t^{l_2-l_1}e^{(q_2-q_1)(t^{-1})}
            & {}^1c_3 t^{l_3-l_1}e^{(q_3-q_1)(t^{-1})}
        \\0 & 1 & {}^1c_1 t^{l_3-l_2}e^{(q_3-q_2)(t^{-1})}
        \\0 & 0 & 1
      \end{pmatrix},
      \cY_0^{-1}
      \begin{pmatrix}
          1 & 0 & 0
        \\{}^2c_2 & 1 & 0
        \\{}^2c_3 & 0 & 1
      \end{pmatrix}
      \cY_0
      ,
    \right.
  \\&\qquad
      \cY_0^{-1}
      \begin{pmatrix}
          1 & {}^3c_2 & {}^3c_3
        \\0 & 1 & 0
        \\0 & 0 & 1
      \end{pmatrix}
      \cY_0
      ,
      \cY_0^{-1}
      \begin{pmatrix}
          1 & 0 & 0
        \\{}^4c_2 & 1 & 0
        \\{}^2c_1{}^4c_2+{}^4c_3 & {}^2c_1 & 1
      \end{pmatrix}
      \cY_0
      ,
  \\&\qquad\qquad
    \left.
      \cY_0^{-1}
      \begin{pmatrix}
          1 & {}^5c_2 & {}^5c_3
        \\0 & 1 & 0
        \\0 & 0 & 1
      \end{pmatrix}
      \cY_0
      ,
      \cY_0^{-1}
      \begin{pmatrix}
          1 & 0 & 0
        \\{}^{\zkTwo}c_2 & 1 & 0
        \\{}^{\zkTwo}c_3 & 0 & 1
      \end{pmatrix}
      \cY_0
    \right)
    \in
    \prod_{\alpha\in\A}\Sto_{\alpha}(A^0) \,.
  \end{align*}
\end{comment}

    \fi

    \ifnum\myHideVariable=0
    %%%%%%%%%%%%%%%%%%%%%%%%%%%%%%%%%%%%%%%%%%%%%%%%%%%%%%%%%%%%%%%%%%%%%%%%%%%%%%%
\section{Further improvements}\label{sec:furtherImprovements}
The set $\prod_{\alpha\in\A}\Sto_\alpha(A^0)$ and thus
$\prod_{\alpha\in\A}\SSto_\alpha(A^0)$ has some bad properties, when small
deformations are applied to $[A^0]$, since under arbitrary small changes, one
Stokes ray can split into two.

In the prove of Theorem~\ref{thm:mainThm2} we have seen, \rewrite{that not only
$\cH(A^0)\cong\prod_{\alpha\in\A}\Sto_\alpha(A^0)$ but} also that
\begin{align*}
  \cH(A^0)&\cong\prod_{k\in\cK}\prod_{\alpha\in\A}\Sto_\alpha^k(A^0)
  \\&\cong\prod_{k\in\cK}\prod_{\alpha\in\A^k}\Sto_\alpha^k(A^0)
  &\text{(since $\SSto_{\alpha}^{k}(A^0)=\{\id\}$ when $k\notin\cK_\alpha$.)}
  \\&\cong\prod_{k\in\cK}\prod_{\alpha\in\A^k}\SSto_\alpha^k(A^0) \,.
\end{align*}
This representation can be used to achieve further improvements, since one is
able to multiply some succeeding Stokes matrices (resp. Stokes germs) of the
same level without loss of information.
This is precisely stated in Corollary~\ref{cor:composeLevelwise} below.

\TODO[move?]
Boalch, who looks in his publications~\cite{boalch,thboalch} only at the
single-leveled case, uses this extensively to obtain better Stokes matrices,
which are stable under small deformations. Our Stokes matrices are in his
publications called \emph{Stokes factors}.

Let us fix a level $k\in\cK$. We want to rewrite the product
$\prod_{\alpha\in\A^k}\SSto_\alpha^k(A^0)$ by collecting the information of a
subsets of these Stokes matrices into one product of these matrices.

\begin{defn}
  A subset of $\A^k$ consisting of $\frac{\#\A^k}{2k}$ consecutive anti-Stokes
  directions of level $k$ will be called a \emph{half-period (of level $k$)}.
\end{defn}
From the definition of anti-Stokes directions of level $k$ it is clear that
every arc of width $\frac{\pi}{k}$, which has no anti-Stokes direction of level
$k$ on its border, contains $\frac{\#\A^k}{2k}$ anti-Stokes directions of level
$k$ since
\begin{einr}
  for every sector $I=U(\theta,\frac{\pi}{k})$ of width $\frac{\pi}{k}$ and
  center $\theta$ which satisfies $\theta\pm\frac{\pi}{2k}\notin\A^k$
  there is to any pair $(q_j,q_l)$ in $\cQ(A^0)$ which has level $k_{jl}=k$
  exactly one direction $\alpha\in \A^k\cap I$ which satisfies that
  $\Im(a_{jl}e^{-ik_{jl\alpha}})=0$.
  At such a direction $\alpha$ corresponding to the pair $(q_j,q_l)$ is then
  either $q_j\myrel{\alpha}q_l$ or $q_l\myrel{\alpha}q_j$ satisfied.
\end{einr}
This implies that for every $\theta$, which satisfies
$\theta\pm\frac{\pi}{2k}\notin\A^k$, is $\A^k\cap U(\theta,\frac{\pi}{k})$
a half-period and by
$\dot\bigcup_{1\leq m\leq 2k}\A^k\cap U(\theta+\frac{m\pi}{k},\frac{\pi}{k})
=\A^k$ is a \rewrite{decomposition} in
distinct half-periods of $\A^k$ given (cf.\ Figure~\ref{fig:halfPeriod}).

\begin{figure}[h!] %{{{
  \begin{center}
    \begin{tikzpicture}[scale=4]
      \pgfmathsetmacro{\k}{1}
      \pgfmathsetmacro{\abs}{0.4}

      \draw[thick, dotted, green!50!black] (0.9,0) -- (0.9,{\abs + 0.1})
        node [above,font=\tiny,] {$\theta$};

      \foreach \x in {-2,-1,...,3}{
        \foreach \argA/\absA in {-0.5/0.4
                                ,0.3/0.35
                                ,0.1/0.27
                                ,0.65/0.23}{
          \begin{scope}
            \clip (-0.65,{\abs+0.1}) rectangle (2.63,{-\abs-0.1});
            \pgfmathsetmacro{\s}{{\argA + \x / \k * 2 - 1/2/\k}};

            \draw[blue!40!white] (\s,0)
                sin ({\s + 1/2/\k},{-\absA})
                cos ({\s + 2/2/\k},0)
                sin ({\s + 3/2/\k},\absA)
                cos ({\s + 4/2/\k},0);
            \draw[dotted] ({\s + 1/2/\k},{-\absA}) -- ({\s + 1/2/\k},-0.5);
            \fill[white] ({\s + 1/2/\k},{-\absA}) circle (1pt);
            \fill[red] ({\s + 1/2/\k},{-\absA}) circle (.4pt);

            \draw[blue!40!black] (\s,0)
                sin ({\s + 1/2/\k},\absA)
                cos ({\s + 2/2/\k},0)
                sin ({\s + 3/2/\k},{-\absA})
                cos ({\s + 4/2/\k},0);
            \draw[dotted] ({\s + 3/2/\k},{-\absA}) -- ({\s + 3/2/\k},-0.5);
            \fill[white] ({\s + 3/2/\k},{-\absA}) circle (1pt);
            \fill[red] ({\s + 3/2/\k},{-\absA}) circle (.4pt);
          \end{scope}
          \begin{scope}[thick]
            \clip (0.4,{\abs+0.1}) rectangle (1.4,{-\abs-0.1});
            \pgfmathsetmacro{\s}{{\argA + \x / \k * 2 - 1/2/\k}};

            \draw[blue!40!white] (\s,0)
                sin ({\s + 1/2/\k},{-\absA})
                cos ({\s + 2/2/\k},0)
                sin ({\s + 3/2/\k},\absA)
                cos ({\s + 4/2/\k},0);
            \fill[white] ({\s + 1/2/\k},{-\absA}) circle (1pt);
            % \draw[dotted] ({\s + 1/2/\k},{-\absA}) -- ({\s + 1/2/\k},0.5);
            \fill[red] ({\s + 1/2/\k},{-\absA}) circle (.4pt);

            \draw[blue!40!black] (\s,0)
                sin ({\s + 1/2/\k},\absA)
                cos ({\s + 2/2/\k},0)
                sin ({\s + 3/2/\k},{-\absA})
                cos ({\s + 4/2/\k},0);
            \fill[white] ({\s + 3/2/\k},{-\absA}) circle (1pt);
            % \draw[dotted] ({\s + 3/2/\k},{-\absA}) -- ({\s + 3/2/\k},0.5);
            \fill[red] ({\s + 3/2/\k},{-\absA}) circle (.4pt);
          \end{scope}
        }
      }
      \node[font=\tiny] at (-0.5,{-\abs - 0.15}) {$\alpha_{\nu-1}$};
      \node[font=\tiny] at (-0.35,{-\abs - 0.15}) {$\alpha_{\nu}$};
      \node[font=\tiny] at (0.1,{-\abs - 0.15}) {$\alpha_1$};
      \node[font=\tiny] at (0.3,{-\abs - 0.15}) {$\alpha_2$};
      \node[font=\tiny] at (0.5,{-\abs - 0.15}) {$\alpha_3$};
      \node[font=\tiny] at (0.65,{-\abs - 0.15}) {$\alpha_4$};
      \node[font=\tiny] at (1.1,{-\abs - 0.15}) {$\alpha_5$};
      \node[font=\tiny] at (1.3,{-\abs - 0.15}) {$\alpha_6$};
      \node[font=\tiny] at (1.5,{-\abs - 0.15}) {$\alpha_7$};
      \node[font=\tiny] at (1.65,{-\abs - 0.15}) {$\alpha_8$};
      \node[font=\tiny] at (2.1,{-\abs - 0.15}) {$\alpha_9$};
      \node[font=\tiny] at (2.3,{-\abs - 0.15}) {$\alpha_{10}$};
      \node[font=\tiny] at (2.5,{-\abs - 0.15}) {$\alpha_{11}$};

      \draw[-latex'] (-0.7,0) -- (2.7,0) node [right] {$S^1$};
      \draw[-latex'] ({0},{-\abs-0.1}) -- ({0},{\abs + 0.2});

      \begin{scope}[dashed]
        \clip (-0.65,{\abs+0.1}) rectangle (2.6,{-\abs-0.4});
        \foreach \x in {-2,-1,1,2} {
          % \draw [purple]
          %   ({0.4 + \x},{-\abs - 0.2}) -- ({0.4 + \x},0);
          \draw [purple]
            ({1.4 + \x},{-\abs - 0.2}) -- ({1.4 + \x},0);
          \draw [purple
                ,decorate
                ,decoration={brace,mirror,amplitude=10pt}
                ,xshift=0pt
                ,yshift=0pt]
            ({0.4 + \x},{-\abs - 0.2}) -- ({1.4 + \x},{-\abs - 0.2});
        }
      \end{scope}
      \draw [purple,dashed]
        (1.4,{-\abs - 0.2}) -- (1.4,0);
      \draw [purple, thick
            ,decorate
            ,decoration={brace,mirror,amplitude=10pt}
            ,xshift=0pt
            ,yshift=0pt]
        (0.4,{-\abs - 0.2}) -- (1.4,{-\abs - 0.2})
        node [midway,yshift=-7pt] {$\tikzmark{g0}$};
      \draw [purple,dashed,
            ,decorate
            ,decoration={brace,mirror,amplitude=10pt}
            ,xshift=0pt
            ,yshift=0pt]
        (1.4,{-\abs - 0.2}) -- (2.4,{-\abs - 0.2})
        node [midway,yshift=-7pt] {$\tikzmark{g1}$};

      % \draw[dotted] (-0.6,\abs)node[left,font=\tiny] {$|a|$} -- (2.6,\abs);
      % \draw[dotted] (-0.6,{-\abs})node[left,font=\tiny] {$-|a|$} -- (2.6,{-\abs});

      % \draw[dotted] ({-0.5},{-\abs}) -- ({-0.5},{\abs + 0.1})
      %   node [above,font=\tiny,] {-0.5 \pi};
      % \foreach \x in {0.5,1,...,2.5}{%
      %   \draw[dotted] ({\x},{-\abs}) -- ({\x},{\abs + 0.1})
      %     node [above,font=\tiny,] {\x \pi};
      % }
    \end{tikzpicture}
  \end{center}
  \begin{flushright}
    \tikzmarkc{n1}{purple}
    $\A^k\cap U(\theta+\frac{\pi}{k},\frac{\pi}{k})=\{\alpha_7,\alpha_8,\alpha_9,\alpha_{10}\}$
  \end{flushright}
  \begin{flushright}
    \tikzmarkc{n0}{purple}
    $\A^k\cap U(\theta,\frac{\pi}{k})=\{\alpha_3,\alpha_4,\alpha_5,\alpha_6\}$
  \end{flushright}
  \begin{tikzpicture}[remember picture,overlay]
    \draw[->,purple!50!white,thick,dashed] (n1) to[out=150,in=270] (g1);
    \draw[->,purple!50!white,thick] (n0) to[out=180,in=270] (g0);
  \end{tikzpicture}
  \caption{Assume that $\cQ(A^0)$ contains exactly eight ordered pairs of degree
    $k$. Choose four of them, which are up to order different. In this picture
    are the real parts of the leading terms corresponding to the pairs in
    \textcolor{blue!40!black}{blue} plotted. In \textcolor{blue!60!white}{light
      blue} are the real parts corresponding to the flipped pairs drawn. The
    points, where the real part is negative and the imaginary part vanishes,
    i.e.\ where the pairs determine an anti-Stokes direction, are marked by the
    \textcolor{red}{red} dots.}\label{fig:halfPeriod}
\end{figure} %}}}
\begin{prop}\label{prop:propAboutCommonStructure}
  Let $I=U(\theta,\frac{\pi}{k})$ be an arc width $\frac{\pi}{k}$ such that
  $\theta\pm\frac{\pi}{2k}\notin\A^k$.
  For every $\alpha$ in the half-period $\A^k\cap I$ we can write the Stokes
  matrix as
  \[
    K^\alpha=(K^\alpha_{jl})_{j,l\in\{1,\dots,s\}}\in\SSto_\alpha^k(A^0)
  \]
  where the $K^\alpha_{jl}\in\C^{n_j\times n_l}$ are blocks corresponding to
  the structure of $Q$ (cf.\ Definition~\ref{defn:groupOfFaithfullReps}).
  We then know for $j\neq l$ that, if $K^\alpha_{jl}\neq0$ then
  \begin{enumerate}
  \item is $K^\alpha_{lj}=0$ and
  \item for every $\alpha'\in(\A^k\cap I)\backslash\{\alpha\}$ is
    $K^{\alpha'}_{jl}=0$ as well as $K^{\alpha'}_{lj}=0$.
  \end{enumerate}
\end{prop}
\begin{proof}
  The first statement was already discussed
  page~\pageref{page:firstStatementOnTheStructure}.
  To see the second statement, look at a pair $(q_j,q_l)\in\cQ(A^0)$ together
  with an anti-Stokes direction $\alpha$ such that $q_j\myrel{\alpha}q_l$. We
  then know (cf.\ Lemma~\ref{lem:rotationalSym}) that the next directions
  $\alpha'$ satisfying $q_j\myrel{\alpha'}q_l$ are $\alpha+\frac{2\pi}{k}$ (on
  the left) and $\alpha-\frac{2\pi}{k}$ (on the right).
  The next directions satisfying the flipped relation $q_l\myrel{\alpha'}q_j$
  are given by $\alpha+\frac{\pi}{k}$ (on the left) and
  $\alpha-\frac{\pi}{k}$ (on the right).

  Thus there is no other anti-Stokes direction but $\alpha$ for the pairs
  $(q_j,q_l)$ and $(q_l,q_j)$ in every arc of width $\frac{\pi}{k}$ containing
  $\alpha$.
\end{proof}
\begin{cor}\label{cor:someCorolaryToLabel}
  In the situation of Proposition~\ref{prop:propAboutCommonStructure}
  there exists a common (block) permutation matrix $P\in\GL_n(\C)$ given by
  $(P)_{jl}=\bdelta_{\pi(j)l}$ where
  \begin{itemize}
  \item $\bdelta_{jl}$ is the block version of Kronecker's delta corresponding
    to the structure of $Q$ which was introduced in
    Definition~\ref{defn:groupOfFaithfullReps} and
  \item $\pi$ is the permutation of $\{1,\dots,s\}$ corresponding to
    $q_j\underset{\theta}{\prec}q_l$ \Rightarrow{} $\pi(j)<\pi(l)$
    \marginnote{\cite[18]{boalch}}
  \end{itemize}
  such that every matrix $P^{-1}K^{\alpha}P$ is upper triangular and the
  observation from Proposition~\ref{prop:propAboutCommonStructure} is still
  satisfied.
  \begin{s-rem}
    After moving the sector $U(\theta,\frac{\pi}{k})$ to
    $U(\theta,\frac{\pi}{k})+\frac{\pi}{k}
    =U(\theta+\frac{\pi}{k},\frac{\pi}{k})$ we obtain also a corresponding
    permutation $\pi'$ which is inverse to $\pi$.
    The corresponding permutation matrix $P'$ is then given by $P'=P^{-1}$.

    Such that the permutation $P$ transforms Stokes matrices on the sector
    $I+\frac{\pi}{k}$ to lower triangular matrices
  \end{s-rem}
  In the single leveled case of this phenomenon is discussed on page 18 of
  Boalch's paper~\cite{boalch}.
\end{cor}

In the paper \cite{BJL1979Birkhoff} from Balser, Jurkat and Lutz is the
following Lemma stated as Lemma 2 on page 75.  It is the fundamental statement,
which allows us to collect the information of multiple Stokes matrices into one
matrix, i.e.\ it states that we will not loose information.
\begin{lem}\label{lem:UniqueDecompositionWotBlocks}
  Let $T\subset\{1,\dots,n\}\times\{1,\dots,n\}$ be a position set, which
  satisfies the \emph{completeness property}:
  \begin{einr}
    if $(j,k)$ and $(k,l)\in T$ then is also $(j,l)\in T$.
  \end{einr}
  Choose a indexing $i:\{1,2,\dots,\mu \} \overset{\cong}{\to}T$ of the position
  set and denote by $\delta_{jl}\in\C$ the \rewrite{ordinary} Kronecker's delta,
  then there exists for every matrix
  \[
    K\in \left\{K=(K_{jl})_{j,l\in\{1,\dots,n\}}\in\GL_n(\C)\mid
      K_{jl}=\delta_{jl} \text{~unless~} (j,l)\in T \right\}
  \]
  unique scalars $t_1,\dots,t_\mu\in\C$ such that
  \[
    K=(\id + t_1E_{i(1)})\cdots(\id + t_{\mu}E_{i(\mu)}) \,.
  \]
\end{lem}
The completeness property is reasonable, since for example the matrices
corresponding to the not complete set $\{(2,3),(3,2)\}$ are not stable under the
product:
\begin{align*}
  \begin{pmatrix}
    1 & a & 0
  \\0 & 1 & 0
  \\0 & 0 & 1
  \end{pmatrix}
  \begin{pmatrix}
    1 & 0 & 0
  \\0 & 1 & b
  \\0 & 0 & 1
  \end{pmatrix}
  =
  \begin{pmatrix}
    1 & a & \textcolor{red!60!black}{ab}
  \\0 & 1 & b
  \\0 & 0 & 1
  \end{pmatrix}
  \,.
\end{align*}
The completeness property looks similar to the transitivity condition on
relations, this yields the following useful proposition.
\begin{cor}
  Every position set defined via an relation $\overset{R}{<}$ as
  $\{(j,l)\mid j\overset{R}{<} l\}$ is complete, if and only if the relation
  $\overset{R}{<}$ is transitive.
\end{cor}
\begin{prop}\label{prop:capAndCupOfRels}
  It is easy to see that
  \begin{enumerate}
  \item $\{(j,l) \mid q_j \myrel{\theta} q_l \} =
    \bigcap_{\theta'\in \A^k\cap U(\theta,\frac{\pi}{k})}
    \{(j,l) \mid q_j \underset{\theta'}{\prec} q_l \}$
    and
  \item
    $\{(j,l) \mid q_j \underset{\theta}{\prec} q_l \} =
    \dot\bigcup_{\theta'\in \A^k\cap U(\theta,\frac{\pi}{k})}
    \{(j,l) \mid q_j \myrel{\theta'} q_l \}$
  \end{enumerate}
  and all the position sets above are complete.
\end{prop}

Another useful fact is given by the following lemma.
\begin{lem}
  Every position set, corresponding to some block, is complete.
\end{lem}
\begin{proof}
  Such a position set, corresponding to some block, is given by
  \[
    T=\{(j,l)\mid j_1\leq j\leq j_2, l_1\leq l\leq l_2\}
  \]
  for some $j_1,j_2,l_2$ and $l_2$ in $\{1,\dots,n\}$.

  Let $(j,k)$, $(k,l)\in T$ be two positions.
  It is obvious that $(j,l)\in T$, since $j_1\leq j\leq j_2$ and
  $l_1\leq l\leq l_2$ are satisfied.
\end{proof}
\begin{cor}
  We can write the Lemma~\ref{lem:UniqueDecompositionWotBlocks} in block form,
  corresponding to the structure of $Q$.
  Let $s$ be the number of diagonal blocks of $Q$ and
  $T\subset\{1,\dots,s\}\times\{1,\dots,s\}$ be a position set.
  \begin{s-defn}
    Define the group of matrices, corresponding to a complete position set, by
    \[
      S_T:=\left\{K=(K_{jl})_{j,l\in\{1,\dots,n\}}\in\GL_n(\C)\mid
      K_{jl}=\bdelta_{jl} \text{~unless~} (j,l)\in T \right\} \,,
    \]
    where $\bdelta_{jl}$ is the block version of Kronecker's delta,
    corresponding to the structure of $Q$.
  \end{s-defn}
  Choose an enumeration
  \[
    i:\{1,2,\dots,\mu\}\overset{\cong}{\to}T
  \]
  of $T$.
  From the Lemma~\ref{lem:UniqueDecompositionWotBlocks} then follows that for
  every $K\in S_T$ there is a unique decomposition
  \[
    K=K_1\cdot K_2\cdots K_\mu \,,
  \]
  where $K_j\in S_{\{i(j)\}}$.
\end{cor}
From the previous corollary we deduce the following corollary.
\begin{cor}\label{cor:composeLevelwise}
  Let $T_1,\dots,T_r\subset\{1,\dots,s\}\times\{1,\dots,s\}$ be a distinct
  position sets, such that
  \begin{einr}
    $T:=T_1\dot\cup\dots\dot\cup T_r$ \rewrite{as well as} every $T_m$ satisfy
    the completeness property.
  \end{einr}
  Then is by
  \begin{align*}
    S_{T_1}\times\cdots\times S_{T_m} &\longrightarrow S_T
  \\(K_1,\dots,K_m) &\longmapsto K_1\cdots K_m
  \end{align*}
  an isomorphism defined.
\end{cor}
We can apply the previous corollary in our situation, i.e.\ let $\theta\in S^1$
be a fixed direction which satisfies $\theta\pm\frac{\pi}{2k}\notin\A^k$.
Let us first define the corresponding groups.
\begin{defn}
  Let us define the group
  \begin{align*}
    \hat\SSto_{\theta}^k(A^0)&:=
      \left\{K=(K_{jl})_{j,l\in\{1,\dots,s\}}\in\GL_n(\C)\mid
        K_{jl}=\delta_{jl} \text{~unless~}
        q_j \underset{\theta}{\prec} q_l \text{, } k_{jl}=k\right\}
    \\&=S_{\{(j,l) \mid q_j \underset{\theta}{\prec} q_l \}}
  \end{align*}
  which contains the matrices with the collected data. They are representation
  of the corresponding elements in
  \[
    \hat\Sto_{\theta}^k(A^0):=
    \{\cY_0^{-1}K\cY_0\mid K\in\hat\SSto_{\theta}^k(A^0) \} \,.
  \]
  Denote the isomorphism
  \begin{align*}
    \hat\Sto_\theta^k(A^0) &\overset{\cong}{\longrightarrow}
                             \hat\SSto_\theta^k(A^0)
    \\ \phi &\longmapsto \cY_0\phi\cY_0^{-1}
  \end{align*}
  by $\hat\rho_\theta^k$.
\end{defn}
The Proposition~\ref{prop:capAndCupOfRels} implies that $\Sto_{\theta}^k(A^0)\cong
\bigcap_{\theta'\in \A^k\cap U(\theta,\frac{\pi}{k})}\hat\Sto_{\theta'}^k(A^0)$.
\begin{thm}\label{thm:theoremForlargerDecomp}
  Let $\theta\in S^1$ be a fixed direction which satisfies
  $\theta\pm\frac{\pi}{2k}\notin\A^k$, then there is a bijection
  \[
    \underset{\cH(A^0)}{%
      \underset{\text{\rotatebox[origin=c]{-90}{$\cong$}}}{\underbrace{%
        \prod_{k\in\cK}\prod_{\alpha\in\A^k}\Sto_\alpha^k(A^0)}}}
    \cong
    \prod_{k\in\cK}\prod_{j\in\{1,\dots,2k\}}
    \hat\Sto_{\theta+j\frac{\pi}{k}}^k(A^0) \,.
  \]
  \begin{s-rem}
    This means that the information of a complete set of Stokes matrices on every level
    $k\in\cK$ can be grouped into $2k$ matrices, which are products of the
    corresponding Stokes matrices.
  \end{s-rem}
\end{thm}
\begin{proof}
  The product map
  \[
    \eta_\theta^k:
    \underset{\alpha\in\A^k\cap U(\theta,\frac{\pi}{k})}\prod
    \underset{S_{\{(j,l)\mid q_j\myrel{\alpha}q_l\}}}{%
      \underset{\text{\rotatebox[origin=c]{-90}{$=$}}}{\underbrace{%
          \SSto_\alpha^k(A^0)}}}
    \overset{\cong}{\longrightarrow}
    \underset{S_{\{(j,l) \mid q_j \underset{\theta}{\prec} q_l \}}}{%
      \underset{\text{\rotatebox[origin=c]{-90}{$=$}}}{\underbrace{%
          \hat\SSto_\theta^k(A^0)}}}
    \,.
  \]
  defined by $\eta_\theta^k(K_1,K_2,\dots,K_m) = K_1\cdot K_2\cdots K_m$ is an
  isomorphism, since Proposition~\ref{prop:capAndCupOfRels} implies the premise
  of Corollary~\ref{cor:composeLevelwise}.

  Denote the product of the isomorphisms $\eta_\theta^k$ as $\eta=
  \prod_{k\in\cK}\prod_{j\in\{1,\dots,2k\}}
  \eta_{\theta+j\frac{\pi}{k}}^k$ and the isomorphism
  $\prod_{k\in\cK}\prod_{j\in\{1,\dots,2k\}}\hat\rho_{\theta+j\frac{\pi}{k}}^k$
  by $\hat\rho$.
  \[ \begin{tikzcd}[column sep=2cm,row sep=1.5cm]
    \prod_{k\in\cK}\prod_{\alpha\in\A^k}\Sto_\alpha^k(A^0)
    \rar[dashed]
    \arrow[d,"\prod_{k\in\cK}\prod_{\alpha\in\A}\rho_{\alpha}^k","\cong"']&
    \underset{k\in\cK}{\prod}\underset{j\in\{1,\dots,2k\}}{\prod}
    \hat\Sto_{\theta+j\frac{\pi}{k}}^k(A^0)
    \arrow[d,"\hat\rho","\cong"']&
    \\
    \prod_{k\in\cK}\prod_{\alpha\in\A^k}\SSto_\alpha^k(A^0)
    \arrow[r,"\eta","\cong"']&
    \underset{k\in\cK}{\prod}\underset{j\in\{1,\dots,2k\}}{\prod}
    \hat\SSto_{\theta+j\frac{\pi}{k}}^k(A^0)
  \end{tikzcd} \]
  The isomorphism, we were looking for, is then given as
  $\bigl(\prod_{k\in\cK}\prod_{\alpha\in\A}\rho_{\alpha}^k\bigr)
  \circ \eta \circ \hat\rho^{-1}$, i.e.\ is the isomorphism making the above
  diagram commute, and we will also denote it by $\eta$.
\end{proof}
\begin{rem}
  The in the previous proof obtained isomorphism
  \[
    \eta: \prod_{\alpha\in\A}\SSto_{\alpha}(A^0) \longrightarrow
    \prod_{k\in\cK} \prod_{j\in\{1,\dots,2k\}}
    \hat\SSto_{\theta+j\frac{\pi}{k}}^k(A^0)
  \]
  is in the single-leveled case with $n=s$, i.e.\ all diagonal elements of $Q$
  are different, given in Lemma 3.2 of Boalch's paper \cite[Lem.3.2]{boalch}.
  In this case is every $\hat\SSto_{\theta+j\frac{\pi}{k}}^k(A^0)$ isomorphic to
  $U_+$, i.e.\ the upper triangular unipotent subgroup of $\Gl_n(\C)$, via the
  map $K\mapsto PKP^{-1}$, where $P$ is the permutation matrix mentioned in
  Corollary~\ref{cor:someCorolaryToLabel}.
  A similar result is also found in Theorem~I of~\cite{BJL1979Birkhoff}.
\end{rem}
But there is in the multileveled case no obvious way to find a subset $J\in S^1$
and an isomorphism
\[
  \prod_{\theta\in J} \hat\Sto_{\theta}(A^0)
  ~~\not{\!\!\!\!\longrightarrow}
  \prod_{k\in\cK}\prod_{j\in\{1,\dots,2k\}}
  \hat\Sto_{\theta+j\frac{\pi}{k}}^k(A^0)\,,
\]
where $\hat\Sto_{\theta}(A^0)$ was defined in Remark~\ref{rem:representation}.

%%% Local Variables:
%%% TeX-master: "Maximilian_Huber-Masters_Thesis-with_notes.tex"
%%% End:

    \fi

    \ifnum\myHideVariable=0
    \section{A taste of summability}\label{app:multisummability}
For the sake of completeness, let us speak a little bit about summability in
this section. We will only mention some facts to give an rough idea how the map
$g$ in the diagram on page~\pageref{page:ofPreDiagram} arises.
The understanding of this section is not necessary for the rest of this
chapter. It is enough to think of the map $g$ (resp.\ $\tilde g$) as a black box
which uses some theory to make the asymptotic expansion unique and yields an
element of $\prod_{\alpha\in\cA}\Sto_\alpha(A^0)$ (resp.\ of
$\Gamma(\dot\cU;\Lambda(A^0))$) corresponding to some marked pair.

\marginnote{\cite[111]{Loday2014}}
The aim of a theory of summation is to associate with any series an asymptotic
function uniquely determined in a way as much natural as possible.
A useful and extensive resources for this topic are Loday-Richaud's
book~\cite{Loday2014}.
A resource which looks only at an special case, which naturally arises from the
previous sections, is Section III of the paper~\cite{Loday1994} and some of the
most important statements, i.e.\ exactly those we need, are found in many places
like Boalch's publications~\cite{boalch,thboalch}, although he looks only at the
single-leveled case.

The following Proposition can be found in~\cite[Prop.III.2.1]{Loday1994}
and~\cite[Thm.4.3.13]{Loday2014}.
\begin{prop}\label{prop:multisummability}
  \marginnote{\cite[Sec.III.2.1]{Loday1994}}
  Let $\cV=(V_j)_{j\in J}$ be a cyclic covering, let $\hat F\in\hat G(A^0)$ be a
  formal transformation and let
  $\dot\phi=(\dot\phi_j)_{j\in J}\in\Gamma(\dot{\mathcal{V}};\Lambda(A^0))$ be a
  $1$-cocycle in the cohomology class $\exp(\hat F)$.

  There exists a \textbf{unique} family of realizations $(F_j)_{j\in J}$ of
  $\hat F$ over $\mathcal{V}$, i.e.\ matrices $F_j$ which are analytic on
  $V_j$, satisfy $[A^0,A]$ and are asymptotic to $\hat F$ on $V_j$, such that
  \[
    \dot\phi_j=F_{j-1}F_j^{-1} \in\Gamma(V_{j-1}\cap V_j;\Lambda(A^0))
  \]
  for every $j\in J$.
\end{prop}
When $\mathcal{V}$ is the, in Section~\ref{sec:proofOfMatrixThm} defined,
cyclic covering $\cU^{\leq k_r}=:\cU$ and $\dot\phi$ is in its Stokes
form\TODO[~(cf.~??)], we call the realizations $F_{\alpha}$, which were obtained
by the previous proposition, the \emph{sums of $\hat F$}.
\begin{defn}\label{defn:sumsLeftRight}
  \marginnote{\cite[Defn.III.2.2]{Loday1994}}
  Denote by $\alpha^+$ the next anti-Stokes direction on the right of $\alpha$
  we can define
  \begin{itemize}
  \item $S_\alpha^-(\hat F):=F_\alpha\in\Gl_n(\cA(U_\alpha))$ as the \emph{sum
      of $\hat F$ on the left of $\alpha$} and
  \item $S_\alpha^+(\hat F):=F_{\alpha^+}\in\Gl_n(\cA(U_{\alpha^+}))$ as the
    \emph{sum of $\hat F$ on the right of $\alpha$}.
  \end{itemize}
\end{defn}
\begin{rem}
  \marginnote{\cite[Defn.3.6]{boalch}}
  $S^+_\alpha(\hat F)t^Le^{Q(t^{-1})}$ is a solution of $[A]$ on the
  corresponding sector.
  \TODO[ or $S^-_\alpha(\hat F)t^Le^{Q(t^{-1})}$?  or
    $t^Le^{Q(t^{-1})}S^+_\alpha(\hat F)$?  or
  $t^Le^{Q(t^{-1})}S^-_\alpha(\hat F)$?]
  Using the equation (\ref{eq:representation}) we can write
  \[
    S^+_\alpha(\hat F)
    \cY_{0,\alpha}(t)
    =
    S^-_\alpha(\hat F)
    \cY_{0,\alpha}(t)C_{\cY_{0,\alpha}}
  \]
  thus the Stokes matrix $C_{\cY_{0,\alpha}}$ is a matrix, which
  describes the \rewrite{blending} between the two adjacent sectors.
  \TODO[see Loday's remarks on thm.4.3.13 in \cite{Loday2014}]
\end{rem}

% \subsubsection{The map
%   $g:\cH(A^0)\to\prod_{\alpha\in\cA}\Sto_\alpha(A^0)$}\label{sect:multisummmap}
There are alternative and more constructive definitions, which are equivalent to
Definition~\ref{defn:sumsLeftRight}.
Since some of them are more direct and do not use the $1$-cocycle, to which they
correspond, they can be used to obtain the corresponding
element in $\prod_{\alpha\in\cA}\Sto_\alpha(A^0)$.
Many different approaches of (multi-)summation can be found in~\cite{Loday2014}.

Let us fix an ambassador $\hat F$ in $G\backslash\hat G(A^0)$ corresponding to
equivalence class of marked pairs in $\cH(A^0)$.
The theory of summation says that there is a map which takes $\hat F$ to a
unique and natural set of functions
\[
  (S_{\alpha_1}^-(\hat F),
  \underset{S_{\alpha_1}^+(\hat F)}{
    \underset{\text{\rotatebox[origin=c]{-90}{$=$}}}{\underbrace{
        S_{\alpha_2}^-(\hat F)
      }}}
  ,\dots,S_{\alpha_\nu}^-(\hat F))\in
  \prod_{\alpha\in\A}\Gl_n(\cA(U_\alpha)) \,,
\]
which have $\hat F$ as asymptotic functions. The corresponding element
$(\phi_{\alpha_1},\dots,\phi_{\alpha_\nu})$ in
$\prod_{\alpha\in\A}\Sto_\alpha(A^0)$ is then found by setting
$\phi_\alpha:=\left(S^-_\alpha(\hat F)\right)̂^{̀-1}S^+_\alpha(\hat F)\in\Gamma(U_\alpha;\Lambda(A^0))$.
\begin{defn}\label{defn:theMapG}
  The corresponding map is denoted by
  \[
    \tilde g:G(\!\{t\}\!)\backslash\hat G(A^0)
    \longrightarrow
    \Gamma(\dot\cU;\Lambda(A^0)) \,.
  \]
  The map $g$ mentioned in the beginning of Chapter~\ref{chap:stokes} is build
  as shown in Section~\ref{sec:theCompleteDiagram}.
\end{defn}

%%% Local Variables:
%%% TeX-master: "Maximilian_Huber-Masters_Thesis-with_notes.tex"
%%% End:

    \fi

    \ifnum\myHideVariable=0
    \pagebreak %%%%%%%%%%%%%%%%%%%%%%%%%%%%%%%%%%%%%%%%%%%%%%%%%%%%%%%%%%%%%%%%%%%%
%%%%%%%%%%%%%%%%%%%%%%%%%%%%%%%%%%%%%%%%%%%%%%%%%%%%%%%%%%%%%%%%%%%%%%%%%%%%%%%
\section{The complete diagram}\label{sec:theCompleteDiagram}
In this section we want to write the improved version of the diagram from
page~\pageref{page:ofPreDiagram} down. It will contain nearly all of the
isomorphisms which were defined in the previous sections. Additionally will
isomorphisms like $g$ be decomposed into their building blocks.

In the following diagram of \textbf{isomorphisms of pointed sets} was the
isomorphism $\exp_{A^0}$ discussed in Section~\ref{sec:mainThm1}. The
\textcolor{purple}{purple} path, which is the isomorphism $h$ from
Theorem~\ref{thm:mainThm2}, together with the lower left part was discussed
in Section~\ref{sec:mainThm2}. The lower right part was discussed in
Section~\label{sec:furtherImprovements} and the \textcolor{blue}{blue} part on
the left was discussed in the previous section.
\begin{center}
  \begin{tikzpicture}[scale=0.95,transform shape]
    \pgfmathsetmacro{\dx}{4.5cm}
    \pgfmathsetmacro{\ldx}{6cm}
    \pgfmathsetmacro{\dy}{2.5cm}
    \pgfmathsetmacro{\mdy}{-1.5cm}

    \node (cH) at (-6,3.5) {$\cH(A^0)$};
    \node (cH2) at ([yshift=-\dy]cH)
      {$G(\!\{t\}\!)\backslash\hat G(A^0)$};
    \node (cH3) at ([yshift=-\dy]cH2)
      {$\Gamma(\dot\cU;\Lambda(A^0))$};
    \node (ProdOfStos) at (-6,-4)
      {$\displaystyle\prod_{\alpha\in\A}\Sto_{\alpha}(A^0)$};
    \node (S3P) at ([xshift={\dx},yshift={\dy}]ProdOfStos)
      {$\displaystyle\prod_{k\in\cK}\prod_{\alpha\in\A}\Sto_{\alpha}^{k}(A^0)$};
    \node (S2P) at ([xshift={\dx},yshift={\dy}]S3P)
      {$\displaystyle\prod_{k\in\cK}\Gamma(\dot\cU^k;\Lambda^k(A^0))$};
    \node (StA0) at ([xshift={\dx},yshift={\dy}]S2P)
      {$\St(A^0)$};
    \node (ProdOfSStos) at ([yshift=-{\dy}]ProdOfStos)
      {$\displaystyle\prod_{\alpha\in\A}\SSto_{\alpha}(A^0)$};
    \node (SS3P) at ([yshift=-{\dy}]S3P)
      {$\displaystyle\prod_{k\in\cK}\prod_{\alpha\in\A}\SSto_{\alpha}^{k}(A^0)$};
    \node (SS2P) at ([yshift={\mdy},xshift={\ldx}]S3P)
      {$\displaystyle\prod_{k\in\cK} \prod_{j\in\{1,\dots,2k\}}
        \hat\Sto_{\theta+j\frac{\pi}{k}}^k(A^0)$};
    \node (SSS2P) at ([yshift={\mdy},xshift={\ldx}]SS3P)
      {$\displaystyle\prod_{k\in\cK} \prod_{j\in\{1,\dots,2k\}}
        \hat\SSto_{\theta+j\frac{\pi}{k}}^k(A^0)$};
    \node[] (SSS3P) at ([yshift={\mdy},xshift={\ldx}]ProdOfSStos)
      {$\displaystyle\prod_{\alpha\in\A}
        \C^{\sum_{q_j\myrel{\alpha}q_l}
          \deg(q_j-q_l)\cdot n_j\cdot n_l}$};

    \draw[->] (cH2) edge [->,transform canvas={xshift=2mm}]
      node[midway,right]
      {\footnotesize(cf.\ Cor.\ref{cor:isomorphyOfClassfset})}
      (cH);
    \draw[->] (cH) -- (StA0) node[midway,above] {$\exp_{A^0}$};
    \draw[->,purple] (ProdOfStos) to
      node[midway,above,sloped]
      {$\chi\circ\prod_{\alpha\in\A}i_\alpha$}
      (S3P);
    \draw[->,purple] (S3P) -- (S2P)
      node[midway,above,sloped] {$\prod_{k\in\cK}i^k$};
    \draw[->,purple] (S2P) -- (StA0)
      node[midway,above left] {$\cT$};
    \draw[->] (ProdOfStos) -- (ProdOfSStos)
      node[midway,left] {$\prod_{\alpha\in\A}\rho_\alpha$};
    \draw[->] (ProdOfSStos) -- (SS3P)
      node[midway,below,sloped] {$\chi\circ\prod_{\alpha\in\A}i_\alpha$};
    \draw[->] (S3P) -- (SS3P)
      node[midway,right] {$\prod_{k\in\cK}\prod_{\alpha\in\A}\rho_{\alpha}^k$};
    \draw[->] (S3P) -- (SS2P) node[midway,above] {$\eta$};
    \draw[->] (SS3P) -- (SSS2P) node[midway,above] {$\eta$};
    \draw[->] (SS2P) -- (SSS2P)
      node[midway,right]
      {$\hat\rho$};
    \draw[->] (SSS3P) -- (ProdOfSStos)
      node[midway,below,sloped] {$\prod_{\alpha\in\A}\vartheta_\alpha$};

    % \draw[->] (cH) to[out=220,in=100]  node[midway,left] {$g$} (ProdOfStos);
    \draw[->,blue] (cH) -- (cH2);
    \draw[->,blue] (cH3) -- (ProdOfStos)
      node[midway,left] {\footnotesize(take germs)};
    \draw[->,blue] (cH2) -- (cH3)
      node[midway,right] {$\tilde g$};

    \draw[->,gray] (ProdOfStos) to[out=60,in=185] node[midway,above] {$\mathfrak{T}$} (S2P);
    \ifnum\myDevelopVariable=1
        \draw[right hook->,dashed,red] (S2P) to[in=39,out=180] node[midway,above] {$\tau$} (cH3);
        \draw[->>,dashed,red] (cH3) to[out=52,in=195] node[midway,above,sloped] {(quotient map)} (StA0);
    \fi
  \end{tikzpicture}
\end{center}
Where we denote
\begin{itemize}
\item $\chi:\displaystyle \prod_{\alpha\in\A}\prod_{k\in\cK}\Sto_\alpha^k(A^0)
  \equiv\displaystyle \prod_{k\in\cK}\prod_{\alpha\in\A}\Sto_\alpha^k(A^0)$
  the reordering and
\item by abuse of notation we also denote the Stokes matrix version in
  the same way
  $\chi:\displaystyle \prod_{\alpha\in\A}\prod_{k\in\cK}\SSto_\alpha^k(A^0)
  \equiv\displaystyle
  \prod_{k\in\cK}\prod_{\alpha\in\A}\SSto_\alpha^k(A^0)$.
\end{itemize}

%%% Local Variables:
%%% TeX-master: "Maximilian_Huber-Masters_Thesis-with_notes.tex"
%%% End:

    \fi
\fi

\pagebreak
\appendix
\ifnum\myHideVariable=0
  %%%%%%%%%%%%%%%%%%%%%%%%%%%%%%%%%%%%%%%%%%%%%%%%%%%%%%%%%%%%%%%%%%%%%
  \ifnum\myDevelopVariable=1
    \addcontentsline{toc}{chapter}{Appendix}
    \input{app-landkarte.tex}
    % generated with mklist.sh
\chapter{Notes}
\section{Meromorphe Zusammenhänge vs. Systeme}
Sei $\Sigma$ eine Riemann Fläche, $D=k_1(a_1)+\dots+k_m(a_m)$ ein effektiver
Divisor darauf.
\begin{paracol}{2}
\begin{defn}[Meromorpher Zusammenhang]
Ein meromorpher Zusammenhang $\nabla$ auf $V$ mit Polen auf $D$ ist eine
Abbildung
\[
  \nabla: V\to V\otimes K(D)
\]
so dass die Leibnitzregel
\begin{equation}
  \nabla(fv)=(df)\otimes v + f\nabla v
\end{equation}
erfüllt ist. Hierbei ist $v$ ein lokaler Schnitt von $V$, $f$ eine lokale
holomorphe Funktion und $K$ die Garbe der holomorphen Eins-Formen.
\end{defn}
\switchcolumn %%%%%%%%%%%%%%%%%%%%%%%%%%%%%%%%%%%%%%%%%%%%%%%%%%%%%%%%%%%%%%%%%
\begin{defn}
Ein \emph{System} ist \dots
\end{defn}
Dieses wird erhalten durch lokale Trivialisierung.
\end{paracol}

\begin{paracol}{2}
\begin{defn}[compatible framing]
\end{defn}
\switchcolumn %%%%%%%%%%%%%%%%%%%%%%%%%%%%%%%%%%%%%%%%%%%%%%%%%%%%%%%%%%%%%%%%%
\begin{defn}[marked pair]
\end{defn}
\end{paracol}

\section{Formal model vs. Formal Decomposition}
\begin{paracol}{2}\sloppy
\switchcolumn[0]\noindent
  Write $A$ as
  \[
    A=\left(\sum^{0}_{j=k_i}A_j\frac{dz}{z^{j}}\right)+A_0dz+\dots
  \]
  \begin{defn}[\cite{boalch} Def 2.2]
    A meromorphic connection is a \emph{nice (or generic) meromorphic connection}
    if
    \begin{itemize}
      \item at each $a_i$
    \end{itemize}
    the leading coefficient ${}^iA_{k_i}$ is
    \begin{itemize}
      \item diagonalisable with distinct eigenvalues and $k_i\geq2$, or
      \item diagonalisable with distinct eigenvalues $\mod\Z$ and $k_i=1$
    \end{itemize}
  \end{defn}
\switchcolumn[1]\noindent
  \begin{thm}[Formal decomposition (II.5.7)]
    Let $(\cM,\nabla)$ be a germ of meromorphic bundle with connection,
    \begin{itemize}
      \item equipped with a basis in which the matrix $\Omega$ takes the form
        \[
          \Omega=t^{-r}A(t)\frac{dt}{t}
        \]
        with $r\geq1$, $A$ having holomorphic entries, and $A_0:=A(0)$
        being regular semisimple\footnote{i.e.\ with pairwise distinct
        eigenvalues.}
    \end{itemize}
    Then there exist
    \begin{itemize}
      \item a model $(\cM^{good},\nabla)
        =\left(\bigoplus_\phi(\cE^\phi\otimes\cR_\phi) \right)$ and
      \item a \textcolor{red!40!black}{`formal'} isomorphism
        \[
          \textcolor{red!40!black}{\hat\cO_{D}\otimes}(\cM,\nabla)
          \overset{\sim}{\longrightarrow}
        \textcolor{red!40!black}{\hat\cO_{D}\otimes}
          \left(\bigoplus_\phi(\cE^\phi\otimes\cR_\phi) \right)
        \]
    \end{itemize}
  \end{thm}
\end{paracol}


\section{Reading text from Loday-Richaud}
\begin{itemize}
  \item \textbf{The paper:}
    \cite{Loday1994}
    \begin{itemize}
      \item Stokes matrices are not all germs of isotropies, only the Sokes
        germs! see 4.3.13
      \item for \textbf{Sheaf $\to$ Matrix}
        \begin{thm}[II.2.1]
          The map
          \[
            h:\prod_{\alpha\in\A}\Sto_\alpha(A_0)\to H^1(S^1;\Lambda(A_0))
          \]
          is bijective and natural.
        \end{thm}
    \end{itemize}
  \item \textbf{The book:} \cite{Loday2014}
    \begin{itemize}
      \item multilevels, multisectors
    \end{itemize}
  \item others:
    \begin{itemize}
      \item \cite{LodayRichaud2004}
    \end{itemize}
\end{itemize}

$k -1 \to k$

\subsection{Determination $\tilde\theta$ of $\theta$}
\begin{center}
  \begin{tikzpicture}[scale=1]
    \node[] (zero) at (0,0) {};
    \fill (zero) circle (1pt);
    \draw (0,-0.5) arc (270:-90:2 and 0.5);

    \node (theta) at ({cos(240) * 2},{sin(240) * 0.5}) {};
    % \node [below left of=theta,blue!40!white] {$\theta$};

    \draw[->,gray!50!white] (-1,1.6) to[out=250,in=120] (theta);
    \draw[->,gray!50!white] (-1,3.6) to[out=250,in=120] (theta);
    \draw[->,gray!50!white] (-1,4.6) to[out=250,in=120] (theta);
    \draw[->,green!60!black,thick] (-1,2.6) to[out=250,in=120] (theta);

    \draw[dotted] (0,0) -- (theta);
    \fill[blue!20!white] (theta) circle (1pt);


    \draw[] (-1,2) arc (270:200:-3 and -0.7);
    \draw[] (-1,2) arc (270:90:1 and -0.2) arc (270:90:-3 and 0.7);
    \fill[blue!20!white] (-1,1.6) circle (1pt);
    \draw[] (-1,3) arc (270:90:1 and -0.2) arc (270:90:-3 and 0.7);
    \fill[green!60!black] (-1,2.6) circle (1pt);
    \draw[] (-1,4) arc (270:90:1 and -0.2) arc (270:90:-3 and 0.7);
    \fill[blue!20!white] (-1,3.6) circle (1pt);
    \draw[] (-1,5) arc (270:90:1 and -0.2) arc (270:200:-3 and 0.7);
    \fill[blue!20!white] (-1,4.6) circle (1pt);
  \end{tikzpicture}
\end{center}

\subsection{Definition of $\Lambda_\theta(A_0)$}
$\Lambda_\theta(A_0)$ $:\Leftrightarrow{}$ sheaf of flat isotropies over $S^1$
\marginnote{p. 854f \cite{Loday1994}}
\begin{defn}
  A germ of $\Lambda_\theta(A_0)$ at $\theta_0\in S^1$ is
  \begin{itemize}
    \item an invertible Matrix $f\in\Gl_n(\cO(U))$
      \begin{itemize}
        \item for suitable arc $U=U(\theta_0,\epsilon,\epsilon')$
      \end{itemize}
      satisfying:
      \begin{enumerate}
        \item \emph{Flatness:}
          \[
            \underset{x\in U}{\underset{x\to0}{\lim}}f(x)=\Id
            \text{ and }
            f\underset{U}{\sim}I
          \]
        \item \emph{Isotropy of $[A_0]$: ${}^f\!A_0=A_0$.}
      \end{enumerate}
  \end{itemize}
\end{defn}

\begin{center}
  \begin{tikzpicture}[scale=1]
    \node[] (zero) at (0,0) {};
    \fill (zero) circle (1pt);

    \fill[blue!20!white] (theta) circle (1pt);

    \filldraw[fill=green!20!white
      ,draw=green!60!black
      ,thick
    ,path fading=west] (0,0)
      -- ({cos( -50 )*4},{sin( -50 )*1}) arc (-50:20:4 and 1)
      -- cycle;

    \filldraw[fill=green!20!white
      ,draw=green!60!black
      ,thick
    ,path fading=west] (0,4)
      -- ({cos( -50 )*0.5},{4 + sin( -50 )*2})
      -- ({cos( -50 )*0.5 + cos( -10 )*3}
        ,{4 + sin( -50 )*2 + sin( -10 )*12})
      arc (-80:40:1 and 3)
      -- ({cos( 20 )*0.5},{4 + sin( 20 )*2})
      -- cycle;

    \fill[fill=white,draw=black] (0,2) arc (270:-90:0.5 and 2);

    \draw[blue!60!black,path fading=east]
      (zero) -- ({cos( -15 )*4},{sin( -15 )*1});

    \draw[dashed] (zero) -- (0,2);
  \end{tikzpicture}
\end{center}

\subsection{Definition of $\Sto_\alpha(A_0)$}
\begin{itemize}
  \item p. 861 in the paper
  \item p. 78 in the book
\end{itemize}

\subsection{Definition of $\cF$}

\subsection{Adequate coverings}
\begin{center}
  \begin{tikzpicture}[scale=2]
    \node (zero) at (0,0) {};
    \draw (zero) circle (1cm);

    \node at (-1,1) {$\cK=\{5\}:$};

    \foreach \w in {55}
    {\foreach \sep in {0,45,90,135,180,225,270,315}
     {\draw[thick,purple] (0,0) -- +({cos( \w + \sep )},{sin( \w + \sep )});}};

    \fill[white] (zero) circle (1.5pt);
    \fill (zero) circle (.7pt);
  \end{tikzpicture}
  \begin{tikzpicture}[scale=2]
    \node (zero) at (0,0) {};
    \draw (zero) circle (1cm);

    \node at (-1,1) {$\cK=\{5,4\}:$};

    \foreach \w in {55}
    {\foreach \sep in {0,45,90,135,180,225,270,315}
     {\draw[thick,purple] (0,0) -- +({cos( \w + \sep )},{sin( \w + \sep )});}};

    \foreach \w in {10,20,55}
    {\foreach \sep in {0,60,120,180,240,300}
     {\draw[blue] (0,0) -- +({cos( \w + \sep )},{sin( \w + \sep )});}};

    \fill[white] (zero) circle (1.5pt);
    \fill (zero) circle (.7pt);
  \end{tikzpicture}
  \begin{tikzpicture}[scale=2]
    \node (zero) at (0,0) {};
    \draw (zero) circle (1cm);

    \node at (-1,1) {$\cK=\{5,4,3\}:$};

    \foreach \w in {55}
    {\foreach \sep in {0,45,90,135,180,225,270,315}
     {\draw[thick,purple] (0,0) -- +({cos( \w + \sep )},{sin( \w + \sep )});}};

    \foreach \w in {10,20,55}
    {\foreach \sep in {0,60,120,180,240,300}
     {\draw[blue] (0,0) -- +({cos( \w + \sep )},{sin( \w + \sep )});}};

    \foreach \w in {30,35}
    {\foreach \sep in {0,90,180,270}
     {\draw[red] (0,0) -- +({cos( \w + \sep )},{sin( \w + \sep )});}};

    \fill[white] (zero) circle (1.5pt);
    \fill (zero) circle (.7pt);
  \end{tikzpicture}
\end{center}

Let $k\in\cK$.
\subsubsection{The \emph{cyclic covering $\cU^k=\{U_\alpha^k;\alpha\in\A^k\}$}}
\begin{center}
  \begin{tikzpicture}[scale=2.5]
    \node (zero) at (0,0) {};
    \draw (zero) circle (1cm);

    \node at (-1,1) {$\textcolor{purple}{\dot\cU^{5}}:$};
    \node at (-1.6,1.6) {$\textcolor{blue}{\dot\cU^{4}}:$};
    \node at (-2,2) {$\textcolor{green!60!black}{\dot\cU^{3}}:$};

    %%%%%%%%%%%%%%%%%%%%%%%%%%%%%%%%%%%%%%%%%%%%%%%%%%%%%%%%%%%%%%%%%%%%%%%%%%%
    %%%%  Purple
    \foreach \w in {10,25}
    {\foreach \sep in {0,45,90,135,180,225,270,315}
     {\draw[thick,purple] (0,0) -- +({cos( \w + \sep )},{sin( \w + \sep )});

      \pgfmathsetmacro\r{{1.1 + mod(\w,10)/100 + mod(\sep,2)/10}}
      \draw[dotted,purple] ({cos( \w + \sep )},{sin( \w + \sep )})
        -- ({cos( \w + \sep) * \r},{sin( \w + \sep) * \r});

      \draw[thick,purple] ({cos( \w + \sep -22) * \r},{sin( \w + \sep -22) * \r})
        arc ({\w + \sep - 22}:{\w + \sep +22}:\r);
      }};

    %%%%%%%%%%%%%%%%%%%%%%%%%%%%%%%%%%%%%%%%%%%%%%%%%%%%%%%%%%%%%%%%%%%%%%%%%%%
    %%%%  Blue
    \foreach \w in {10,20,55}
    {\foreach \sep in {0,60,120,180,240,300}
     {\draw[blue] (0,0) -- +({cos( \w + \sep ) * 0.9},{sin( \w + \sep ) * 0.9});}};

    \foreach \sep in {0,180}
    {\pgfmathsetmacro\rr{1.6}
     \foreach \w in {10,20,55
                    ,70,80,115
                    ,130,140,175}
     {\pgfmathsetmacro\rr{{\rr + (mod(\w,60) - 5*mod(\w,10))/150 
                               + mod(\w + \sep - mod(\w,60),120)/300}}
      \draw[dotted,blue] ({cos( \w + \sep )},{sin( \w + \sep )})
        -- ({cos( \w + \sep) * \rr},{sin( \w + \sep) * \rr});
      \draw[blue,thick] ({cos(\sep+\w -29) * \rr},{sin(\sep+\w -29) * \rr})
        arc ({\sep+\w -29}:{\sep+\w +29}:\rr);
     }
    }

    %%%%%%%%%%%%%%%%%%%%%%%%%%%%%%%%%%%%%%%%%%%%%%%%%%%%%%%%%%%%%%%%%%%%%%%%%%%
    %%%%  Green
    \foreach \w in {30,35}
    {\foreach \sep in {0,90,180,270}
     {\draw[green!60!black] (0,0) -- +({cos( \w + \sep ) * 0.8},{sin( \w + \sep) * 0.8});}};

    \foreach \sep in {0,180}
    {\pgfmathsetmacro\rr{2.4}
     \foreach \w in {30,35
                    ,120,125}
     {\pgfmathsetmacro\rr{{\rr + mod(\w,10)/100 + mod(\w-mod(\w,90),180)/900}}
      \draw[dotted,green!60!black] ({cos( \w + \sep )},{sin( \w + \sep )})
        -- ({cos( \w + \sep) * \rr},{sin( \w + \sep) * \rr});
      \draw[green!60!black,thick] ({cos(\sep+\w -44) * \rr},{sin(\sep+\w -44) * \rr})
        arc ({\sep+\w -44}:{\sep+\w +44}:\rr);
     }
    }

    \fill[white] (zero) circle (1.5pt);
    \fill (zero) circle (.7pt);
  \end{tikzpicture}
\end{center}

\subsubsection{The \emph{cyclic covering
  $\cU^{\leq k}=\{U_\alpha^{\leq k};\alpha\in\A^{\leq k}\}$}}
\begin{center}
  \begin{tikzpicture}[scale=2.5]
    \node (zero) at (0,0) {};
    \draw (zero) circle (1cm);

    \node at (-1.05,1.05) {$\textcolor{purple}{\dot\cU^{\leq 5}}:$};
    \node at (-1.65,1.65) {$\textcolor{blue}{\dot\cU^{\leq 4}}:$};
    \node at (-2.05,2.05) {$\textcolor{green!60!black}{\dot\cU^{\leq 3}}:$};

    %%%%%%%%%%%%%%%%%%%%%%%%%%%%%%%%%%%%%%%%%%%%%%%%%%%%%%%%%%%%%%%%%%%%%%%%%%%
    %%%%  Purple
    \foreach \w in {10,25}
    % {\foreach \sep in {0,45,90,135,180,225,270,315}
    {\foreach \sep in {0,36,72,108,144,180,216,252,288,324}
     {\draw[thick,purple] (0,0) -- +({cos( \w + \sep )},{sin( \w + \sep )});

      \pgfmathsetmacro\r{{1.1 + mod(\w,10)/100 + mod(\sep,72)/360}}
      \draw[dotted,purple] ({cos( \w + \sep )},{sin( \w + \sep )})
        -- ({cos( \w + \sep) * \r},{sin( \w + \sep) * \r});

      \draw[thick,purple] ({cos( \w + \sep -17.5) * \r},{sin( \w + \sep -17.5) * \r})
        arc ({\w + \sep - 17.5}:{\w + \sep +17.5}:\r);

      % \draw[purple] ({cos( \w + \sep -17.5) * \r},{sin( \w + \sep -17.5) * \r})
      %   -- ({cos( \w + \sep -17.5) * 1.5},{sin( \w + \sep -17.5) * 1.5});
      % \draw[purple] ({cos( \w + \sep +17.5) * \r},{sin( \w + \sep +17.5) * \r})
      %   -- ({cos( \w + \sep +17.5) * 1.5},{sin( \w + \sep +17.5) * 1.5});
      }};

    %%%%%%%%%%%%%%%%%%%%%%%%%%%%%%%%%%%%%%%%%%%%%%%%%%%%%%%%%%%%%%%%%%%%%%%%%%%
    %%%%  Blue
    \foreach \w in {10,20,55}
    {\foreach \sep in {0,60,120,180,240,300}
     {\draw[blue] (0,0) -- +({cos( \w + \sep ) * 0.9},{sin( \w + \sep ) * 0.9});}};

    \foreach \sep in {0,180}
    {\foreach \w in {10,20,55
                    ,70,80,115
                    ,130,140,175}
     {\pgfmathsetmacro\rr{{1.7 + (mod(\w,60) - 5*mod(\w,10))/150 
                               + mod(\w + \sep - mod(\w,60),120)/300}}
      \draw[dotted,blue] ({cos( \w + \sep )},{sin( \w + \sep )})
        -- ({cos( \w + \sep) * \rr},{sin( \w + \sep) * \rr});
      \draw[gray!42!white] ({cos(\sep+\w -29) * \rr},{sin(\sep+\w -29) * \rr})
        arc ({\sep+\w -29}:{\sep+\w +29}:\rr);

      % \draw[gray!42!white] ({cos( \w + \sep -29) * \rr},{sin( \w + \sep -29) * \rr})
      %   -- ({cos( \w + \sep -29) * 1.5},{sin( \w + \sep -29) * 1.5});
      % \draw[gray!42!white] ({cos( \w + \sep +29) * \rr},{sin( \w + \sep +29) * \rr})
      %   -- ({cos( \w + \sep +29) * 1.5},{sin( \w + \sep +29) * 1.5});
     }

     \pgfmathsetmacro\w{10}
     \pgfmathsetmacro\rr{{1.7 + (mod(\w,60) - 5*mod(\w,10))/150 
                              + mod(\w + \sep - mod(\w,60),120)/300}}
     \draw[thick,blue] ({cos(\sep+\w -17.5) * \rr},{sin(\sep+\w -17.5) * \rr})
       arc ({\sep+\w -17.5}:{\sep+\w +17.5}:\rr);

     \pgfmathsetmacro\w{20}
     \pgfmathsetmacro\rr{{1.7 + (mod(\w,60) - 5*mod(\w,10))/150 
                              + mod(\w + \sep - mod(\w,60),120)/300}}
     \draw[thick,blue] ({cos(\sep+ -7.5) * \rr},{sin(\sep+ -7.5) * \rr})
       arc ({\sep+ -7.5}:{\sep+ 42.5}:\rr);

     \pgfmathsetmacro\w{55}
     \pgfmathsetmacro\rr{{1.7 + (mod(\w,60) - 5*mod(\w,10))/150 
                              + mod(\w + \sep - mod(\w,60),120)/300}}
     \draw[thick,blue] ({cos(\sep+ 28.5) * \rr},{sin(\sep+ 28.5) * \rr})
       arc ({\sep+ 28.5}:{\sep+ 78.5}:\rr);

     \pgfmathsetmacro\w{70}
     \pgfmathsetmacro\rr{{1.7 + (mod(\w,60) - 5*mod(\w,10))/150 
                              + mod(\w + \sep - mod(\w,60),120)/300}}
     \draw[thick,blue] ({cos(\sep+ 43.5) * \rr},{sin(\sep+ 43.5) * \rr})
       arc ({\sep+ 43.5}:{\sep + \w + 29}:\rr);

     \pgfmathsetmacro\w{80}
     \pgfmathsetmacro\rr{{1.7 + (mod(\w,60) - 5*mod(\w,10))/150 
                              + mod(\w + \sep - mod(\w,60),120)/300}}
     \draw[thick,blue] ({cos(\sep+ \w - 29) * \rr},{sin(\sep+ \w - 29) * \rr})
       arc ({\sep+ \w - 29}:{\sep + 99.5}:\rr);

     \pgfmathsetmacro\w{115}
     \pgfmathsetmacro\rr{{1.7 + (mod(\w,60) - 5*mod(\w,10))/150 
                              + mod(\w + \sep - mod(\w,60),120)/300}}
     \draw[thick,blue] ({cos(\sep+ \w - 29) * \rr},{sin(\sep+ \w - 29) * \rr})
       arc ({\sep+ \w - 29}:{\sep + 135.5}:\rr);

     \pgfmathsetmacro\w{130}
     \pgfmathsetmacro\rr{{1.7 + (mod(\w,60) - 5*mod(\w,10))/150 
                              + mod(\w + \sep - mod(\w,60),120)/300}}
     \draw[thick,blue] ({cos(\sep+ \w - 29) * \rr},{sin(\sep+ \w - 29) * \rr})
       arc ({\sep+ \w - 29}:{\sep + 150.5}:\rr);

     \pgfmathsetmacro\w{140}
     \pgfmathsetmacro\rr{{1.7 + (mod(\w,60) - 5*mod(\w,10))/150 
                              + mod(\w + \sep - mod(\w,60),120)/300}}
     \draw[thick,blue] ({cos(\sep+ 115.5) * \rr},{sin(\sep+ 115.5) * \rr})
       arc ({\sep+ 115.5}:{\sep + \w + 29}:\rr);

     \pgfmathsetmacro\w{175}
     \pgfmathsetmacro\rr{{1.7 + (mod(\w,60) - 5*mod(\w,10))/150 
                              + mod(\w + \sep - mod(\w,60),120)/300}}
     \draw[thick,blue] ({cos(\sep+ 151.5) * \rr},{sin(\sep+ 151.5) * \rr})
       arc ({\sep+ 151.5}:{\sep + \w + 29}:\rr);
    }

    %%%%%%%%%%%%%%%%%%%%%%%%%%%%%%%%%%%%%%%%%%%%%%%%%%%%%%%%%%%%%%%%%%%%%%%%%%%
    %%%%  Green
    \foreach \w in {30,35}
    {\foreach \sep in {0,90,180,270}
     {\draw[green!60!black] (0,0) -- +({cos( \w + \sep ) * 0.8},{sin( \w + \sep) * 0.8});}};

    \foreach \sep in {0,180}
    {\foreach \w in {30,35
                    ,120,125}
     {\pgfmathsetmacro\rr{{2.5 + mod(\w,10)/100 + mod(\w-mod(\w,90),180)/900}}
      \draw[dotted,green!60!black] ({cos( \w + \sep )},{sin( \w + \sep )})
        -- ({cos( \w + \sep) * \rr},{sin( \w + \sep) * \rr});
      \draw[gray!42!white] ({cos(\sep+\w -44) * \rr},{sin(\sep+\w -44) * \rr})
        arc ({\sep+\w -44}:{\sep+\w +44}:\rr);
     }


     \pgfmathsetmacro\w{30}
     \pgfmathsetmacro\rr{{2.5 + mod(\w,10)/100 + mod(\w-mod(\w,90),180)/900}}
     \draw[thick,green!60!black] ({cos(\sep+ 7.5) * \rr},{sin(\sep+ 7.5) * \rr})
       arc ({\sep+ 7.5}:{\sep +63.5}:\rr);
     \pgfmathsetmacro\w{35}
     \pgfmathsetmacro\rr{{2.5 + mod(\w,10)/100 + mod(\w-mod(\w,90),180)/900}}
     \draw[thick,green!60!black] ({cos(\sep+ 7.5) * \rr},{sin(\sep+ 7.5) * \rr})
       arc ({\sep+ 7.5}:{\sep +63.5}:\rr);

     \pgfmathsetmacro\w{120}
     \pgfmathsetmacro\rr{{2.5 + mod(\w,10)/100 + mod(\w-mod(\w,90),180)/900}}
     \draw[thick,green!60!black] ({cos(\sep+ 100.5) * \rr},{sin(\sep+ 100.5) * \rr})
       arc ({\sep+ 100.5}:{\sep +150.5}:\rr);
     \pgfmathsetmacro\w{125}
     \pgfmathsetmacro\rr{{2.5 + mod(\w,10)/100 + mod(\w-mod(\w,90),180)/900}}
     \draw[thick,green!60!black] ({cos(\sep+ 100.5) * \rr},{sin(\sep+ 100.5) * \rr})
       arc ({\sep+ 100.5}:{\sep +150.5}:\rr);
    }

    \fill[white] (zero) circle (1.5pt);
    \fill (zero) circle (.7pt);
  \end{tikzpicture}
\end{center}

\subsection{Proof}
\subsubsection{Unique level} \marginnote{p. 872}
\subsubsection{Multi level} \marginnote{p. 872}


    % generated with mklist.sh
\chapter{Notes}
\section{Meromorphe Zusammenhänge vs. Systeme}
Sei $\Sigma$ eine Riemann Fläche, $D=k_1(a_1)+\dots+k_m(a_m)$ ein effektiver
Divisor darauf.
\begin{paracol}{2}
\begin{defn}[Meromorpher Zusammenhang]
Ein meromorpher Zusammenhang $\nabla$ auf $V$ mit Polen auf $D$ ist eine
Abbildung
\[
  \nabla: V\to V\otimes K(D)
\]
so dass die Leibnitzregel
\begin{equation}
  \nabla(fv)=(df)\otimes v + f\nabla v
\end{equation}
erfüllt ist. Hierbei ist $v$ ein lokaler Schnitt von $V$, $f$ eine lokale
holomorphe Funktion und $K$ die Garbe der holomorphen Eins-Formen.
\end{defn}
\switchcolumn %%%%%%%%%%%%%%%%%%%%%%%%%%%%%%%%%%%%%%%%%%%%%%%%%%%%%%%%%%%%%%%%%
\begin{defn}
Ein \emph{System} ist \dots
\end{defn}
Dieses wird erhalten durch lokale Trivialisierung.
\end{paracol}

\begin{paracol}{2}
\begin{defn}[compatible framing]
\end{defn}
\switchcolumn %%%%%%%%%%%%%%%%%%%%%%%%%%%%%%%%%%%%%%%%%%%%%%%%%%%%%%%%%%%%%%%%%
\begin{defn}[marked pair]
\end{defn}
\end{paracol}

\section{Formal model vs. Formal Decomposition}
\begin{paracol}{2}\sloppy
\switchcolumn[0]\noindent
  Write $A$ as
  \[
    A=\left(\sum^{0}_{j=k_i}A_j\frac{dz}{z^{j}}\right)+A_0dz+\dots
  \]
  \begin{defn}[\cite{boalch} Def 2.2]
    A meromorphic connection is a \emph{nice (or generic) meromorphic connection}
    if
    \begin{itemize}
      \item at each $a_i$
    \end{itemize}
    the leading coefficient ${}^iA_{k_i}$ is
    \begin{itemize}
      \item diagonalisable with distinct eigenvalues and $k_i\geq2$, or
      \item diagonalisable with distinct eigenvalues $\mod\Z$ and $k_i=1$
    \end{itemize}
  \end{defn}
\switchcolumn[1]\noindent
  \begin{thm}[Formal decomposition (II.5.7)]
    Let $(\cM,\nabla)$ be a germ of meromorphic bundle with connection,
    \begin{itemize}
      \item equipped with a basis in which the matrix $\Omega$ takes the form
        \[
          \Omega=t^{-r}A(t)\frac{dt}{t}
        \]
        with $r\geq1$, $A$ having holomorphic entries, and $A_0:=A(0)$
        being regular semisimple\footnote{i.e.\ with pairwise distinct
        eigenvalues.}
    \end{itemize}
    Then there exist
    \begin{itemize}
      \item a model $(\cM^{good},\nabla)
        =\left(\bigoplus_\phi(\cE^\phi\otimes\cR_\phi) \right)$ and
      \item a \textcolor{red!40!black}{`formal'} isomorphism
        \[
          \textcolor{red!40!black}{\hat\cO_{D}\otimes}(\cM,\nabla)
          \overset{\sim}{\longrightarrow}
        \textcolor{red!40!black}{\hat\cO_{D}\otimes}
          \left(\bigoplus_\phi(\cE^\phi\otimes\cR_\phi) \right)
        \]
    \end{itemize}
  \end{thm}
\end{paracol}


\section{Reading text from Loday-Richaud}
\begin{itemize}
  \item \textbf{The paper:}
    \cite{Loday1994}
    \begin{itemize}
      \item Stokes matrices are not all germs of isotropies, only the Sokes
        germs! see 4.3.13
      \item for \textbf{Sheaf $\to$ Matrix}
        \begin{thm}[II.2.1]
          The map
          \[
            h:\prod_{\alpha\in\A}\Sto_\alpha(A_0)\to H^1(S^1;\Lambda(A_0))
          \]
          is bijective and natural.
        \end{thm}
    \end{itemize}
  \item \textbf{The book:} \cite{Loday2014}
    \begin{itemize}
      \item multilevels, multisectors
    \end{itemize}
  \item others:
    \begin{itemize}
      \item \cite{LodayRichaud2004}
    \end{itemize}
\end{itemize}

$k -1 \to k$

\subsection{Determination $\tilde\theta$ of $\theta$}
\begin{center}
  \begin{tikzpicture}[scale=1]
    \node[] (zero) at (0,0) {};
    \fill (zero) circle (1pt);
    \draw (0,-0.5) arc (270:-90:2 and 0.5);

    \node (theta) at ({cos(240) * 2},{sin(240) * 0.5}) {};
    % \node [below left of=theta,blue!40!white] {$\theta$};

    \draw[->,gray!50!white] (-1,1.6) to[out=250,in=120] (theta);
    \draw[->,gray!50!white] (-1,3.6) to[out=250,in=120] (theta);
    \draw[->,gray!50!white] (-1,4.6) to[out=250,in=120] (theta);
    \draw[->,green!60!black,thick] (-1,2.6) to[out=250,in=120] (theta);

    \draw[dotted] (0,0) -- (theta);
    \fill[blue!20!white] (theta) circle (1pt);


    \draw[] (-1,2) arc (270:200:-3 and -0.7);
    \draw[] (-1,2) arc (270:90:1 and -0.2) arc (270:90:-3 and 0.7);
    \fill[blue!20!white] (-1,1.6) circle (1pt);
    \draw[] (-1,3) arc (270:90:1 and -0.2) arc (270:90:-3 and 0.7);
    \fill[green!60!black] (-1,2.6) circle (1pt);
    \draw[] (-1,4) arc (270:90:1 and -0.2) arc (270:90:-3 and 0.7);
    \fill[blue!20!white] (-1,3.6) circle (1pt);
    \draw[] (-1,5) arc (270:90:1 and -0.2) arc (270:200:-3 and 0.7);
    \fill[blue!20!white] (-1,4.6) circle (1pt);
  \end{tikzpicture}
\end{center}

\subsection{Definition of $\Lambda_\theta(A_0)$}
$\Lambda_\theta(A_0)$ $:\Leftrightarrow{}$ sheaf of flat isotropies over $S^1$
\marginnote{p. 854f \cite{Loday1994}}
\begin{defn}
  A germ of $\Lambda_\theta(A_0)$ at $\theta_0\in S^1$ is
  \begin{itemize}
    \item an invertible Matrix $f\in\Gl_n(\cO(U))$
      \begin{itemize}
        \item for suitable arc $U=U(\theta_0,\epsilon,\epsilon')$
      \end{itemize}
      satisfying:
      \begin{enumerate}
        \item \emph{Flatness:}
          \[
            \underset{x\in U}{\underset{x\to0}{\lim}}f(x)=\Id
            \text{ and }
            f\underset{U}{\sim}I
          \]
        \item \emph{Isotropy of $[A_0]$: ${}^f\!A_0=A_0$.}
      \end{enumerate}
  \end{itemize}
\end{defn}

\begin{center}
  \begin{tikzpicture}[scale=1]
    \node[] (zero) at (0,0) {};
    \fill (zero) circle (1pt);

    \fill[blue!20!white] (theta) circle (1pt);

    \filldraw[fill=green!20!white
      ,draw=green!60!black
      ,thick
    ,path fading=west] (0,0)
      -- ({cos( -50 )*4},{sin( -50 )*1}) arc (-50:20:4 and 1)
      -- cycle;

    \filldraw[fill=green!20!white
      ,draw=green!60!black
      ,thick
    ,path fading=west] (0,4)
      -- ({cos( -50 )*0.5},{4 + sin( -50 )*2})
      -- ({cos( -50 )*0.5 + cos( -10 )*3}
        ,{4 + sin( -50 )*2 + sin( -10 )*12})
      arc (-80:40:1 and 3)
      -- ({cos( 20 )*0.5},{4 + sin( 20 )*2})
      -- cycle;

    \fill[fill=white,draw=black] (0,2) arc (270:-90:0.5 and 2);

    \draw[blue!60!black,path fading=east]
      (zero) -- ({cos( -15 )*4},{sin( -15 )*1});

    \draw[dashed] (zero) -- (0,2);
  \end{tikzpicture}
\end{center}

\subsection{Definition of $\Sto_\alpha(A_0)$}
\begin{itemize}
  \item p. 861 in the paper
  \item p. 78 in the book
\end{itemize}

\subsection{Definition of $\cF$}

\subsection{Adequate coverings}
\begin{center}
  \begin{tikzpicture}[scale=2]
    \node (zero) at (0,0) {};
    \draw (zero) circle (1cm);

    \node at (-1,1) {$\cK=\{5\}:$};

    \foreach \w in {55}
    {\foreach \sep in {0,45,90,135,180,225,270,315}
     {\draw[thick,purple] (0,0) -- +({cos( \w + \sep )},{sin( \w + \sep )});}};

    \fill[white] (zero) circle (1.5pt);
    \fill (zero) circle (.7pt);
  \end{tikzpicture}
  \begin{tikzpicture}[scale=2]
    \node (zero) at (0,0) {};
    \draw (zero) circle (1cm);

    \node at (-1,1) {$\cK=\{5,4\}:$};

    \foreach \w in {55}
    {\foreach \sep in {0,45,90,135,180,225,270,315}
     {\draw[thick,purple] (0,0) -- +({cos( \w + \sep )},{sin( \w + \sep )});}};

    \foreach \w in {10,20,55}
    {\foreach \sep in {0,60,120,180,240,300}
     {\draw[blue] (0,0) -- +({cos( \w + \sep )},{sin( \w + \sep )});}};

    \fill[white] (zero) circle (1.5pt);
    \fill (zero) circle (.7pt);
  \end{tikzpicture}
  \begin{tikzpicture}[scale=2]
    \node (zero) at (0,0) {};
    \draw (zero) circle (1cm);

    \node at (-1,1) {$\cK=\{5,4,3\}:$};

    \foreach \w in {55}
    {\foreach \sep in {0,45,90,135,180,225,270,315}
     {\draw[thick,purple] (0,0) -- +({cos( \w + \sep )},{sin( \w + \sep )});}};

    \foreach \w in {10,20,55}
    {\foreach \sep in {0,60,120,180,240,300}
     {\draw[blue] (0,0) -- +({cos( \w + \sep )},{sin( \w + \sep )});}};

    \foreach \w in {30,35}
    {\foreach \sep in {0,90,180,270}
     {\draw[red] (0,0) -- +({cos( \w + \sep )},{sin( \w + \sep )});}};

    \fill[white] (zero) circle (1.5pt);
    \fill (zero) circle (.7pt);
  \end{tikzpicture}
\end{center}

Let $k\in\cK$.
\subsubsection{The \emph{cyclic covering $\cU^k=\{U_\alpha^k;\alpha\in\A^k\}$}}
\begin{center}
  \begin{tikzpicture}[scale=2.5]
    \node (zero) at (0,0) {};
    \draw (zero) circle (1cm);

    \node at (-1,1) {$\textcolor{purple}{\dot\cU^{5}}:$};
    \node at (-1.6,1.6) {$\textcolor{blue}{\dot\cU^{4}}:$};
    \node at (-2,2) {$\textcolor{green!60!black}{\dot\cU^{3}}:$};

    %%%%%%%%%%%%%%%%%%%%%%%%%%%%%%%%%%%%%%%%%%%%%%%%%%%%%%%%%%%%%%%%%%%%%%%%%%%
    %%%%  Purple
    \foreach \w in {10,25}
    {\foreach \sep in {0,45,90,135,180,225,270,315}
     {\draw[thick,purple] (0,0) -- +({cos( \w + \sep )},{sin( \w + \sep )});

      \pgfmathsetmacro\r{{1.1 + mod(\w,10)/100 + mod(\sep,2)/10}}
      \draw[dotted,purple] ({cos( \w + \sep )},{sin( \w + \sep )})
        -- ({cos( \w + \sep) * \r},{sin( \w + \sep) * \r});

      \draw[thick,purple] ({cos( \w + \sep -22) * \r},{sin( \w + \sep -22) * \r})
        arc ({\w + \sep - 22}:{\w + \sep +22}:\r);
      }};

    %%%%%%%%%%%%%%%%%%%%%%%%%%%%%%%%%%%%%%%%%%%%%%%%%%%%%%%%%%%%%%%%%%%%%%%%%%%
    %%%%  Blue
    \foreach \w in {10,20,55}
    {\foreach \sep in {0,60,120,180,240,300}
     {\draw[blue] (0,0) -- +({cos( \w + \sep ) * 0.9},{sin( \w + \sep ) * 0.9});}};

    \foreach \sep in {0,180}
    {\pgfmathsetmacro\rr{1.6}
     \foreach \w in {10,20,55
                    ,70,80,115
                    ,130,140,175}
     {\pgfmathsetmacro\rr{{\rr + (mod(\w,60) - 5*mod(\w,10))/150 
                               + mod(\w + \sep - mod(\w,60),120)/300}}
      \draw[dotted,blue] ({cos( \w + \sep )},{sin( \w + \sep )})
        -- ({cos( \w + \sep) * \rr},{sin( \w + \sep) * \rr});
      \draw[blue,thick] ({cos(\sep+\w -29) * \rr},{sin(\sep+\w -29) * \rr})
        arc ({\sep+\w -29}:{\sep+\w +29}:\rr);
     }
    }

    %%%%%%%%%%%%%%%%%%%%%%%%%%%%%%%%%%%%%%%%%%%%%%%%%%%%%%%%%%%%%%%%%%%%%%%%%%%
    %%%%  Green
    \foreach \w in {30,35}
    {\foreach \sep in {0,90,180,270}
     {\draw[green!60!black] (0,0) -- +({cos( \w + \sep ) * 0.8},{sin( \w + \sep) * 0.8});}};

    \foreach \sep in {0,180}
    {\pgfmathsetmacro\rr{2.4}
     \foreach \w in {30,35
                    ,120,125}
     {\pgfmathsetmacro\rr{{\rr + mod(\w,10)/100 + mod(\w-mod(\w,90),180)/900}}
      \draw[dotted,green!60!black] ({cos( \w + \sep )},{sin( \w + \sep )})
        -- ({cos( \w + \sep) * \rr},{sin( \w + \sep) * \rr});
      \draw[green!60!black,thick] ({cos(\sep+\w -44) * \rr},{sin(\sep+\w -44) * \rr})
        arc ({\sep+\w -44}:{\sep+\w +44}:\rr);
     }
    }

    \fill[white] (zero) circle (1.5pt);
    \fill (zero) circle (.7pt);
  \end{tikzpicture}
\end{center}

\subsubsection{The \emph{cyclic covering
  $\cU^{\leq k}=\{U_\alpha^{\leq k};\alpha\in\A^{\leq k}\}$}}
\begin{center}
  \begin{tikzpicture}[scale=2.5]
    \node (zero) at (0,0) {};
    \draw (zero) circle (1cm);

    \node at (-1.05,1.05) {$\textcolor{purple}{\dot\cU^{\leq 5}}:$};
    \node at (-1.65,1.65) {$\textcolor{blue}{\dot\cU^{\leq 4}}:$};
    \node at (-2.05,2.05) {$\textcolor{green!60!black}{\dot\cU^{\leq 3}}:$};

    %%%%%%%%%%%%%%%%%%%%%%%%%%%%%%%%%%%%%%%%%%%%%%%%%%%%%%%%%%%%%%%%%%%%%%%%%%%
    %%%%  Purple
    \foreach \w in {10,25}
    % {\foreach \sep in {0,45,90,135,180,225,270,315}
    {\foreach \sep in {0,36,72,108,144,180,216,252,288,324}
     {\draw[thick,purple] (0,0) -- +({cos( \w + \sep )},{sin( \w + \sep )});

      \pgfmathsetmacro\r{{1.1 + mod(\w,10)/100 + mod(\sep,72)/360}}
      \draw[dotted,purple] ({cos( \w + \sep )},{sin( \w + \sep )})
        -- ({cos( \w + \sep) * \r},{sin( \w + \sep) * \r});

      \draw[thick,purple] ({cos( \w + \sep -17.5) * \r},{sin( \w + \sep -17.5) * \r})
        arc ({\w + \sep - 17.5}:{\w + \sep +17.5}:\r);

      % \draw[purple] ({cos( \w + \sep -17.5) * \r},{sin( \w + \sep -17.5) * \r})
      %   -- ({cos( \w + \sep -17.5) * 1.5},{sin( \w + \sep -17.5) * 1.5});
      % \draw[purple] ({cos( \w + \sep +17.5) * \r},{sin( \w + \sep +17.5) * \r})
      %   -- ({cos( \w + \sep +17.5) * 1.5},{sin( \w + \sep +17.5) * 1.5});
      }};

    %%%%%%%%%%%%%%%%%%%%%%%%%%%%%%%%%%%%%%%%%%%%%%%%%%%%%%%%%%%%%%%%%%%%%%%%%%%
    %%%%  Blue
    \foreach \w in {10,20,55}
    {\foreach \sep in {0,60,120,180,240,300}
     {\draw[blue] (0,0) -- +({cos( \w + \sep ) * 0.9},{sin( \w + \sep ) * 0.9});}};

    \foreach \sep in {0,180}
    {\foreach \w in {10,20,55
                    ,70,80,115
                    ,130,140,175}
     {\pgfmathsetmacro\rr{{1.7 + (mod(\w,60) - 5*mod(\w,10))/150 
                               + mod(\w + \sep - mod(\w,60),120)/300}}
      \draw[dotted,blue] ({cos( \w + \sep )},{sin( \w + \sep )})
        -- ({cos( \w + \sep) * \rr},{sin( \w + \sep) * \rr});
      \draw[gray!42!white] ({cos(\sep+\w -29) * \rr},{sin(\sep+\w -29) * \rr})
        arc ({\sep+\w -29}:{\sep+\w +29}:\rr);

      % \draw[gray!42!white] ({cos( \w + \sep -29) * \rr},{sin( \w + \sep -29) * \rr})
      %   -- ({cos( \w + \sep -29) * 1.5},{sin( \w + \sep -29) * 1.5});
      % \draw[gray!42!white] ({cos( \w + \sep +29) * \rr},{sin( \w + \sep +29) * \rr})
      %   -- ({cos( \w + \sep +29) * 1.5},{sin( \w + \sep +29) * 1.5});
     }

     \pgfmathsetmacro\w{10}
     \pgfmathsetmacro\rr{{1.7 + (mod(\w,60) - 5*mod(\w,10))/150 
                              + mod(\w + \sep - mod(\w,60),120)/300}}
     \draw[thick,blue] ({cos(\sep+\w -17.5) * \rr},{sin(\sep+\w -17.5) * \rr})
       arc ({\sep+\w -17.5}:{\sep+\w +17.5}:\rr);

     \pgfmathsetmacro\w{20}
     \pgfmathsetmacro\rr{{1.7 + (mod(\w,60) - 5*mod(\w,10))/150 
                              + mod(\w + \sep - mod(\w,60),120)/300}}
     \draw[thick,blue] ({cos(\sep+ -7.5) * \rr},{sin(\sep+ -7.5) * \rr})
       arc ({\sep+ -7.5}:{\sep+ 42.5}:\rr);

     \pgfmathsetmacro\w{55}
     \pgfmathsetmacro\rr{{1.7 + (mod(\w,60) - 5*mod(\w,10))/150 
                              + mod(\w + \sep - mod(\w,60),120)/300}}
     \draw[thick,blue] ({cos(\sep+ 28.5) * \rr},{sin(\sep+ 28.5) * \rr})
       arc ({\sep+ 28.5}:{\sep+ 78.5}:\rr);

     \pgfmathsetmacro\w{70}
     \pgfmathsetmacro\rr{{1.7 + (mod(\w,60) - 5*mod(\w,10))/150 
                              + mod(\w + \sep - mod(\w,60),120)/300}}
     \draw[thick,blue] ({cos(\sep+ 43.5) * \rr},{sin(\sep+ 43.5) * \rr})
       arc ({\sep+ 43.5}:{\sep + \w + 29}:\rr);

     \pgfmathsetmacro\w{80}
     \pgfmathsetmacro\rr{{1.7 + (mod(\w,60) - 5*mod(\w,10))/150 
                              + mod(\w + \sep - mod(\w,60),120)/300}}
     \draw[thick,blue] ({cos(\sep+ \w - 29) * \rr},{sin(\sep+ \w - 29) * \rr})
       arc ({\sep+ \w - 29}:{\sep + 99.5}:\rr);

     \pgfmathsetmacro\w{115}
     \pgfmathsetmacro\rr{{1.7 + (mod(\w,60) - 5*mod(\w,10))/150 
                              + mod(\w + \sep - mod(\w,60),120)/300}}
     \draw[thick,blue] ({cos(\sep+ \w - 29) * \rr},{sin(\sep+ \w - 29) * \rr})
       arc ({\sep+ \w - 29}:{\sep + 135.5}:\rr);

     \pgfmathsetmacro\w{130}
     \pgfmathsetmacro\rr{{1.7 + (mod(\w,60) - 5*mod(\w,10))/150 
                              + mod(\w + \sep - mod(\w,60),120)/300}}
     \draw[thick,blue] ({cos(\sep+ \w - 29) * \rr},{sin(\sep+ \w - 29) * \rr})
       arc ({\sep+ \w - 29}:{\sep + 150.5}:\rr);

     \pgfmathsetmacro\w{140}
     \pgfmathsetmacro\rr{{1.7 + (mod(\w,60) - 5*mod(\w,10))/150 
                              + mod(\w + \sep - mod(\w,60),120)/300}}
     \draw[thick,blue] ({cos(\sep+ 115.5) * \rr},{sin(\sep+ 115.5) * \rr})
       arc ({\sep+ 115.5}:{\sep + \w + 29}:\rr);

     \pgfmathsetmacro\w{175}
     \pgfmathsetmacro\rr{{1.7 + (mod(\w,60) - 5*mod(\w,10))/150 
                              + mod(\w + \sep - mod(\w,60),120)/300}}
     \draw[thick,blue] ({cos(\sep+ 151.5) * \rr},{sin(\sep+ 151.5) * \rr})
       arc ({\sep+ 151.5}:{\sep + \w + 29}:\rr);
    }

    %%%%%%%%%%%%%%%%%%%%%%%%%%%%%%%%%%%%%%%%%%%%%%%%%%%%%%%%%%%%%%%%%%%%%%%%%%%
    %%%%  Green
    \foreach \w in {30,35}
    {\foreach \sep in {0,90,180,270}
     {\draw[green!60!black] (0,0) -- +({cos( \w + \sep ) * 0.8},{sin( \w + \sep) * 0.8});}};

    \foreach \sep in {0,180}
    {\foreach \w in {30,35
                    ,120,125}
     {\pgfmathsetmacro\rr{{2.5 + mod(\w,10)/100 + mod(\w-mod(\w,90),180)/900}}
      \draw[dotted,green!60!black] ({cos( \w + \sep )},{sin( \w + \sep )})
        -- ({cos( \w + \sep) * \rr},{sin( \w + \sep) * \rr});
      \draw[gray!42!white] ({cos(\sep+\w -44) * \rr},{sin(\sep+\w -44) * \rr})
        arc ({\sep+\w -44}:{\sep+\w +44}:\rr);
     }


     \pgfmathsetmacro\w{30}
     \pgfmathsetmacro\rr{{2.5 + mod(\w,10)/100 + mod(\w-mod(\w,90),180)/900}}
     \draw[thick,green!60!black] ({cos(\sep+ 7.5) * \rr},{sin(\sep+ 7.5) * \rr})
       arc ({\sep+ 7.5}:{\sep +63.5}:\rr);
     \pgfmathsetmacro\w{35}
     \pgfmathsetmacro\rr{{2.5 + mod(\w,10)/100 + mod(\w-mod(\w,90),180)/900}}
     \draw[thick,green!60!black] ({cos(\sep+ 7.5) * \rr},{sin(\sep+ 7.5) * \rr})
       arc ({\sep+ 7.5}:{\sep +63.5}:\rr);

     \pgfmathsetmacro\w{120}
     \pgfmathsetmacro\rr{{2.5 + mod(\w,10)/100 + mod(\w-mod(\w,90),180)/900}}
     \draw[thick,green!60!black] ({cos(\sep+ 100.5) * \rr},{sin(\sep+ 100.5) * \rr})
       arc ({\sep+ 100.5}:{\sep +150.5}:\rr);
     \pgfmathsetmacro\w{125}
     \pgfmathsetmacro\rr{{2.5 + mod(\w,10)/100 + mod(\w-mod(\w,90),180)/900}}
     \draw[thick,green!60!black] ({cos(\sep+ 100.5) * \rr},{sin(\sep+ 100.5) * \rr})
       arc ({\sep+ 100.5}:{\sep +150.5}:\rr);
    }

    \fill[white] (zero) circle (1.5pt);
    \fill (zero) circle (.7pt);
  \end{tikzpicture}
\end{center}

\subsection{Proof}
\subsubsection{Unique level} \marginnote{p. 872}
\subsubsection{Multi level} \marginnote{p. 872}

\input{notes/2014-12-19_MeromZsh.tex}
\input{notes/2014-12-20_Asymptotic.tex}
\input{notes/2014-12-21_Stokes.tex}

    \input{app-aufgaben.tex}
  \fi

  % \nocite{*}
  \makeatletter
  \@mainmatterfalse
  \makeatother
  \addcontentsline{toc}{chapter}{Bibliography}
  \printbibliography[keyword={prim},title={Bibliography}]
  \KOMAoptions{open=any}
  \printbibliography[notkeyword={prim},title={Additional references}]
\fi
\newpage\thispagestyle{empty}
\mbox{}\vfill
\begin{center}
  \includegraphics[width=0.2\textwidth]{qrCode.eps}
\end{center}
\begin{center}
  \url{https://github.com/maximilianhuber/MastersThesis}
\end{center}

%%% Local Variables:
%%% TeX-master: "Maximilian_Huber-Masters_Thesis-with_notes.tex"
%%% End:


% \ifnum\myDevelopVariable=1
%   \pagebreak
%   \layout
% \fi
\end{document}
% vim:set ft=tex foldmethod=marker foldmarker={{{,}}}:

%%% Local Variables:
%%% mode: latex
%%% TeX-master: t
%%% End:
